%% capacity_appendix.tex
%% Self-contained appendix: complete proof of the capacity theorem.
%% This is a FRAGMENT: no preamble, no document environment.
%% Requires: amsmath, amssymb, amsthm; theorem environments
%%   theorem, lemma, proposition, remark, definition;
%%   the paper macros T, R, Z, PP.  (NN and Cov are not used here.)
%% Intended use:  %% capacity_appendix.tex
%% Self-contained appendix: complete proof of the capacity theorem.
%% This is a FRAGMENT: no preamble, no document environment.
%% Requires: amsmath, amssymb, amsthm; theorem environments
%%   theorem, lemma, proposition, remark, definition;
%%   the paper macros T, R, Z, PP.  (NN and Cov are not used here.)
%% Intended use:  %% capacity_appendix.tex
%% Self-contained appendix: complete proof of the capacity theorem.
%% This is a FRAGMENT: no preamble, no document environment.
%% Requires: amsmath, amssymb, amsthm; theorem environments
%%   theorem, lemma, proposition, remark, definition;
%%   the paper macros T, R, Z, PP.  (NN and Cov are not used here.)
%% Intended use:  %% capacity_appendix.tex
%% Self-contained appendix: complete proof of the capacity theorem.
%% This is a FRAGMENT: no preamble, no document environment.
%% Requires: amsmath, amssymb, amsthm; theorem environments
%%   theorem, lemma, proposition, remark, definition;
%%   the paper macros T, R, Z, PP.  (NN and Cov are not used here.)
%% Intended use:  \input{capacity_appendix}  inside the paper's \appendix.

\section{Complete proof of the capacity bound (Theorem~\ref{capacity})}
\label{app:capacity}

This appendix proves Theorem~\ref{capacity} in full. Nothing below is
quoted from elsewhere except four classical facts, each isolated and
labelled as such (Remarks~\ref{app:classical-riesz},
\ref{app:classical-homog}, \ref{app:classical-density} and
\ref{app:classical-leray}); no numerical statement is used anywhere.

The chain is: the family $u_N$ has an exact scaling
$\hat u_N(k)=N^{-2}F(k/N)$ (\S\ref{app:profile}), so that the lattice
sums reconstructing $u_N$ near the origin are \emph{literally} Riemann
sums, of spacing $1/N$, for the inverse Fourier transform $V$ of $F$
(\S\ref{app:riemann}); pairing $\PP(u_N\cdot\nabla u_N)$ against a test
field concentrated at scale $1/N$ therefore converts, with an explicit
error, into the fixed continuum pairing
$\int_{\R^3}(V\cdot\nabla V)\cdot\Psi$ times $N^{3/2}$
(\S\ref{app:duality}); and that continuum pairing is nonzero for some
$\Psi$ because $\PP_{\R^3}(V\cdot\nabla V)\not\equiv0$, which is the
Fourier form of the elementary fact that the Oseen field is not a
stationary Euler flow (\S\ref{app:nondeg}). The bookkeeping is assembled
in \S\ref{app:final}.

\subsection{Conventions}
\label{app:conventions}

On $\T^3=(\R/2\pi\Z)^3$ we use the normalised measure
$d\mu=(2\pi)^{-3}dx$, so that for $u:\T^3\to\R^3$
\[
u(x)=\sum_{k\in\Z^3}\hat u(k)e^{ik\cdot x},\qquad
\hat u(k)=\int_{\T^3}u(x)e^{-ik\cdot x}\,d\mu(x),
\]
\[
\langle f,g\rangle=\int_{\T^3}f\cdot\bar g\,d\mu
=\sum_{k\in\Z^3}\hat f(k)\cdot\overline{\hat g(k)},
\qquad
H_r(u)=\sum_{k\in\Z^3}|k|^{2r}|\hat u(k)|^2 .
\]
On $\R^3$ we use the forward transform and the inverse-transform
bracket
\[
\mathcal F f(\zeta)=\int_{\R^3}f(y)e^{-i\zeta\cdot y}\,dy,
\qquad
\widetilde G(y)=\int_{\R^3}G(\xi)e^{i\xi\cdot y}\,d\xi ,
\]
so that $\mathcal F\widetilde G=(2\pi)^3G$ and
$\|\widetilde G\|_{L^2(\R^3)}^2=(2\pi)^3\|G\|_{L^2(\R^3)}^2$.

For $\xi\in\R^3\setminus\{0\}$ put
$P_\xi=I-\dfrac{\xi\otimes\xi}{|\xi|^2}$. Thus $P_\xi$ is the orthogonal
projection onto $\xi^\perp$: it is symmetric, idempotent, $P_\xi\xi=0$,
$\|P_\xi\|\le1$, $P_{-\xi}=P_\xi$, it is homogeneous of degree $0$ and
smooth on $\R^3\setminus\{0\}$, and consequently
$P_{k}=P_{k/N}$ for every $k\neq0$ and every $N>0$. The Leray
projection $\PP$ acts on $\T^3$ by
$\widehat{\PP f}(k)=P_k\hat f(k)$ for $k\neq0$, leaving the mean
untouched, and on $\R^3$ by the same symbol $P_\zeta$. On both
$L^2(\T^3)$ and $L^2(\R^3)$, $\PP$ is an orthogonal projection, hence
self-adjoint.

Throughout, $\|v_0\|$ is the Euclidean norm of $v_0\in\R^3$, and
$e_1,e_2,e_3$ is the standard basis. Repeated Latin indices are summed.
One notational warning: in this appendix $V$ always denotes the
continuum profile of Definition~\ref{app:Vdef}; the unrelated use of
$V$ for the modal variance $\mathrm{Var}_p(|k|^2)$ in
Theorem~\ref{main}(a) does not occur here.

\subsection{The cutoff class, the family, and its exact laws}
\label{app:laws}

We restate the definitions of the main text for completeness.

\begin{definition}[admissible cutoff; restatement of
Definition~\ref{admissible}]\label{app:admissible}
$\chi\in C^4([0,\infty);[0,1])$ with $\chi\equiv1$ on
$[0,\tfrac12]$ and $\operatorname{supp}\chi\subset[0,1]$. No
monotonicity is assumed --- it is used nowhere. Write
$c_\chi:=1+\|\chi'\|_{L^\infty}\ (\ge1)$.
\end{definition}

Such $\chi$ exist in abundance. The standard mollifier construction
gives a $C^\infty$ example, $\chi(t)=\Theta\big(2t-1\big)$ with
$\Theta(s)=1$ for $s\le0$, $\Theta(s)=0$ for $s\ge1$ and
$\Theta(s)=\frac{e^{-1/(1-s)}}{e^{-1/s}+e^{-1/(1-s)}}$ on $(0,1)$; the
exactly rational cutoff of Appendix~\ref{app:certificates}, built from
the degree-$9$ smoothstep, lies in $C^4\setminus C^5$ and is likewise
admissible.

\begin{remark}[what smoothness is actually used]\label{app:c1}
Definition~\ref{app:admissible} is still stronger than the proof
requires. The only place any derivative of $\chi$ is used
load-bearingly is Theorem~\ref{app:rate}, and it enters solely through
$c_\chi=1+\|\chi'\|_\infty$; Sections~\ref{app:duality} and
\ref{app:nondeg} use only $0\le\chi\le1$, $\chi\equiv1$ on
$[0,\tfrac12]$ and $\operatorname{supp}\chi\subset[0,1]$, and
Sections~\ref{app:laws}, \ref{app:lattice}, \ref{app:profile} use no
derivative of $\chi$ at all. Every statement of this appendix on which
Theorem~\ref{app:main} depends therefore holds verbatim for
$\chi\in C^1([0,\infty);[0,1])$ --- indeed for Lipschitz $\chi$ ---
with those support properties. The four derivatives are demanded only
so that Theorem~\ref{app:farfield} holds as stated with $M=4$; that
theorem is consumed only by the inert Remark~\ref{app:noperiod}, and
nothing in the proof of Theorem~\ref{app:main} uses it. In particular
$V\in C^\infty$ (Theorem~\ref{app:Vprops}(3)) and the Paley--Wiener
statement there are bought by the \emph{compact support} of $F$, not by
smoothness of $\chi$, and are unaffected.
\end{remark}

\begin{definition}\label{app:family}
Fix $v_0\in\R^3\setminus\{0\}$, an admissible $\chi$, and an integer
$N\ge2$. Define $u_N=u_N^{\chi,v_0}$ on $\T^3$ by
\begin{equation}\label{app:famdef}
\hat u_N(k)=\chi\Big(\frac{|k|}{N}\Big)\frac{P_kv_0}{|k|^{2}}
\quad(k\neq0),\qquad \hat u_N(0)=0 .
\end{equation}
\end{definition}

Since $\hat u_N(k)=0$ for $|k|>N$, $u_N$ is a finite trigonometric
polynomial banded at $|k|\le N$; its coefficients are real vectors and
no integrality of $v_0$ is required.

\begin{lemma}[basic structure]\label{app:structure}
$u_N$ is real, zero-mean, exactly divergence-free, $C^\infty$, and
$u_N\not\equiv0$.
\end{lemma}

\begin{proof}
Zero mean is $\hat u_N(0)=0$. Reality: the coefficients are real
vectors and $P_{-k}=P_k$, so
$\hat u_N(-k)=\hat u_N(k)=\overline{\hat u_N(k)}$, which is exactly the
reality condition $\hat u(-k)=\overline{\hat u(k)}$. Divergence-free:
$\widehat{\nabla\cdot u_N}(k)=ik\cdot\hat u_N(k)$ and $k\cdot P_kv_0=0$
for every $k\neq0$. Smoothness: a finite sum of exponentials. Finally,
among $e_1,e_2,e_3$ at least two are not parallel to $v_0$; pick such a
$k$ with $|k|=1$. Then $\chi(1/N)=1$ because $1/N\le\tfrac12$, and
$P_kv_0\neq0$, so $\hat u_N(k)\neq0$.
\end{proof}

The exact laws rest on a band-symmetry identity which, because $\chi$ is
not assumed monotone and the truncation is not a sharp ball, must be
available for a \emph{general} radial weight, with no positivity
assumption.

\begin{lemma}[band symmetry, lattice, general radial weight]
\label{app:bandsym}
Let $f:\Z^3\setminus\{0\}\to\R$ be radial, i.e.\ $f(k)=\varphi(|k|^2)$
for some $\varphi$, and suppose $\sum_{k\neq0}|k|^2|f(k)|<\infty$. Then
for all $i,j$
\[
\sum_{k\neq0}k_ik_j\,f(k)
=\delta_{ij}\,\frac13\sum_{k\neq0}|k|^2f(k).
\]
\end{lemma}

\begin{proof}
Absolute summability of $|k|^2f$ dominates every summand
$|k_ik_jf(k)|\le|k|^2|f(k)|$, so all rearrangements below are
legitimate. The index set $\Z^3\setminus\{0\}$ is invariant under each
sign flip $\sigma_i:k_i\mapsto-k_i$ and under all coordinate
permutations, and $f$, depending on $|k|^2$ only, is invariant under
both. For $i\neq j$ the bijection $\sigma_i$ negates the summand
$k_ik_jf(k)$, so that sum equals its own negative and vanishes. For
$i=j$, coordinate permutations carry the three sums $\sum_kk_i^2f(k)$
into one another, so they are equal; their sum is $\sum_k|k|^2f(k)$.
\end{proof}

\begin{lemma}[band symmetry, continuum]\label{app:bandsymcont}
Let $g:\R^3\to\R$ be radial with $\int_{\R^3}|\xi|^2|g(\xi)|\,d\xi<\infty$.
Then $\int\xi_i\xi_jg\,d\xi=\delta_{ij}\frac13\int|\xi|^2g\,d\xi$.
\end{lemma}

\begin{proof}
Identical: Lebesgue measure on $\R^3$ and $g$ are both invariant under
sign flips and permutations of the coordinates, and the hypothesis
dominates the integrand.
\end{proof}

\begin{theorem}[exact family laws]\label{app:exactlaws}
Set
\[
S_N^\chi:=\sum_{k\neq0}\frac{\chi(|k|/N)^2}{|k|^2},\qquad
T_N^\chi:=\sum_{k\neq0}\frac{\chi(|k|/N)^2}{|k|^4},\qquad
\Sigma_N^\chi:=\sum_{k\neq0}\frac{\chi(|k|/N)}{|k|^2},
\]
all three finite sums. Then
\[
H_0(u_N)=\tfrac23\|v_0\|^2\,T_N^\chi,\qquad
H_1(u_N)=\tfrac23\|v_0\|^2\,S_N^\chi,\qquad
N_0^2(u_N)=\frac{S_N^\chi}{T_N^\chi},\qquad
u_N(0)=\tfrac23\,\Sigma_N^\chi\,v_0 ,
\]
and consequently
$\|u_N\|_\infty^2/H_1(u_N)\ge 2(\Sigma_N^\chi)^2/(3S_N^\chi)$.
\end{theorem}

\begin{proof}
Since $P_k$ is an orthogonal projection,
$|P_kv_0|^2=\|v_0\|^2-(k\cdot v_0)^2/|k|^2$. Hence
\[
H_0=\sum_{k\neq0}\frac{\chi(|k|/N)^2}{|k|^4}\,|P_kv_0|^2
=\|v_0\|^2\sum_{k\neq0}\frac{\chi^2}{|k|^4}
-\sum_{k\neq0}\frac{\chi^2\,(k\cdot v_0)^2}{|k|^6}.
\]
Apply Lemma~\ref{app:bandsym} with the radial weight
$f(k)=\chi(|k|/N)^2|k|^{-6}$ (a finite sum, so the hypothesis is
trivially met):
\[
\sum_{k\neq0}\frac{\chi^2(k\cdot v_0)^2}{|k|^6}
=v_{0i}v_{0j}\sum_{k\neq0}k_ik_jf(k)
=\frac{\|v_0\|^2}{3}\sum_{k\neq0}\frac{\chi^2}{|k|^4}
=\frac{\|v_0\|^2}{3}T_N^\chi .
\]
Therefore $H_0=\|v_0\|^2T_N^\chi-\frac13\|v_0\|^2T_N^\chi
=\frac23\|v_0\|^2T_N^\chi$. The same computation with
$f=\chi^2|k|^{-4}$ applied to
$H_1=\sum_{k}|k|^2|\hat u_N(k)|^2$ gives
$H_1=\frac23\|v_0\|^2S_N^\chi$, and $N_0^2=H_1/H_0$ follows. For the
point value,
\[
u_N(0)=\sum_{k\neq0}\hat u_N(k)
=v_0\,\Sigma_N^\chi-\sum_{k\neq0}\frac{\chi\,k\,(k\cdot v_0)}{|k|^4},
\]
and Lemma~\ref{app:bandsym} with $f=\chi|k|^{-4}$ turns the last sum
into $\frac13v_0\Sigma_N^\chi$; hence
$u_N(0)=\frac23\Sigma_N^\chi v_0$. Finally
$\|u_N\|_\infty\ge|u_N(0)|=\frac23\Sigma_N^\chi\|v_0\|$, so
$\|u_N\|_\infty^2/H_1\ge\frac49(\Sigma_N^\chi)^2\|v_0\|^2
\big/\big(\frac23\|v_0\|^2S_N^\chi\big)
=\frac23(\Sigma_N^\chi)^2/S_N^\chi$.
\end{proof}

\subsection{Two-sided lattice bounds for the bandwidth}
\label{app:lattice}

The constants below are crude and explicit; no sharpness is attempted.

\begin{lemma}[dyadic shell count]\label{app:shell}
For every integer $j\ge0$,
$\#\{k\in\Z^3:2^j\le|k|<2^{j+1}\}\le64\cdot8^{j}$.
\end{lemma}

\begin{proof}
Such $k$ satisfy $|k|_\infty\le|k|<2^{j+1}$, hence
$|k_i|\le2^{j+1}-1$ for each $i$, so the count is at most
$(2^{j+2}-1)^3<2^{3j+6}=64\cdot8^{j}$.
\end{proof}

\begin{lemma}[upper bounds]\label{app:upper}
With $S_M:=\sum_{1\le|k|\le M}|k|^{-2}$ and
$T_\infty:=\sum_{k\neq0}|k|^{-4}$ one has
$S_M\le128M$ for every real $M\ge1$, and $T_\infty\le128$.
\end{lemma}

\begin{proof}
Let $J=\lfloor\log_2M\rfloor$. Every $k$ with $1\le|k|\le M$ lies in a
shell $2^j\le|k|<2^{j+1}$ with $0\le j\le J$, where $|k|^{-2}\le4^{-j}$.
By Lemma~\ref{app:shell},
$S_M\le\sum_{j=0}^{J}64\cdot8^{j}4^{-j}=64(2^{J+1}-1)<128\cdot2^{J}\le128M$.
Similarly $T_\infty\le\sum_{j\ge0}64\cdot8^{j}16^{-j}=64\cdot2=128$.
\end{proof}

\begin{lemma}[lower bound]\label{app:lower}
For every real $M\ge16$ one has $S_M\ge M/87$.
\end{lemma}

\begin{proof}
Consider
$\mathcal B:=\{k\in\Z^3:\ \tfrac{M}{2\sqrt3}<k_i\le\tfrac{M}{\sqrt3}\
\text{for }i=1,2,3\}$. For $k\in\mathcal B$ we have
$|k|^2\le3(M/\sqrt3)^2=M^2$ and $|k|^2>3M^2/12=M^2/4$, so every
$k\in\mathcal B$ is counted in $S_M$ (in particular $k\neq0$) and
contributes $|k|^{-2}\ge M^{-2}$. Each coordinate ranges over a
half-open interval of length $M/(2\sqrt3)$, which contains at least
$\frac{M}{2\sqrt3}-1$ integers; for $M\ge16$,
$\frac{M}{2\sqrt3}-1\ge0.28868M-0.0625M\ge0.226M$. Hence
$\#\mathcal B\ge(0.226M)^3>0.01154M^3$ and
$S_M\ge0.01154M^3\cdot M^{-2}=0.01154M\ge M/87$.
\end{proof}

\begin{proposition}[two-sided bandwidth bounds]\label{app:twosided}
For every admissible $\chi$, every $v_0\neq0$ and every integer
$N\ge32$,
\[
\frac{N}{348}\ \le\ S_N^\chi\ \le\ 128N,\qquad
6\ \le\ T_N^\chi\ \le\ 128,\qquad
\frac{N}{44544}\ \le\ N_0^2(u_N)\ \le\ \frac{64}{3}N .
\]
\end{proposition}

\begin{proof}
Upper bounds: $0\le\chi\le1$ and $\chi(|k|/N)=0$ for $|k|>N$, so
$S_N^\chi\le S_N\le128N$ and $T_N^\chi\le T_\infty\le128$ by
Lemma~\ref{app:upper}. Lower bound for $T_N^\chi$: for $N\ge2$ one has
$\chi(1/N)=1$, and the six lattice points with $|k|=1$ each contribute
$1$, so $T_N^\chi\ge6$. Lower bound for $S_N^\chi$: $\chi(|k|/N)=1$
whenever $|k|\le N/2$, hence $S_N^\chi\ge S_{N/2}$; for $N\ge32$ we have
$N/2\ge16$, so Lemma~\ref{app:lower} gives
$S_N^\chi\ge (N/2)/87\ge N/348$. Combining with
$N_0^2=S_N^\chi/T_N^\chi$ (Theorem~\ref{app:exactlaws}) yields
$N_0^2\ge\frac{N/348}{128}=\frac{N}{44544}$ and
$N_0^2\le\frac{128N}{6}=\frac{64}{3}N$.
\end{proof}

\subsection{The profile $V$}
\label{app:profile}

\begin{definition}\label{app:Vdef}
On $\R^3$ put
\[
h(\xi):=\frac{P_\xi v_0}{|\xi|^2}\quad(\xi\neq0),\qquad
F(\xi):=\chi(|\xi|)\,h(\xi)\quad(\xi\neq0),\quad F(0):=0,
\]
\[
V(y):=\widetilde F(y)=\int_{\R^3}F(\xi)e^{i\xi\cdot y}\,d\xi .
\]
\end{definition}

Because $P_\xi$ is homogeneous of degree $0$ we have
$P_{k/N}=P_k$ and therefore
$F(k/N)=\chi(|k|/N)\,N^2P_kv_0/|k|^2$, i.e.\ the exact scaling identity
\begin{equation}\label{app:scaling}
\hat u_N(k)=N^{-2}F(k/N)\qquad\text{for every }k\in\Z^3
\end{equation}
(the case $k=0$ holds by the conventions $\hat u_N(0)=0=F(0)$).
Identity~\eqref{app:scaling} is the engine of the whole argument: it
says that the family is a single fixed continuum profile sampled on the
lattice $N^{-1}\Z^3$.

\begin{theorem}[elementary properties of $V$]\label{app:Vprops}
\begin{enumerate}
\item[(1)] $|F(\xi)|\le\|v_0\|\,|\xi|^{-2}\mathbf 1_{\{|\xi|\le1\}}$;
hence $F\in L^1(\R^3)\cap L^p(\R^3)$ for every $p<3/2$, with
$\|F\|_{L^1}\le4\pi\|v_0\|$.
\item[(2)] $V$ is real, even, bounded and continuous, with
$\|V\|_\infty\le4\pi\|v_0\|$.
\item[(3)] $V\in C^\infty(\R^3)$, and for every multi-index $\beta$
\[
\partial^\beta V(y)=\int_{\R^3}(i\xi)^\beta F(\xi)e^{i\xi\cdot y}\,d\xi,
\qquad
\|\partial^\beta V\|_\infty\le\|v_0\|\int_{|\xi|\le1}|\xi|^{|\beta|-2}d\xi
=\frac{4\pi\|v_0\|}{|\beta|+1} .
\]
In particular each $\|\partial_iV_j\|_\infty\le2\pi\|v_0\|$, so
$\|\nabla V\|_\infty\le6\pi\|v_0\|$ for the Frobenius norm. Moreover $V$
extends to an entire function on $\mathbb C^3$ of exponential type
$\le1$.
\item[(4)] $\nabla\cdot V\equiv0$ on $\R^3$.
\item[(5)] $V(0)=\frac{8\pi}{3}\big(\int_0^\infty\chi(r)\,dr\big)v_0\neq0$.
\end{enumerate}
\end{theorem}

\begin{proof}
(1) $|P_\xi v_0|\le\|v_0\|$, $0\le\chi\le1$ and
$\operatorname{supp}\chi\subset[0,1]$ give the pointwise bound;
$\int_{|\xi|\le1}|\xi|^{-2}d\xi=4\pi\int_0^1dr=4\pi$, and
$\int_{|\xi|\le1}|\xi|^{-2p}d\xi<\infty$ exactly when $2p<3$.

(2) $F\in L^1$ gives, by dominated convergence, continuity of $V$ and
$\|V\|_\infty\le\|F\|_{L^1}\le4\pi\|v_0\|$. $F$ is real-valued and even
($P_{-\xi}=P_\xi$), so
$V(y)=\int F(\xi)\cos(\xi\cdot y)\,d\xi$ is real and even.

(3) By (1), $\xi^\beta F\in L^1$ for every $\beta$, since the exponent
$|\beta|-2>-3$ is locally integrable in $\R^3$ and the support is
bounded. Differentiation under the integral sign is then justified by
dominated convergence, giving the displayed formula and, by (1) again,
the bound
$\|v_0\|\int_{|\xi|\le1}|\xi|^{|\beta|-2}d\xi
=4\pi\|v_0\|\int_0^1r^{|\beta|}dr=4\pi\|v_0\|/(|\beta|+1)$. Since $F$
is an $L^1$ function supported in $\{|\xi|\le1\}$, the integral
$\int F(\xi)e^{i\xi\cdot y}d\xi$ converges for all complex $y$ and is
entire in $y$, with $|V(y)|\le\|F\|_{L^1}e^{|\operatorname{Im}y|}$;
this is exponential type $\le1$.

(4) $\nabla\cdot V(y)=\int i\,\xi\cdot F(\xi)e^{i\xi\cdot y}d\xi$ and
$\xi\cdot F(\xi)=\chi(|\xi|)\,\xi\cdot P_\xi v_0/|\xi|^2=0$ pointwise.

(5) $V(0)=\int F
=v_0\int\frac{\chi(|\xi|)}{|\xi|^2}d\xi
-\int\frac{\chi(|\xi|)\,\xi\,(\xi\cdot v_0)}{|\xi|^4}d\xi$.
Lemma~\ref{app:bandsymcont} with the radial
$g(\xi)=\chi(|\xi|)|\xi|^{-4}$ (integrable against $|\xi|^2$ by (1))
gives $\int\xi_i\xi_jg=\frac{\delta_{ij}}3\int|\xi|^2g
=\frac{\delta_{ij}}3\int\chi(|\xi|)|\xi|^{-2}d\xi$, so the second
integral equals $\frac13v_0\int\chi(|\xi|)|\xi|^{-2}d\xi$. Hence
\[
V(0)=\frac23v_0\int_{\R^3}\frac{\chi(|\xi|)}{|\xi|^2}\,d\xi
=\frac23v_0\cdot4\pi\int_0^\infty\chi(r)\,dr
=\frac{8\pi}{3}\Big(\int_0^\infty\chi\Big)v_0 ,
\]
and $\int_0^\infty\chi\ge\int_0^{1/2}1=\frac12>0$ because $\chi\ge0$ and
$\chi\equiv1$ on $[0,\tfrac12]$.
\end{proof}

\begin{remark}[structural form]\label{app:oseenform}
$F=P_\xi\big(v_0\chi(|\xi|)|\xi|^{-2}\big)$, so
$V=\PP\big(v_0\,g_\chi\big)$ where $g_\chi$ is the inverse transform of
the smoothed Newtonian symbol $\chi(|\cdot|)|\cdot|^{-2}$; solving
$\Delta\phi=g_\chi$ with $\phi$ radial exhibits $V$ in the Oseen normal
form $V(y)=a(r)v_0+b(r)(v_0\cdot\hat y)\hat y$ with
$a=\phi''+\phi'/r$, $b=\phi'/r-\phi''$. This form is not used below; it
explains why the far field of $V$ is exactly the Oseen profile
(Theorem~\ref{app:farfield}).
\end{remark}

We record here the two classical transform facts used later.

\begin{remark}[classical input: Riesz-potential transforms]
\label{app:classical-riesz}
For $0<\alpha<3$, and by analytic continuation for every $\alpha$ at
which the Gamma factors have no pole, the homogeneous function
$|y|^{-\alpha}$ on $\R^3$ is a tempered distribution whose Fourier
transform is the homogeneous distribution
\[
\mathcal F\big[|y|^{-\alpha}\big](\zeta)=c_{3,\alpha}|\zeta|^{\alpha-3},
\qquad
c_{3,\alpha}=\frac{2^{3-\alpha}\pi^{3/2}\,\Gamma\!\big(\frac{3-\alpha}{2}\big)}
{\Gamma\!\big(\frac{\alpha}{2}\big)} .
\]
See, e.g., \cite{SteinWeiss,GelfandShilov,Grafakos} (we give no section
locators, not having verified them). We use
$\Gamma(\tfrac12)=\sqrt\pi$, $\Gamma(-\tfrac12)=-2\sqrt\pi$,
$\Gamma(-\tfrac32)=\tfrac43\sqrt\pi$, which give
\[
c_{3,-1}=-8\pi,\quad c_{3,1}=4\pi,\quad c_{3,2}=2\pi^2,\quad
c_{3,4}=-\pi^2,\quad c_{3,6}=\tfrac{\pi^2}{12},
\]
together with the elementary rules
$\mathcal F[y_ay_bf]=-\partial_a\partial_b\mathcal F[f]$ and
$\mathcal F[y_ay_by_cy_df]=\partial_a\partial_b\partial_c\partial_d\mathcal F[f]$,
valid for tempered $f$.
\end{remark}

\begin{remark}[classical input: products and convolutions of homogeneous
distributions]\label{app:classical-homog}
Let $A,B\in L^1_{\rm loc}(\R^3)$ be homogeneous of degrees $-a,-b$ with
$0<a,b<3$ and $a+b<3$, and smooth away from the origin. Then $AB$ is
locally integrable and homogeneous of degree $-(a+b)$, hence tempered;
$\mathcal FA,\mathcal FB$ are locally integrable, homogeneous of degrees
$a-3,b-3$ and smooth away from the origin; the convolution
$\mathcal FA*\mathcal FB$ converges absolutely at every $\zeta\neq0$
(because $(3-a)+(3-b)>3$) and defines a locally integrable homogeneous
function of degree $3-a-b$; and the exchange formula
\[
\mathcal F[AB]=(2\pi)^{-3}\,\mathcal FA*\mathcal FB
\]
holds as tempered distributions. See, e.g.,
\cite{GelfandShilov,Grafakos}. We use this only in
Theorem~\ref{app:tauformula}, with $a=b=1$.
\end{remark}

\subsection{The far field, and why periodisation is unavailable}
\label{app:farfield-sec}

\begin{theorem}[far-field asymptotics]\label{app:farfield}
Let
\[
W(y):=\frac{v_0}{|y|}+\frac{(v_0\cdot y)\,y}{|y|^3}\qquad(y\neq0).
\]
Then $\widetilde h=\pi^2W$ as tempered distributions, and
\[
V(y)=\pi^2W(y)+\rho(y),\qquad
|\rho(y)|\le C_M(\chi,v_0)|y|^{-M}\quad\text{for }|y|\ge1
\text{ and every }2\le M\le4 .
\]
(Four derivatives are all that Definition~\ref{app:admissible}
supplies, and $M=4$ is all that Remark~\ref{app:noperiod} consumes; if
$\chi\in C^m$ with $m\ge2$ the same proof gives every
$2\le M\le m$.)
\end{theorem}

\begin{proof}
\emph{(a) The identity $\widetilde h=\pi^2W$.} Let
$\mathcal O_{ij}(y)=\frac1{8\pi}\big(\frac{\delta_{ij}}{|y|}
+\frac{y_iy_j}{|y|^3}\big)$ be the Oseen tensor, i.e.\ the fundamental
solution of the steady Stokes system
\cite{LadyzhenskayaBook,Galdi2011}, and
$h_{ij}(\xi)=\frac{\delta_{ij}}{|\xi|^2}-\frac{\xi_i\xi_j}{|\xi|^4}$, so
that $h_i=h_{ij}v_{0j}$. We verify $\mathcal F[\mathcal O_{ij}]=h_{ij}$
from Remark~\ref{app:classical-riesz}. First,
$\mathcal F[|y|^{-1}]=c_{3,1}|\zeta|^{-2}=4\pi|\zeta|^{-2}$, so
$\mathcal F\big[\tfrac{\delta_{ij}}{8\pi|y|}\big]
=\tfrac{\delta_{ij}}{2|\zeta|^{2}}$. Second, differentiating
$\partial_i|y|=y_i/|y|$ once more gives the pointwise identity
\[
\partial_i\partial_j|y|=\frac{\delta_{ij}}{|y|}-\frac{y_iy_j}{|y|^3},
\qquad\text{i.e.}\qquad
\frac{y_iy_j}{|y|^3}=\frac{\delta_{ij}}{|y|}-\partial_i\partial_j|y| .
\]
With $\mathcal F[|y|]=c_{3,-1}|\zeta|^{-4}=-8\pi|\zeta|^{-4}$ we get
$\mathcal F[\partial_i\partial_j|y|]
=(i\zeta_i)(i\zeta_j)\mathcal F[|y|]=8\pi\zeta_i\zeta_j|\zeta|^{-4}$, hence
\[
\mathcal F\Big[\frac{y_iy_j}{|y|^3}\Big]
=\frac{4\pi\delta_{ij}}{|\zeta|^2}-\frac{8\pi\zeta_i\zeta_j}{|\zeta|^4}.
\]
Therefore
$\mathcal F[\mathcal O_{ij}]
=\frac1{8\pi}\big(\frac{4\pi\delta_{ij}}{|\zeta|^2}
+\frac{4\pi\delta_{ij}}{|\zeta|^2}
-\frac{8\pi\zeta_i\zeta_j}{|\zeta|^4}\big)
=\frac{\delta_{ij}}{|\zeta|^2}-\frac{\zeta_i\zeta_j}{|\zeta|^4}=h_{ij}$.
All these are identities between homogeneous tempered distributions;
possible ambiguities are supported at the origin only and are excluded
because both sides are homogeneous of degree $-2>-3$, hence locally
integrable, with no room for a delta term. Fourier inversion
($\mathcal F\widetilde G=(2\pi)^3G$) now gives
$\widetilde h=(2\pi)^3\mathcal O\,v_0=\frac{(2\pi)^3}{8\pi}W=\pi^2W$,
where we used $8\pi\,\mathcal O_{ij}v_{0j}=W_i$.

\emph{(b) The remainder.} Write $G:=(1-\chi(|\cdot|))h$, so that
$F-h=-G$ and $\rho:=-\widetilde G$. Then $G$ vanishes on
$\{|\xi|<\tfrac12\}$ and is $C^4$ on all of $\R^3$; since $h$ is
homogeneous of degree $-2$ and smooth off the origin, and
$1-\chi(|\cdot|)$ is $C^4$ with its first four derivatives bounded, we
get
\[
|\partial^\beta G(\xi)|\le C_\beta(\chi)\|v_0\|\,\langle\xi\rangle^{-2-|\beta|}
\qquad(\xi\in\R^3,\ |\beta|\le4),
\]
which for $2\le|\beta|\le4$ places $\partial^\beta G$ in $L^1(\R^3)$
(the exponent $2+|\beta|>3$). From
the distributional identity
$y^\beta\widetilde G=i^{|\beta|}\,\widetilde{\partial^\beta G}$, the
right-hand side for $2\le|\beta|\le4$ is a bounded continuous function
with supremum at most $\|\partial^\beta G\|_{L^1}$. Choosing
$\beta=(M,0,0),(0,M,0),(0,0,M)$ for a given $2\le M\le4$ and using
$|y|^M\le3^{M/2}\max_i|y_i|^M$ shows that $\widetilde G$ agrees on
$\{|y|\ge1\}$ with a continuous function obeying
$|\widetilde G(y)|\le C_M(\chi,v_0)|y|^{-M}$. Since
$V=\widetilde F=\widetilde h-\widetilde G=\pi^2W+\rho$, the claim
follows.
\end{proof}

Thus the decay of $V$ is \emph{exactly} $|y|^{-1}$, with a
$\chi$-independent leading profile $\pi^2W$ satisfying
\begin{equation}\label{app:Wpos}
v_0\cdot W(y)=\frac{\|v_0\|^2}{|y|}+\frac{(v_0\cdot y)^2}{|y|^3}
\ \ge\ \frac{\|v_0\|^2}{|y|}\ >\ 0 .
\end{equation}

\begin{remark}[the naive periodisation argument is unavailable]
\label{app:noperiod}
It is tempting to reconstruct $u_N$ from $V$ by Poisson summation,
$u_N(x)\stackrel{?}{=}N\sum_{m\in\Z^3}V(N(x+2\pi m))$ with the $m\neq0$
terms uniformly $O(1/N)$. \textbf{The naive absolutely convergent
periodisation argument is unavailable}: the series of translates
diverges, because $V$ decays only like $|y|^{-1}$ and its leading
profile has a sign-definite spherical mean in the $v_0$ direction, so
there is no cancellation to exploit and no absolute convergence to
appeal to. Indeed, by Theorem~\ref{app:farfield} with $M=4$, for
$m\neq0$ and $x\in[-\pi,\pi)^3$ one has $|N(x+2\pi m)|\ge N\pi\ge1$ and
therefore, using \eqref{app:Wpos},
\[
v_0\cdot V\big(N(x+2\pi m)\big)
\ \ge\ \frac{\pi^2\|v_0\|^2}{N|x+2\pi m|}
-\frac{C_4\|v_0\|}{N^4|x+2\pi m|^4}.
\]
The shell $|m|_\infty=n$ contains $O(n^2)$ points with
$|x+2\pi m|\asymp n$, so $\sum_{m\neq0}|x+2\pi m|^{-4}<\infty$ while
$\sum_{m\neq0}|x+2\pi m|^{-1}=\infty$; hence
$\sum_{m\neq0}v_0\cdot V(N(x+2\pi m))=+\infty$ for every $x$ and every
$N$. The terms are moreover eventually positive, so the divergence is
not of a kind that a regular linear summation method would remove: for
$a_m\ge0$ with $\sum a_m=\infty$, monotone convergence gives
$\lim_{\varepsilon\downarrow0}\sum_me^{-\varepsilon|m|}a_m=+\infty$, so
Abel summation in particular does not converge here. We make no
assertion beyond this, and none is needed: the point is only that the
naive absolutely convergent periodisation argument is not available,
because $V$ decays only like $|y|^{-1}$ with a non-oscillating leading
profile whose spherical mean, $\tfrac43v_0/r$, is sign-definite in the
$v_0$ direction. If a periodisation formula
is wanted purely for orientation, the correct \emph{renormalised} one,
\[
u_N(x)=N\Big[V(Nx)+\sum_{m\neq0}\big(V(N(x+2\pi m))-V(2\pi Nm)\big)\Big]
+\text{const},
\]
does converge absolutely, the bracketed differences being $O(|m|^{-2})$
by Theorem~\ref{app:farfield}. \textbf{Nothing in this appendix uses any
periodisation identity}: \S\ref{app:riemann} replaces it by an exact
lattice-Riemann-sum identity, which is elementary, quantitative, and
additionally delivers the derivative statement. This remark, and with it
Theorem~\ref{app:farfield}, is inert for the proof of
Theorem~\ref{capacity}.
\end{remark}

\subsection{The inner rescaling: an exact lattice-Riemann identity with
an explicit rate}
\label{app:riemann}

For $\Phi:\R^3\to\mathbb C^d$ put
\[
\mathcal R_N[\Phi]:=N^{-3}\sum_{k\in\Z^3}\Phi(k/N)
-\int_{\R^3}\Phi(\xi)\,d\xi ,
\]
whenever both terms converge absolutely.

\begin{lemma}[exact inner rescaling]\label{app:innerexact}
For every $y\in\R^3$ and every integer $N\ge2$,
\[
u_N(y/N)=N\big(V(y)+E_N(y)\big),\qquad
\partial_iu_{N,j}(y/N)=N^2\big(\partial_iV_j(y)+E'_{N,ij}(y)\big),
\]
where, with $\Phi_y(\xi):=F(\xi)e^{i\xi\cdot y}$ and
$\Phi'_{y,ij}(\xi):=i\xi_iF_j(\xi)e^{i\xi\cdot y}$ (both set to $0$ at
$\xi=0$),
\[
E_N(y)=\mathcal R_N[\Phi_y],\qquad
E'_{N,ij}(y)=\mathcal R_N[\Phi'_{y,ij}] .
\]
These are identities, not approximations.
\end{lemma}

\begin{proof}
Both series are finite (support $|k|\le N$) and both integrals converge
absolutely by Theorem~\ref{app:Vprops}(1),(3). By \eqref{app:scaling},
\[
u_N(y/N)=\sum_{k\in\Z^3}\hat u_N(k)e^{i(k/N)\cdot y}
=\sum_{k}N^{-2}F(k/N)e^{i(k/N)\cdot y}
=N\cdot N^{-3}\sum_k\Phi_y(k/N),
\]
which equals $N\big(\int\Phi_y+\mathcal R_N[\Phi_y]\big)
=N\big(V(y)+E_N(y)\big)$, using $V(y)=\int\Phi_y$
(Definition~\ref{app:Vdef}) and $\Phi_y(0)=F(0)=0$. For the gradient,
$\partial_iu_{N,j}(x)=\sum_kik_i\hat u_{N,j}(k)e^{ik\cdot x}$; at
$x=y/N$ this is
\[
\sum_kiN^{-2}k_iF_j(k/N)e^{i(k/N)\cdot y}
=N^{-1}\sum_ki(k_i/N)F_j(k/N)e^{i(k/N)\cdot y}
=N^{2}\cdot N^{-3}\sum_k\Phi'_{y,ij}(k/N),
\]
and $\int\Phi'_{y,ij}=\partial_iV_j(y)$ by
Theorem~\ref{app:Vprops}(3).
\end{proof}

Everything now rests on bounding a Riemann-sum defect for an explicit,
compactly supported integrand with a single integrable algebraic
singularity.

\begin{theorem}[$C^1_{\rm loc}$ convergence with explicit rate]
\label{app:rate}
There is an absolute constant $A$ (the proof gives $A=3000$) such that
for every real $R\ge1$ and every integer $N\ge8R$,
\[
\varepsilon_N:=\sup_{|y|\le R}|E_N(y)|
\ \le\ A\,\|v_0\|\,\frac{c_\chi(1+\log_2N)+R}{N},
\qquad
\varepsilon'_N:=\sup_{|y|\le R}|E'_N(y)|
\ \le\ A\,\|v_0\|\,\frac{c_\chi+R}{N} .
\]
In particular $\varepsilon_N+\varepsilon'_N=O(N^{-1}\log N)\to0$ as
$N\to\infty$, for each fixed $R$.
\end{theorem}

\begin{proof}
Partition $\R^3$ into the half-open cubes
$Q_k=\frac kN+[0,\frac1N)^3$, $k\in\Z^3$, of side $1/N$, volume
$N^{-3}$ and diameter $\sqrt3/N$. For any $\Phi$ for which both the sum
and the integral converge absolutely,
\begin{equation}\label{app:cellwise}
\mathcal R_N[\Phi]=\sum_{k\in\Z^3}\int_{Q_k}
\big[\Phi(k/N)-\Phi(\xi)\big]\,d\xi ,
\end{equation}
since $|Q_k|=N^{-3}$. Absolute convergence holds for
$\Phi=\Phi_y$ and $\Phi=\Phi'_{y,ij}$: both are supported in
$\{|\xi|\le1\}$ and dominated there by $\|v_0\||\xi|^{-2}$ respectively
$\|v_0\||\xi|^{-1}$, which are integrable on $\R^3$; and the lattice
sums are finite sums, bounded by
$N^{-3}\sum_{0<|k|\le N}\|v_0\|N^{2}|k|^{-2}
=\|v_0\|N^{-1}S_N\le128\|v_0\|$ and by
$N^{-3}\sum_{0<|k|\le N}\|v_0\|N|k|^{-1}
\le\|v_0\|N^{-2}\cdot N\,S_N\le128\|v_0\|$ respectively, using
$|k|^{-1}\le N|k|^{-2}$ for $|k|\le N$ and Lemma~\ref{app:upper}.

\emph{Pointwise bounds.} On $0<|\xi|\le1$ we have
$|F(\xi)|\le\|v_0\||\xi|^{-2}$, and by the product rule applied to
$F=\chi(|\xi|)|\xi|^{-2}P_\xi v_0$,
\[
|\nabla F(\xi)|\ \le\
\underbrace{c_\chi\|v_0\||\xi|^{-2}}_{\nabla\chi(|\xi|)}
+\underbrace{2\|v_0\||\xi|^{-3}}_{\chi\,\nabla(|\xi|^{-2})}
+\underbrace{4\|v_0\||\xi|^{-3}}_{\chi|\xi|^{-2}\nabla(P_\xi v_0)}
\ \le\ 7c_\chi\|v_0\|\,|\xi|^{-3},
\]
where we used $c_\chi\ge1$, $|\xi|\le1$, $|\nabla(|\xi|^{-2})|
=2|\xi|^{-3}$, and, from
$\partial_a\big(\xi_i\xi_j/|\xi|^2\big)
=(\delta_{ai}\xi_j+\delta_{aj}\xi_i)/|\xi|^2
-2\xi_i\xi_j\xi_a/|\xi|^4$, the honest bound
$|\nabla(P_\xi v_0)|\le4\|v_0\|/|\xi|$. Consequently, for $|y|\le R$,
\begin{equation}\label{app:phibounds}
|\Phi_y|\le\frac{\|v_0\|}{|\xi|^{2}},\quad
|\nabla\Phi_y|\le\frac{7c_\chi\|v_0\|}{|\xi|^{3}}+\frac{R\|v_0\|}{|\xi|^{2}},
\qquad
|\Phi'_y|\le\frac{\|v_0\|}{|\xi|},\quad
|\nabla\Phi'_y|\le\frac{8c_\chi\|v_0\|}{|\xi|^{2}}+\frac{R\|v_0\|}{|\xi|} .
\end{equation}
(For the last one, $\nabla(i\xi_iF_je^{i\xi y})$ has the three terms
$|F|\le\|v_0\||\xi|^{-2}$, $|\xi||\nabla F|\le7c_\chi\|v_0\||\xi|^{-2}$
and $R|\xi||F|\le R\|v_0\||\xi|^{-1}$.)

\emph{The core.} Let
$\mathcal K_0:=\{k:\ Q_k\cap\{|\xi|\le4\sqrt3/N\}\neq\emptyset\}$. Any
such $Q_k$ lies in $\{|\xi|\le5\sqrt3/N\}$ and has $|k|\le5\sqrt3<9$,
i.e.\ $|k|\le8$. Note that $k=0$ is in the core, and that although
$\Phi_y(0)=0$ by convention while $\int_{Q_0}\Phi_y\neq0$, both are
bounded by the two terms below, so no separate treatment is needed. The
core contributes to \eqref{app:cellwise} at most
\[
\int_{|\xi|\le5\sqrt3/N}|\Phi_y|\,d\xi
+N^{-3}\!\!\sum_{0<|k|\le8}\!|\Phi_y(k/N)|
\ \le\ \|v_0\|\Big(\frac{4\pi\cdot5\sqrt3}{N}
+\frac1N\!\!\sum_{0<|k|\le8}\!\frac1{|k|^{2}}\Big)
\ \le\ \frac{1133\,\|v_0\|}{N},
\]
using $\int_{|\xi|\le a}|\xi|^{-2}d\xi=4\pi a$, then
$|\Phi_y(k/N)|\le\|v_0\|N^2|k|^{-2}$, and finally
$\sum_{0<|k|\le8}|k|^{-2}\le S_8\le1024$ by Lemma~\ref{app:upper}
($4\pi\cdot5\sqrt3\le109$, and $109+1024=1133$).

\emph{Dyadic annuli.} Put $J:=\lfloor\log_2\big(N/(4\sqrt3)\big)\rfloor$,
which is $\ge0$ because $N\ge8R\ge8>4\sqrt3$. Let $k\notin\mathcal K_0$
and set $r_k:=\inf_{\xi\in \overline{Q_k}}|\xi|$, so $r_k>4\sqrt3/N>0$.
If $r_k>1$ then $\Phi_y$ vanishes identically on $\overline{Q_k}$ and
that cell contributes $0$; otherwise choose the unique integer $j(k)$
with $2^{-j(k)-1}<r_k\le2^{-j(k)}$. Then $j(k)\ge0$, and
$2^{-j(k)}\ge r_k>4\sqrt3/N$ forces $2^{j(k)}<N/(4\sqrt3)$, i.e.\
$0\le j(k)\le J$; so the annuli
$A_j:=\{2^{-j-1}<|\xi|\le2^{-j}\}$, $0\le j\le J$, together with the
core, do cover $\operatorname{supp}\Phi_y$.

Fix $j$ and consider the cells with $j(k)=j$. They are pairwise
disjoint. Each satisfies $\overline{Q_k}\subset\{|\xi|\ge2^{-j-1}\}$ and
$\overline{Q_k}\subset\{|\xi|\le r_k+\sqrt3/N\le\tfrac54 2^{-j}\}$,
because $\sqrt3/N\le\tfrac14 2^{-j}$ (which follows from
$2^{-j}\ge2^{-J}\ge4\sqrt3/N$). Hence all of them lie in the fixed
annulus $B_j:=\{2^{-j-1}\le|\xi|\le\tfrac54 2^{-j}\}$, whose volume is
$\frac{4\pi}3\big((\tfrac54)^3-(\tfrac12)^3\big)2^{-3j}\le9\cdot2^{-3j}$;
therefore $\sum_{j(k)=j}|Q_k|\le9\cdot2^{-3j}$.

Since $\overline{Q_k}$ is convex, avoids the origin, and $\Phi_y$ is
$C^1$ there, the mean value inequality gives
$\big|\int_{Q_k}[\Phi_y(k/N)-\Phi_y(\xi)]d\xi\big|
\le\frac{\sqrt3}{N}\sup_{\overline{Q_k}}|\nabla\Phi_y|\cdot|Q_k|$. By
\eqref{app:phibounds} and $|\xi|\ge2^{-j-1}$ on $\overline{Q_k}$,
\[
\sup_{\overline{Q_k}}|\nabla\Phi_y|
\le 7c_\chi\|v_0\|2^{3(j+1)}+R\|v_0\|2^{2(j+1)}
=56c_\chi\|v_0\|8^{j}+4R\|v_0\|4^{j}.
\]
Summing over the cells with $j(k)=j$,
\[
\sum_{j(k)=j}\Big|\int_{Q_k}[\Phi_y(k/N)-\Phi_y]\Big|
\le\frac{\sqrt3}{N}\big(56c_\chi\|v_0\|8^{j}+4R\|v_0\|4^{j}\big)
\cdot9\cdot2^{-3j}
\le\frac{\|v_0\|}{N}\big(873\,c_\chi+63\,R\,2^{-j}\big).
\]
Summing over $0\le j\le J\le\log_2N$ and adding the core bound,
\[
|E_N(y)|\le\frac{\|v_0\|}{N}
\Big(1133+873\,c_\chi(1+\log_2N)+126R\Big)
\le\frac{3000\,\|v_0\|}{N}\big(c_\chi(1+\log_2N)+R\big),
\]
using $c_\chi\ge1$ and $1+\log_2N\ge1$. This is the first display.

\emph{The gradient.} Repeat verbatim with $\Phi'_{y,ij}$. The core
contributes at most
\[
\int_{|\xi|\le5\sqrt3/N}\frac{\|v_0\|}{|\xi|}\,d\xi
+N^{-3}\!\!\sum_{0<|k|\le8}\!\frac{\|v_0\|N}{|k|}
=\|v_0\|\Big(\frac{2\pi\cdot75}{N^{2}}
+\frac1{N^{2}}\!\!\sum_{0<|k|\le8}\!\frac1{|k|}\Big)
\le\frac{1816\,\|v_0\|}{N^{2}}\le\frac{1816\,\|v_0\|}{N},
\]
where $\int_{|\xi|\le a}|\xi|^{-1}d\xi=2\pi a^2$ and, by
Lemma~\ref{app:shell},
$\sum_{0<|k|\le8}|k|^{-1}\le\sum_{j=0}^{2}64\cdot8^{j}2^{-j}=1344$
(and $2\pi\cdot75\le472$, $472+1344=1816$). On the annuli,
$\sup_{\overline{Q_k}}|\nabla\Phi'_y|
\le8c_\chi\|v_0\|2^{2(j+1)}+R\|v_0\|2^{j+1}
=32c_\chi\|v_0\|4^{j}+2R\|v_0\|2^{j}$, so the level-$j$ contribution is
at most
\[
\frac{\sqrt3}{N}\big(32c_\chi\|v_0\|4^{j}+2R\|v_0\|2^{j}\big)
\cdot9\cdot2^{-3j}
\le\frac{\|v_0\|}{N}\big(499\,c_\chi2^{-j}+32\,R\,4^{-j}\big),
\]
and now the geometric series converge \emph{without} a logarithm,
summing over all $j\ge0$ to at most
$\frac{\|v_0\|}{N}(998c_\chi+43R)$. Adding the core,
$|E'_N(y)|\le\frac{\|v_0\|}{N}(1816+998c_\chi+43R)
\le\frac{3000\|v_0\|}{N}(c_\chi+R)$. The absence of a logarithm is
structural: the singularity of $\Phi'_y$ is only $|\xi|^{-1}$, the extra
factor $\xi$ taming the origin. The derivative estimate is thus not a
consequence of the $C^0$ one, but it is the easier of the two.
\end{proof}

\subsection{The concentrated test field and the duality lower bound}
\label{app:duality}

\begin{definition}\label{app:testfield}
Let $\Psi\in C_c^\infty(\R^3;\R^3)$ be \emph{real}, with
$\nabla\cdot\Psi=0$, $\operatorname{supp}\Psi\subset\{|y|\le R\}$ for
some $R\ge1$, and $\Psi\not\equiv0$. For $N>R/\pi$ define on $\T^3$
\[
\psi_N(x):=N^{3/2}\sum_{m\in\Z^3}\Psi\big(N(x+2\pi m)\big).
\]
\end{definition}

\begin{lemma}[properties of $\psi_N$]\label{app:psiprops}
Let $N>R/\pi$. Then $\psi_N$ is a well-defined real $C^\infty$ vector
field on $\T^3$; it is divergence-free and has zero mean; on the
fundamental domain $[-\pi,\pi)^3$ exactly the term $m=0$ is nonzero, so
$\operatorname{supp}\psi_N\subset\{|x|\le R/N\}$ there; and
\[
\|\psi_N\|_{L^2(\T^3)}=(2\pi)^{-3/2}\|\Psi\|_{L^2(\R^3)},
\]
which is independent of $N$.
\end{lemma}

\begin{proof}
Local finiteness: $\Psi(N(x+2\pi m))\neq0$ forces
$|x+2\pi m|\le R/N<\pi$; but for $x\in[-\pi,\pi)^3$ and $m\neq0$,
$|x+2\pi m|\ge2\pi|m|_\infty-\pi\ge\pi$, a contradiction. Hence on a
neighbourhood of any point the sum has at most one nonzero term, so
$\psi_N$ is smooth, real, and $2\pi\Z^3$-periodic by construction.
Divergence-free: locally
$\nabla\cdot\psi_N(x)=N^{5/2}(\nabla\cdot\Psi)(N(x+2\pi m))=0$.
Zero mean: with $d\mu=(2\pi)^{-3}dx$ and the substitution $y=Nx$,
\[
\int_{\T^3}\psi_N\,d\mu
=(2\pi)^{-3}N^{3/2}\int_{\R^3}\Psi(Nx)\,dx
=(2\pi)^{-3}N^{-3/2}\int_{\R^3}\Psi ,
\]
and $\int_{\R^3}\Psi_i=\int_{\R^3}\nabla\cdot(y_i\Psi)=0$ since
$\nabla\cdot\Psi=0$ and $y_i\Psi\in C^\infty_c$. (No curl
representation of $\Psi$ is needed: zero mean is automatic for a
compactly supported divergence-free field.) Norm:
\[
\|\psi_N\|_{L^2(\T^3)}^2
=(2\pi)^{-3}\int_{|x|\le R/N}N^{3}|\Psi(Nx)|^2dx
=(2\pi)^{-3}N^{3}N^{-3}\|\Psi\|_{L^2(\R^3)}^2 ,
\]
which is the stated identity.
\end{proof}

\begin{theorem}[duality identity with controlled error]\label{app:pairing}
Let $\Psi$ be as in Definition~\ref{app:testfield} and put
\[
I_\Psi:=\int_{\R^3}\big(V\cdot\nabla V\big)\cdot\Psi\,dy .
\]
Then for every integer $N\ge\max(8R,32)$,
\[
\big\langle\PP(u_N\cdot\nabla u_N),\psi_N\big\rangle_{L^2(\T^3)}
=(2\pi)^{-3}N^{3/2}\big(I_\Psi+\mathcal E_N\big),
\]
where
\[
|\mathcal E_N|\ \le\ \|\Psi\|_{L^1(\R^3)}
\Big[\varepsilon_N\big(6\pi\|v_0\|+\varepsilon'_N\big)
+4\pi\|v_0\|\,\varepsilon'_N\Big]
\ =\ O\Big(\frac{\log N}{N}\Big)
\]
with $\varepsilon_N,\varepsilon'_N$ as in Theorem~\ref{app:rate}.
\end{theorem}

\begin{proof}
By Lemma~\ref{app:psiprops}, $\psi_N$ is zero-mean and divergence-free,
so $\hat\psi_N(0)=0$ and $k\cdot\hat\psi_N(k)=0$ for all $k$, whence
$\PP\psi_N=\psi_N$. Since $\PP$ is an orthogonal projection on
$L^2(\T^3)$, hence self-adjoint,
\[
\big\langle\PP(u_N\cdot\nabla u_N),\psi_N\big\rangle
=\big\langle u_N\cdot\nabla u_N,\PP\psi_N\big\rangle
=\big\langle u_N\cdot\nabla u_N,\psi_N\big\rangle .
\]
Both fields are real. Because $N\ge8R$ and $R\ge1$ we have
$R/N\le\tfrac18<\pi$, so Lemma~\ref{app:psiprops} applies and
$\psi_N(x)=N^{3/2}\Psi(Nx)$ on the fundamental domain, supported in
$\{|x|\le R/N\}$. Using $d\mu=(2\pi)^{-3}dx$ and substituting $x=y/N$
($dx=N^{-3}dy$),
\[
\big\langle u_N\cdot\nabla u_N,\psi_N\big\rangle
=(2\pi)^{-3}N^{3/2}N^{-3}\int_{|y|\le R}
(u_N\cdot\nabla u_N)(y/N)\cdot\Psi(y)\,dy .
\]
By Lemma~\ref{app:innerexact}, for $|y|\le R$,
\[
(u_N\cdot\nabla u_N)_j(y/N)
=u_{N,i}(y/N)\,\partial_iu_{N,j}(y/N)
=N^{3}\big(V_i(y)+E_{N,i}(y)\big)\big(\partial_iV_j(y)+E'_{N,ij}(y)\big),
\]
so the factor $N^{-3}N^{3}$ cancels and
\[
\big\langle u_N\cdot\nabla u_N,\psi_N\big\rangle
=(2\pi)^{-3}N^{3/2}\big(I_\Psi+\mathcal E_N\big),
\qquad
\mathcal E_N=\int_{|y|\le R}
\big[E_{N,i}\partial_iV_j+V_iE'_{N,ij}+E_{N,i}E'_{N,ij}\big]\Psi_j\,dy ,
\]
where we used $\operatorname{supp}\Psi\subset\{|y|\le R\}$ to write the
main term as $I_\Psi$. Bounding each factor by its supremum and using
Theorem~\ref{app:Vprops}(2),(3) --- $\|V\|_\infty\le4\pi\|v_0\|$,
$\|\nabla V\|_\infty\le6\pi\|v_0\|$ --- gives the stated estimate;
$\varepsilon_N,\varepsilon'_N=O(N^{-1}\log N)$ by
Theorem~\ref{app:rate}, $R$ and $\Psi$ being fixed.
\end{proof}

\begin{proposition}[the $N^3$ lower bound, given a good test field]
\label{app:N3}
Suppose some $\Psi$ as in Definition~\ref{app:testfield} has
$I_\Psi\neq0$. Then
\[
\big\|\PP(u_N\cdot\nabla u_N)\big\|_{L^2(\T^3)}^2
\ \ge\ \frac{|I_\Psi+\mathcal E_N|^2}
{(2\pi)^{3}\,\|\Psi\|_{L^2(\R^3)}^{2}}\;N^{3}
\qquad(N\ge\max(8R,32)),
\]
and consequently
$\|\PP(u_N\cdot\nabla u_N)\|_{2}^2\ge c_0N^3$ for all $N\ge N_*$, where
\[
c_0=\frac{|I_\Psi|^2}{2\,(2\pi)^{3}\,\|\Psi\|_{L^2(\R^3)}^{2}},
\qquad
N_*=\min\Big\{N_0\ge\max(8R,32):\ |\mathcal E_N|\le(1-2^{-1/2})|I_\Psi|
\ \ \forall N\ge N_0\Big\},
\]
which is finite because $\mathcal E_N\to0$.
\end{proposition}

\begin{proof}
By Cauchy--Schwarz in $L^2(\T^3)$,
\[
\big\|\PP(u_N\cdot\nabla u_N)\big\|_{2}
\ \ge\ \frac{\big|\langle\PP(u_N\cdot\nabla u_N),\psi_N\rangle\big|}
{\|\psi_N\|_{2}} .
\]
Insert Theorem~\ref{app:pairing} and Lemma~\ref{app:psiprops}: the
right-hand side equals
\[
\frac{(2\pi)^{-3}N^{3/2}|I_\Psi+\mathcal E_N|}
{(2\pi)^{-3/2}\|\Psi\|_{L^2(\R^3)}}
=(2\pi)^{-3/2}N^{3/2}\frac{|I_\Psi+\mathcal E_N|}{\|\Psi\|_{L^2(\R^3)}} .
\]
Squaring gives the display. For $N\ge N_*$ we have
$|I_\Psi+\mathcal E_N|\ge2^{-1/2}|I_\Psi|$, whence
$|I_\Psi+\mathcal E_N|^2\ge\tfrac12|I_\Psi|^2$ and the constant $c_0$ as
stated.
\end{proof}

Everything is now reduced to the single question of whether some
admissible $\Psi$ has $I_\Psi\neq0$.

\subsection{Non-degeneracy of the continuum nonlinearity}
\label{app:nondeg}

\subsubsection*{Reduction to a symbol statement}

\begin{lemma}[symbol of the nonlinearity]\label{app:symbol}
$V\cdot\nabla V=\nabla\cdot(V\otimes V)$ has the representation
\[
(V\cdot\nabla V)(y)=\int_{\R^3}M(\zeta)e^{i\zeta\cdot y}\,d\zeta,
\qquad
M_i(\zeta)=i\,\zeta_j\,T_{ij}(\zeta),\qquad T_{ij}:=F_i*F_j ,
\]
where $T$ is real, symmetric, supported in $\{|\zeta|\le2\}$, and
$T,M\in L^1(\R^3)\cap L^2(\R^3)$. Consequently
$\PP(V\cdot\nabla V)\in L^2(\R^3)$ with symbol $P_\zeta M(\zeta)$, and
\[
\PP(V\cdot\nabla V)\equiv0
\quad\Longleftrightarrow\quad
\zeta\times\big(T(\zeta)\zeta\big)=0\ \text{ for a.e. }\zeta .
\]
\end{lemma}

\begin{proof}
$V=\widetilde F$ with $F\in L^1$, and $\partial_jV_i=\widetilde{i\xi_jF_i}$
with $\xi F\in L^1$ (Theorem~\ref{app:Vprops}(1),(3)). The pointwise
product of inverse transforms of $L^1$ functions is the inverse
transform of the convolution, so
$V_j\partial_jV_i=\widetilde{F_j*(i\,\cdot_j F_i)}$, i.e.\ the symbol of
$V\cdot\nabla V$ is
$i\int F_j(\zeta-\eta)\,\eta_j\,F_i(\eta)\,d\eta$. Since
$(\zeta-\eta)\cdot F(\zeta-\eta)=0$ identically (Definition~\ref{app:Vdef}),
we may replace $\eta_jF_j(\zeta-\eta)$ by $\zeta_jF_j(\zeta-\eta)$
inside the integral, obtaining
$M_i(\zeta)=i\zeta_j(F_i*F_j)(\zeta)=i\zeta_jT_{ij}(\zeta)$.

$T$ is real because $F$ is real, symmetric because convolution is
commutative, and supported in $\{|\zeta|\le2\}$ because
$\operatorname{supp}F\subset\{|\xi|\le1\}$. Integrability: by
Theorem~\ref{app:Vprops}(1), $F\in L^{4/3}$; Young's inequality with
$\frac1r=\frac2{4/3}-1=\frac12$ gives $T=F*F\in L^{2}$, and since $T$
has compact support, $T\in L^1\cap L^2$; the same holds for
$M=i\zeta T$, as $|\zeta|\le2$ on the support. By Plancherel,
$V\cdot\nabla V\in L^2(\R^3)$ and, since $|P_\zeta M|\le|M|$,
$\PP(V\cdot\nabla V)\in L^2(\R^3)$ with symbol $P_\zeta M$. Finally
$\PP(V\cdot\nabla V)=0$ in $L^2$ iff $P_\zeta M(\zeta)=0$ for a.e.
$\zeta$, and for a vector $w$ one has $P_\zeta w=0$ iff $w\parallel\zeta$
iff $\zeta\times w=0$; here $w=i\,T(\zeta)\zeta$.
\end{proof}

\begin{remark}[classical input: density of smooth compactly supported
solenoidal fields]\label{app:classical-density}
The space $C^\infty_{c,\sigma}(\R^3;\R^3)$ of real, smooth, compactly
supported, divergence-free vector fields is dense in
$L^2_\sigma(\R^3;\R^3)$, the closed subspace of real $L^2$ fields that
are distributionally divergence-free; indeed $L^2_\sigma$ is by the
standard definition the $L^2$-closure of $C^\infty_{c,\sigma}$, and on
$\R^3$ the two descriptions agree. See, e.g.,
\cite{Temam,LadyzhenskayaBook,ConstantinFoias}.
\end{remark}

\begin{remark}[classical input: Helmholtz--Leray decomposition on $\R^3$]
\label{app:classical-leray}
For $f\in L^2(\R^3;\R^3)$ one has $f-\PP f=\nabla p$ for some
$p\in L^2_{\rm loc}(\R^3)$ with $\nabla p\in L^2$; equivalently, the
symbol of $f-\PP f$ is $\zeta(\zeta\cdot\hat f)/|\zeta|^2$. Same
references.
\end{remark}

\begin{lemma}[from non-degeneracy to a test field]\label{app:testexists}
If $\PP(V\cdot\nabla V)\not\equiv0$, then there exists a real
$\Psi\in C^\infty_c(\R^3;\R^3)$ with $\nabla\cdot\Psi=0$ and
$I_\Psi=\int_{\R^3}(V\cdot\nabla V)\cdot\Psi\neq0$; after enlarging $R$
if necessary, $\Psi$ satisfies Definition~\ref{app:testfield}.
\end{lemma}

\begin{proof}
By Lemma~\ref{app:symbol}, $\PP(V\cdot\nabla V)$ lies in $L^2(\R^3)$,
is real (its symbol $P_\zeta M$ satisfies the reality condition because
$F$ is real and even), is distributionally divergence-free (its symbol
is orthogonal to $\zeta$), and is nonzero by hypothesis. Hence it is a
nonzero element of the real Hilbert space $L^2_\sigma(\R^3)$. By
Remark~\ref{app:classical-density} there is a real
$\Psi\in C^\infty_{c,\sigma}(\R^3)$ with
$\langle\PP(V\cdot\nabla V),\Psi\rangle_{L^2(\R^3)}\neq0$. By
Remark~\ref{app:classical-leray},
$V\cdot\nabla V-\PP(V\cdot\nabla V)=\nabla p$ with $\nabla p\in L^2$, and
$\int\nabla p\cdot\Psi=-\int p\,\nabla\cdot\Psi=0$ because $\Psi$ is
compactly supported and divergence-free. Therefore
$I_\Psi=\langle\PP(V\cdot\nabla V),\Psi\rangle\neq0$. Choosing
$R\ge\max(1,\sup\{|y|:y\in\operatorname{supp}\Psi\})$ puts $\Psi$ in
Definition~\ref{app:testfield}.
\end{proof}

\subsubsection*{The $\chi$-independent leading singularity}

By rotational equivariance we may and do assume from now on that
$v_0=\|v_0\|e_3$: $F$ and $V$ are linear in $v_0$ and rotate with it,
$T$ and $\tau$ are quadratic, and the conclusion
$\PP(V\cdot\nabla V)\not\equiv0$ is rotation invariant.

\begin{definition}\label{app:tau}
$\tau_{ij}:=h_i*h_j$, with $h(\xi)=P_\xi v_0/|\xi|^2$ the uncut symbol
of Definition~\ref{app:Vdef}.
\end{definition}

\begin{lemma}[$\tau$ is well defined and homogeneous]\label{app:tauwd}
The integral $\tau_{ij}(\zeta)=\int_{\R^3}h_i(\eta)h_j(\zeta-\eta)\,d\eta$
converges absolutely for every $\zeta\neq0$, and $\tau$ is homogeneous of
degree $-1$ and smooth on $\R^3\setminus\{0\}$.
\end{lemma}

\begin{proof}
$|h(\xi)|\le\|v_0\||\xi|^{-2}$. Near $\eta=0$ the integrand is
$O(|\eta|^{-2})$ and near $\eta=\zeta$ it is $O(|\zeta-\eta|^{-2})$, both
locally integrable on $\R^3$; for $|\eta|\ge2|\zeta|$ it is
$O(|\eta|^{-4})$, integrable at infinity. Substituting
$\eta=\lambda\eta'$ and using $h(\lambda\eta)=\lambda^{-2}h(\eta)$ gives
$\tau(\lambda\zeta)=\lambda^{-1}\tau(\zeta)$ for $\lambda>0$. Smoothness
away from $0$ follows from differentiating under the integral sign on
the region where the two singularities are separated.
\end{proof}

\begin{theorem}[closed form of the leading obstruction]
\label{app:tauformula}
With $v_0=\|v_0\|e_3$, for every $\zeta\neq0$,
\[
\zeta\times\big(\tau(\zeta)\zeta\big)
=\frac{3\pi^3}{8}\,\|v_0\|^2\,\frac{\zeta_3}{|\zeta|}\,
\big(\zeta\times e_3\big).
\]
In particular $\zeta\times(\tau(\zeta)\zeta)\neq0$ whenever
$\zeta_3\neq0$ and $\zeta\not\parallel e_3$.
\end{theorem}

\begin{proof}
Both sides are quadratic in $\|v_0\|$, so we normalise $\|v_0\|=1$,
$v_0=e_3$. By Theorem~\ref{app:farfield}, $\widetilde h=\pi^2W$ with
\[
W_i(y)=\frac{\delta_{i3}}{r}+\frac{y_3y_i}{r^3},\qquad r=|y| .
\]
$W_i$ is locally integrable, homogeneous of degree $-1$ and smooth away
from the origin, so Remark~\ref{app:classical-homog} applies with
$a=b=1$ and gives
\[
\tau_{ij}=h_i*h_j=(2\pi)^{-3}\,\mathcal F\big[\widetilde h_i\widetilde h_j\big]
=(2\pi)^{-3}\pi^4\,\mathcal F\big[W_iW_j\big],
\]
where we used $\mathcal F\widetilde h_i=(2\pi)^3h_i$ and the evenness of
$h$. (Both sides are homogeneous of degree $-1$ and locally integrable,
so no distribution supported at the origin can be added.) Expanding,
\[
W_iW_j=\frac{\delta_{i3}\delta_{j3}}{r^{2}}
+\frac{\delta_{i3}y_3y_j+\delta_{j3}y_3y_i}{r^{4}}
+\frac{y_3^{2}y_iy_j}{r^{6}} .
\]
Each bracket is locally integrable and homogeneous of degree $-2$, so by
Remark~\ref{app:classical-riesz} its transform is an honest homogeneous
distribution of degree $-1$, given by
\[
\mathcal F\Big[\frac1{r^{2}}\Big]=\frac{2\pi^{2}}{s},
\qquad
\mathcal F\Big[\frac{y_iy_j}{r^{4}}\Big]
=\pi^{2}\Big(\frac{\delta_{ij}}{s}-\frac{\zeta_i\zeta_j}{s^{3}}\Big),
\qquad
\mathcal F\Big[\frac{y_3^{2}y_iy_j}{r^{6}}\Big]
=\frac{\pi^{2}}{12}\,\partial_3^{2}\partial_i\partial_j\,s^{3},
\]
where $s:=|\zeta|$. (Indeed $\mathcal F[r^{-2}]=c_{3,2}s^{-1}$;
$\mathcal F[y_iy_jr^{-4}]=-\partial_i\partial_j(c_{3,4}s)
=\pi^2\partial_i\partial_js=\pi^2(\delta_{ij}/s-\zeta_i\zeta_j/s^3)$;
$\mathcal F[y_3^2y_iy_jr^{-6}]
=\partial_3^2\partial_i\partial_j(c_{3,6}s^3)$. Trace check on the
middle formula: $\pi^2(3/s-1/s)=2\pi^2/s=\mathcal F[r^{-2}]$, as it must
be.) Direct differentiation of $s^3$ gives
\[
\partial_3^{2}\partial_i\partial_j s^{3}
=3\Big[\frac{\delta_{ij}+2\delta_{i3}\delta_{j3}}{s}
-\frac{\zeta_3^{2}\delta_{ij}
+2\zeta_3(\delta_{i3}\zeta_j+\delta_{j3}\zeta_i)+\zeta_i\zeta_j}{s^{3}}
+\frac{3\zeta_i\zeta_j\zeta_3^{2}}{s^{5}}\Big].
\]
Assembling the three pieces (the middle one contributing
$\pi^2[2\delta_{i3}\delta_{j3}/s-\zeta_3(\delta_{i3}\zeta_j
+\delta_{j3}\zeta_i)/s^3]$) yields
$\mathcal F[W_iW_j]=\pi^{2}G_{ij}$ with
\[
G_{ij}=\frac{\delta_{ij}}{4s}
+\frac{9\,\delta_{i3}\delta_{j3}}{2s}
-\frac{\zeta_3^{2}\delta_{ij}}{4s^{3}}
-\frac{3\zeta_3(\delta_{i3}\zeta_j+\delta_{j3}\zeta_i)}{2s^{3}}
-\frac{\zeta_i\zeta_j}{4s^{3}}
+\frac{3\zeta_i\zeta_j\zeta_3^{2}}{4s^{5}} ,
\]
so that $\tau_{ij}=(2\pi)^{-3}\pi^{4}\cdot\pi^{2}G_{ij}
=\frac{\pi^{3}}{8}G_{ij}$.

Contract with $\zeta_j$, using $\zeta_j\zeta_j=s^2$. The terms
$\frac{\delta_{ij}}{4s}\zeta_j=\frac{\zeta_i}{4s}$ and
$-\frac{\zeta_i\zeta_j}{4s^3}\zeta_j=-\frac{\zeta_i}{4s}$ cancel. The
$\delta_{i3}\zeta_3/s$ coefficients are $\frac92$ from the second term
and $-\frac32$ from the fourth, total $3$. The $\zeta_i\zeta_3^2/s^3$
coefficients are $-\frac14$ (third term), $-\frac32$ (fourth term) and
$+\frac34$ (sixth term), total $-1$. Hence
\[
G(\zeta)\zeta=\frac{3\zeta_3}{s}\,e_3-\frac{\zeta_3^{2}}{s^{3}}\,\zeta ,
\qquad\text{so}\qquad
\zeta\times\big(G(\zeta)\zeta\big)=\frac{3\zeta_3}{s}\,(\zeta\times e_3),
\]
the second term being parallel to $\zeta$. Multiplying by
$\frac{\pi^3}{8}$ and restoring $\|v_0\|^2$ gives the claim. Finally
$\zeta\times e_3=0$ exactly when $\zeta\parallel e_3$.
\end{proof}

\begin{remark}[physical-space cross-check, free of any regularisation]
\label{app:crosscheck}
Theorem~\ref{app:tauformula} passes through the finite-part
continuations of Remark~\ref{app:classical-riesz}. The following
elementary computation confirms it independently, with no
regularisation and no Fourier transform, and is worth recording because
it identifies what the obstruction \emph{is}: the statement that the
Oseen field is not a stationary Euler flow.

For $W(y)=\frac{e_3}{r}+\frac{y_3y}{r^{3}}$ on $\R^3\setminus\{0\}$,
\[
\nabla\cdot W=0,\qquad
\omega:=\nabla\times W=\frac{2\,(e_3\times y)}{r^{3}},\qquad
\nabla\times\big(W\times\omega\big)
=\frac{12\,y_3\,(e_3\times y)}{r^{6}}\ \not\equiv0 .
\]
\emph{Proof.} $\nabla\cdot(e_3/r)=e_3\cdot\nabla r^{-1}=-y_3/r^3$, and
$\partial_i(y_3y_ir^{-3})=y_3r^{-3}+3y_3r^{-3}-3y_3r^{-3}=y_3/r^3$, so
$\nabla\cdot W=0$. Next
$\nabla\times(r^{-1}e_3)=\nabla(r^{-1})\times e_3=-r^{-3}y\times e_3
=r^{-3}(e_3\times y)$, while
$[\nabla\times(y_3y\,r^{-3})]_a
=\epsilon_{abi}\partial_b(y_3y_ir^{-3})
=\epsilon_{abi}(\delta_{b3}y_i+y_3\delta_{bi})r^{-3}
-3\epsilon_{abi}y_3y_iy_br^{-5}
=\epsilon_{a3i}y_ir^{-3}=(e_3\times y)_ar^{-3}$, the last two terms
vanishing by antisymmetry. Adding gives $\omega$. Then, using
$e_3\times(e_3\times y)=y_3e_3-y$ and $y\times(e_3\times y)=r^2e_3-y_3y$,
\[
W\times\omega
=\frac{2(y_3e_3-y)}{r^{4}}+\frac{2y_3e_3}{r^{4}}-\frac{2y_3^{2}y}{r^{6}}
=\frac{4y_3e_3}{r^{4}}-\frac{2y}{r^{4}}-\frac{2y_3^{2}y}{r^{6}} .
\]
Now curl each piece. For any radial $f$, $\nabla\times(f(r)y)
=\nabla f\times y=f'(r)\hat y\times y=0$, so $-2y/r^4$ contributes
nothing. With $g=-2y_3^{2}/r^{6}$ one has
$\nabla g=-4y_3e_3/r^{6}+12y_3^{2}y/r^{8}$, so
$\nabla\times(gy)=\nabla g\times y=-4y_3(e_3\times y)/r^{6}$. With
$f=4y_3/r^{4}$ one has $\nabla f=4e_3/r^{4}-16y_3y/r^{6}$, so
$\nabla\times(fe_3)=\nabla f\times e_3
=-16y_3(y\times e_3)/r^{6}=16y_3(e_3\times y)/r^{6}$. The total is
$12y_3(e_3\times y)/r^{6}$. Since
$W\cdot\nabla W=\nabla\frac{|W|^2}{2}-W\times\omega$, it follows that
$\nabla\times(W\cdot\nabla W)\not\equiv0$, i.e.\ $W\cdot\nabla W$ is not
a gradient, i.e.\ $\PP(W\cdot\nabla W)\neq0$. $\square$

\emph{Consistency with Theorem~\ref{app:tauformula}.} The curl of
$\PP(\widetilde h\cdot\nabla\widetilde h)$ has symbol
$i\zeta\times\big(iP_\zeta\tau\zeta\big)=-\zeta\times(\tau\zeta)$, using
$\zeta\times P_\zeta A=\zeta\times A$. With $\|v_0\|=1$,
Theorem~\ref{app:tauformula} makes this
$-\frac{3\pi^{3}}{8}\zeta_3(\zeta\times e_3)/|\zeta|$. Writing its
$a$-component as $-\frac{3\pi^3}{8}\epsilon_{ab3}\zeta_b\zeta_3/|\zeta|$
and using $\widetilde{|\zeta|^{-1}}=c_{3,1}r^{-2}=4\pi r^{-2}$ together
with $\widetilde{\zeta_b\zeta_3G}=-\partial_b\partial_3\widetilde G$,
\[
\widetilde{\Big[\frac{\zeta_b\zeta_3}{|\zeta|}\Big]}
=-4\pi\,\partial_b\partial_3r^{-2}
=\frac{8\pi\delta_{b3}}{r^{4}}-\frac{32\pi y_3y_b}{r^{6}} ;
\]
since $\epsilon_{ab3}\delta_{b3}=0$, only the second term survives,
leaving $-32\pi y_3(y\times e_3)_a/r^{6}
=+32\pi y_3(e_3\times y)_a/r^{6}$. Multiplying by
$-\frac{3\pi^{3}}{8}$ gives $-12\pi^{4}y_3(e_3\times y)/r^{6}$. On the
other hand $\widetilde h=\pi^2W$ gives
$\widetilde h\cdot\nabla\widetilde h=\pi^{4}W\cdot\nabla W$ and hence
$\nabla\times(\widetilde h\cdot\nabla\widetilde h)
=-\pi^{4}\nabla\times(W\times\omega)
=-12\pi^{4}y_3(e_3\times y)/r^{6}$. The two derivations agree exactly,
constants included.
\end{remark}

\subsubsection*{Transfer to the cutoff family}

\begin{lemma}[the cutoff perturbation is bounded near $\zeta=0$]
\label{app:perturb}
For all $|\zeta|\le\frac18$ and all $i,j$,
\[
\big|T_{ij}(\zeta)-\tau_{ij}(\zeta)\big|\ \le\ C_1:=64\pi\|v_0\|^{2},
\]
uniformly over all admissible $\chi$.
\end{lemma}

\begin{proof}
Since $F=\chi(|\cdot|)h$,
\[
T_{ij}(\zeta)-\tau_{ij}(\zeta)
=\int_{\R^3}h_i(\eta)h_j(\zeta-\eta)
\big[\chi(|\eta|)\chi(|\zeta-\eta|)-1\big]\,d\eta .
\]
The bracket vanishes whenever $|\eta|\le\frac12$ \emph{and}
$|\zeta-\eta|\le\frac12$, because $\chi\equiv1$ on $[0,\frac12]$. For
$|\zeta|\le\frac18$, the condition $|\eta|\le\frac14$ implies
$|\zeta-\eta|\le\frac18+\frac14=\frac38<\frac12$; hence the integrand is
supported in $\{|\eta|\ge\frac14\}$. There
$|\zeta-\eta|\ge|\eta|-\frac18\ge|\eta|-\frac{|\eta|}{2}
=\frac{|\eta|}{2}$. With $|h|\le\|v_0\||\cdot|^{-2}$ and
$|\chi\chi-1|\le1$,
\[
|T_{ij}-\tau_{ij}|
\le\|v_0\|^{2}\int_{|\eta|\ge1/4}\frac{d\eta}{|\eta|^{2}|\zeta-\eta|^{2}}
\le4\|v_0\|^{2}\int_{|\eta|\ge1/4}\frac{d\eta}{|\eta|^{4}}
=4\|v_0\|^{2}\cdot4\pi\int_{1/4}^{\infty}\frac{dr}{r^{2}}
=64\pi\|v_0\|^{2},
\]
which is $C_1$.
\end{proof}

\begin{theorem}[non-degeneracy]\label{app:vnondeg}
For every $v_0\in\R^3\setminus\{0\}$ and every admissible $\chi$,
\[
\PP_{\R^3}\big(V\cdot\nabla V\big)\ \not\equiv\ 0 .
\]
Explicitly, with $v_0=\|v_0\|e_3$, one has
$\zeta\times(T(\zeta)\zeta)\neq0$ --- hence $P_\zeta M(\zeta)\neq0$ ---
at every point of the open set of positive measure
\[
\mathcal U:=\Big\{\zeta\in\R^3:\ 0<|\zeta|<\tfrac{\pi^{2}}{2048},\
|\zeta_3|\ge\tfrac{|\zeta|}{2},\
|\zeta\times e_3|\ge\tfrac{|\zeta|}{2}\Big\}.
\]
\end{theorem}

\begin{proof}
Let $\zeta\in\mathcal U$. By Theorem~\ref{app:tauformula},
\[
\big|\zeta\times(\tau(\zeta)\zeta)\big|
=\frac{3\pi^{3}}{8}\|v_0\|^{2}\,
\frac{|\zeta_3|\,|\zeta\times e_3|}{|\zeta|}
\ \ge\ \frac{3\pi^{3}}{8}\|v_0\|^{2}\cdot\frac{|\zeta|}{4}
=\frac{3\pi^{3}}{32}\|v_0\|^{2}|\zeta| .
\]
Since $\frac{\pi^{2}}{2048}\approx4.8\times10^{-3}<\frac18$,
Lemma~\ref{app:perturb} applies, and the Frobenius norm of
$T(\zeta)-\tau(\zeta)$ is at most $3C_1$ (nine entries each at most
$C_1$). Hence
\[
\big|\zeta\times\big((T(\zeta)-\tau(\zeta))\zeta\big)\big|
\le|\zeta|\cdot3C_1|\zeta|=192\pi\|v_0\|^{2}|\zeta|^{2},
\]
and therefore
\[
\big|\zeta\times(T(\zeta)\zeta)\big|
\ \ge\ \|v_0\|^{2}|\zeta|
\Big(\frac{3\pi^{3}}{32}-192\pi|\zeta|\Big)\ >\ 0
\qquad\text{whenever }
|\zeta|<\frac{3\pi^{3}}{32\cdot192\pi}=\frac{3\pi^{2}}{6144}
=\frac{\pi^{2}}{2048},
\]
which holds on $\mathcal U$. The set $\mathcal U$ is open and nonempty
(it contains, e.g., all small $\zeta$ in the direction
$(1,0,1)/\sqrt2$), hence of positive measure. By
Lemma~\ref{app:symbol}, $P_\zeta M\neq0$ on a set of positive measure,
so $\PP(V\cdot\nabla V)\neq0$ in $L^2(\R^3)$.
\end{proof}

\begin{remark}[uniformity, and where the obstruction lives]
\label{app:uniformity}
The proof of Theorem~\ref{app:vnondeg} is uniform in $\chi$ and in
$v_0$: the factor $\|v_0\|^2$ cancels between the main term and the
perturbation, and the only properties of $\chi$ used are $0\le\chi\le1$
and $\chi\equiv1$ on $[0,\frac12]$ --- no smoothness of $\chi$ enters
\S\ref{app:nondeg} at all. Neither a specific $\chi$ nor any numerically
evaluated convolution integral is required: the obstruction is carried
entirely by the $\chi$-independent leading singularity of $F*F$ at
$\zeta=0$, which is the same for every admissible cutoff. Moreover
Theorem~\ref{app:tauformula} shows the obstruction degenerates exactly
on the two directions $\zeta\parallel v_0$ and $\zeta\perp v_0$, so the
two cone conditions in $\mathcal U$ are essentially optimal.

This uniformity does \textbf{not} propagate to the constants of
Theorem~\ref{app:main}. It is the \emph{qualitative} conclusion
$\PP(V\cdot\nabla V)\not\equiv0$ that holds uniformly in $(\chi,v_0)$,
and that is what licenses the quantifier ``for every admissible $\chi$
and every $v_0\neq0$'' there. The quantitative constants $c_0$ and
$N_*$ are extracted only after a test field $\Psi$ has been chosen for
the particular profile $V=V^{\chi,v_0}$
(Lemma~\ref{app:testexists}), and they depend on $(\chi,v_0)$; no
$(\chi,v_0)$-independent value of either is claimed, and neither is
effective (Remark~\ref{app:noneffective}).
\end{remark}

\subsection{Proof of Theorem~\ref{capacity}}
\label{app:final}

\begin{theorem}[capacity bound; restatement of Theorem~\ref{capacity}]
\label{app:main}
Let $v_0\in\R^3\setminus\{0\}$, let $\chi$ be admissible
(Definition~\ref{app:admissible}), and let $u_N=u_N^{\chi,v_0}$ be the
family \eqref{app:famdef}. Then:
\begin{enumerate}
\item[(1)] $u_N$ is a real, zero-mean, exactly divergence-free
trigonometric polynomial banded at $|k|\le N$, and the exact laws of
Theorem~\ref{app:exactlaws} hold;
\item[(2)] $\dfrac{N}{44544}\le N_0^2(u_N)\le\dfrac{64}{3}N$ for every
$N\ge32$;
\item[(3)] there exist $c_0=c_0(\chi,v_0)>0$ and
$N_*=N_*(\chi,v_0)<\infty$ such that
\[
\big\|\PP(u_N\cdot\nabla u_N)\big\|_{L^2(\T^3)}^{2}\ \ge\ c_0\,N^{3}
\qquad\text{for all }N\ge N_* .
\]
\end{enumerate}
\end{theorem}

\begin{proof}
(1) is Lemma~\ref{app:structure} together with
Theorem~\ref{app:exactlaws}. (2) is Proposition~\ref{app:twosided}.

(3) By Theorem~\ref{app:vnondeg}, $\PP_{\R^3}(V\cdot\nabla V)\not\equiv0$,
where $V$ is the profile of Definition~\ref{app:Vdef} attached to the
given $(\chi,v_0)$. By Lemma~\ref{app:testexists} there is a real
$\Psi\in C^\infty_c(\R^3;\R^3)$, divergence-free, supported in
$\{|y|\le R\}$ for some $R\ge1$, with $I_\Psi\neq0$. Fix such a $\Psi$.
Lemma~\ref{app:innerexact} and Theorem~\ref{app:rate} give the exact
inner rescaling with $\varepsilon_N+\varepsilon'_N=O(N^{-1}\log N)$;
Lemma~\ref{app:psiprops} and Theorem~\ref{app:pairing} convert the
pairing of $\PP(u_N\cdot\nabla u_N)$ against the concentrated field
$\psi_N$ into $(2\pi)^{-3}N^{3/2}(I_\Psi+\mathcal E_N)$ with
$\mathcal E_N\to0$; and Proposition~\ref{app:N3} then yields the stated
bound with
\[
c_0=c_0(\chi,v_0)=\frac{|I_\Psi|^{2}}
{2\,(2\pi)^{3}\,\|\Psi\|_{L^2(\R^3)}^{2}},
\qquad
N_*=N_*(\chi,v_0)
=\min\Big\{N_0\ge\max(8R,32):|\mathcal E_N|\le(1-2^{-1/2})|I_\Psi|
\ \forall N\ge N_0\Big\}.
\]
Both depend on $(\chi,v_0)$ alone, through $V=V^{\chi,v_0}$ and through
the test field $\Psi$ selected for it.
\end{proof}

\begin{remark}[what the constants depend on, and that they are
non-effective]\label{app:noneffective}
The constants $c_0$ and $N_*$ produced above depend on the pair
$(\chi,v_0)$ only, through the profile $V=V^{\chi,v_0}$ and through the
test field $\Psi$ selected in Lemma~\ref{app:testexists}. They are
\textbf{not effective}. The reason is localised precisely at
Lemma~\ref{app:testexists}: $\Psi$ is obtained from the \emph{density}
of $C^\infty_{c,\sigma}(\R^3)$ in $L^2_\sigma(\R^3)$
(Remark~\ref{app:classical-density}), a pure existence argument, so
neither $\Psi$, nor its support radius $R$, nor $I_\Psi$, nor
$\|\Psi\|_{L^2}$, nor $c_0$, nor $N_*$ is exhibited. Every other
constant in this appendix is explicit, namely those of
Lemmas~\ref{app:shell}--\ref{app:lower},
Proposition~\ref{app:twosided}, Theorems~\ref{app:exactlaws},
\ref{app:Vprops}, \ref{app:farfield}, \ref{app:rate},
\ref{app:pairing}, \ref{app:tauformula}, \ref{app:vnondeg},
Lemmas~\ref{app:innerexact}, \ref{app:psiprops}, \ref{app:symbol},
\ref{app:tauwd}, \ref{app:perturb}, and Proposition~\ref{app:N3}
\emph{as a function of} $(\Psi,R,I_\Psi)$.

An effective variant is available in principle --- take $\Psi$ with
$\widehat\Psi$ a smooth bump supported in the explicit set $\mathcal U$
of Theorem~\ref{app:vnondeg}, mollified to compact support in $y$ ---
but it is expensive: non-vanishing is certified only for
$|\zeta|<\pi^2/2048\approx4.8\times10^{-3}$, so such a $\Psi$ has
Fourier support at that scale, hence spatial radius $R\gtrsim10^{3}$,
hence $N_*\ge8R\gtrsim10^{4}$, with a correspondingly small $c_0$. This
is the sole source of non-effectivity in the paper.
\end{remark}

\begin{remark}[the exponent $3$ is sharp, and the sharp cutoff is not
covered]\label{app:sharpness}
The exponent cannot be improved: by Theorem~\ref{main}(d) and
Theorem~\ref{app:exactlaws},
$\|\PP(u_N\cdot\nabla u_N)\|_2^2\le\|u_N\|_\infty^2H_1(u_N)$, and both
$\|u_N\|_\infty$ and $H_1(u_N)$ are $O(N)$ by
Proposition~\ref{app:twosided} and
$\|u_N\|_\infty\le\sum_{k\neq0}|\hat u_N(k)|\le\|v_0\|\Sigma_N^\chi
\le\|v_0\|S_N\le128\|v_0\|N$; so the left-hand side is $O(N^3)$.

Smoothness of $\chi$ is used in exactly one place, namely
Theorem~\ref{app:rate}, through $|\nabla F|\lesssim c_\chi|\xi|^{-3}$.
For the sharply truncated family $\chi=\mathbf 1_{[0,1]}$ this bound
fails, $F$ jumping across $|\xi|=1$; the Riemann-sum defect then carries
a lattice-point (Gauss-circle) term, and $V$ loses the rapidly decaying
tail correction of Theorem~\ref{app:farfield}, acquiring instead an
oscillatory $|y|^{-2}$ tail. Nothing in \S\ref{app:duality} or
\S\ref{app:nondeg} would change, both being insensitive to $\chi$; the
sharp-cutoff case is thus a technical gap in the inner-rescaling step,
not a structural one. It is left open because no statement in this paper
requires it (see Remark~\ref{lstarremark}).
\end{remark}
  inside the paper's \appendix.

\section{Complete proof of the capacity bound (Theorem~\ref{capacity})}
\label{app:capacity}

This appendix proves Theorem~\ref{capacity} in full. Nothing below is
quoted from elsewhere except four classical facts, each isolated and
labelled as such (Remarks~\ref{app:classical-riesz},
\ref{app:classical-homog}, \ref{app:classical-density} and
\ref{app:classical-leray}); no numerical statement is used anywhere.

The chain is: the family $u_N$ has an exact scaling
$\hat u_N(k)=N^{-2}F(k/N)$ (\S\ref{app:profile}), so that the lattice
sums reconstructing $u_N$ near the origin are \emph{literally} Riemann
sums, of spacing $1/N$, for the inverse Fourier transform $V$ of $F$
(\S\ref{app:riemann}); pairing $\PP(u_N\cdot\nabla u_N)$ against a test
field concentrated at scale $1/N$ therefore converts, with an explicit
error, into the fixed continuum pairing
$\int_{\R^3}(V\cdot\nabla V)\cdot\Psi$ times $N^{3/2}$
(\S\ref{app:duality}); and that continuum pairing is nonzero for some
$\Psi$ because $\PP_{\R^3}(V\cdot\nabla V)\not\equiv0$, which is the
Fourier form of the elementary fact that the Oseen field is not a
stationary Euler flow (\S\ref{app:nondeg}). The bookkeeping is assembled
in \S\ref{app:final}.

\subsection{Conventions}
\label{app:conventions}

On $\T^3=(\R/2\pi\Z)^3$ we use the normalised measure
$d\mu=(2\pi)^{-3}dx$, so that for $u:\T^3\to\R^3$
\[
u(x)=\sum_{k\in\Z^3}\hat u(k)e^{ik\cdot x},\qquad
\hat u(k)=\int_{\T^3}u(x)e^{-ik\cdot x}\,d\mu(x),
\]
\[
\langle f,g\rangle=\int_{\T^3}f\cdot\bar g\,d\mu
=\sum_{k\in\Z^3}\hat f(k)\cdot\overline{\hat g(k)},
\qquad
H_r(u)=\sum_{k\in\Z^3}|k|^{2r}|\hat u(k)|^2 .
\]
On $\R^3$ we use the forward transform and the inverse-transform
bracket
\[
\mathcal F f(\zeta)=\int_{\R^3}f(y)e^{-i\zeta\cdot y}\,dy,
\qquad
\widetilde G(y)=\int_{\R^3}G(\xi)e^{i\xi\cdot y}\,d\xi ,
\]
so that $\mathcal F\widetilde G=(2\pi)^3G$ and
$\|\widetilde G\|_{L^2(\R^3)}^2=(2\pi)^3\|G\|_{L^2(\R^3)}^2$.

For $\xi\in\R^3\setminus\{0\}$ put
$P_\xi=I-\dfrac{\xi\otimes\xi}{|\xi|^2}$. Thus $P_\xi$ is the orthogonal
projection onto $\xi^\perp$: it is symmetric, idempotent, $P_\xi\xi=0$,
$\|P_\xi\|\le1$, $P_{-\xi}=P_\xi$, it is homogeneous of degree $0$ and
smooth on $\R^3\setminus\{0\}$, and consequently
$P_{k}=P_{k/N}$ for every $k\neq0$ and every $N>0$. The Leray
projection $\PP$ acts on $\T^3$ by
$\widehat{\PP f}(k)=P_k\hat f(k)$ for $k\neq0$, leaving the mean
untouched, and on $\R^3$ by the same symbol $P_\zeta$. On both
$L^2(\T^3)$ and $L^2(\R^3)$, $\PP$ is an orthogonal projection, hence
self-adjoint.

Throughout, $\|v_0\|$ is the Euclidean norm of $v_0\in\R^3$, and
$e_1,e_2,e_3$ is the standard basis. Repeated Latin indices are summed.
One notational warning: in this appendix $V$ always denotes the
continuum profile of Definition~\ref{app:Vdef}; the unrelated use of
$V$ for the modal variance $\mathrm{Var}_p(|k|^2)$ in
Theorem~\ref{main}(a) does not occur here.

\subsection{The cutoff class, the family, and its exact laws}
\label{app:laws}

We restate the definitions of the main text for completeness.

\begin{definition}[admissible cutoff; restatement of
Definition~\ref{admissible}]\label{app:admissible}
$\chi\in C^4([0,\infty);[0,1])$ with $\chi\equiv1$ on
$[0,\tfrac12]$ and $\operatorname{supp}\chi\subset[0,1]$. No
monotonicity is assumed --- it is used nowhere. Write
$c_\chi:=1+\|\chi'\|_{L^\infty}\ (\ge1)$.
\end{definition}

Such $\chi$ exist in abundance. The standard mollifier construction
gives a $C^\infty$ example, $\chi(t)=\Theta\big(2t-1\big)$ with
$\Theta(s)=1$ for $s\le0$, $\Theta(s)=0$ for $s\ge1$ and
$\Theta(s)=\frac{e^{-1/(1-s)}}{e^{-1/s}+e^{-1/(1-s)}}$ on $(0,1)$; the
exactly rational cutoff of Appendix~\ref{app:certificates}, built from
the degree-$9$ smoothstep, lies in $C^4\setminus C^5$ and is likewise
admissible.

\begin{remark}[what smoothness is actually used]\label{app:c1}
Definition~\ref{app:admissible} is still stronger than the proof
requires. The only place any derivative of $\chi$ is used
load-bearingly is Theorem~\ref{app:rate}, and it enters solely through
$c_\chi=1+\|\chi'\|_\infty$; Sections~\ref{app:duality} and
\ref{app:nondeg} use only $0\le\chi\le1$, $\chi\equiv1$ on
$[0,\tfrac12]$ and $\operatorname{supp}\chi\subset[0,1]$, and
Sections~\ref{app:laws}, \ref{app:lattice}, \ref{app:profile} use no
derivative of $\chi$ at all. Every statement of this appendix on which
Theorem~\ref{app:main} depends therefore holds verbatim for
$\chi\in C^1([0,\infty);[0,1])$ --- indeed for Lipschitz $\chi$ ---
with those support properties. The four derivatives are demanded only
so that Theorem~\ref{app:farfield} holds as stated with $M=4$; that
theorem is consumed only by the inert Remark~\ref{app:noperiod}, and
nothing in the proof of Theorem~\ref{app:main} uses it. In particular
$V\in C^\infty$ (Theorem~\ref{app:Vprops}(3)) and the Paley--Wiener
statement there are bought by the \emph{compact support} of $F$, not by
smoothness of $\chi$, and are unaffected.
\end{remark}

\begin{definition}\label{app:family}
Fix $v_0\in\R^3\setminus\{0\}$, an admissible $\chi$, and an integer
$N\ge2$. Define $u_N=u_N^{\chi,v_0}$ on $\T^3$ by
\begin{equation}\label{app:famdef}
\hat u_N(k)=\chi\Big(\frac{|k|}{N}\Big)\frac{P_kv_0}{|k|^{2}}
\quad(k\neq0),\qquad \hat u_N(0)=0 .
\end{equation}
\end{definition}

Since $\hat u_N(k)=0$ for $|k|>N$, $u_N$ is a finite trigonometric
polynomial banded at $|k|\le N$; its coefficients are real vectors and
no integrality of $v_0$ is required.

\begin{lemma}[basic structure]\label{app:structure}
$u_N$ is real, zero-mean, exactly divergence-free, $C^\infty$, and
$u_N\not\equiv0$.
\end{lemma}

\begin{proof}
Zero mean is $\hat u_N(0)=0$. Reality: the coefficients are real
vectors and $P_{-k}=P_k$, so
$\hat u_N(-k)=\hat u_N(k)=\overline{\hat u_N(k)}$, which is exactly the
reality condition $\hat u(-k)=\overline{\hat u(k)}$. Divergence-free:
$\widehat{\nabla\cdot u_N}(k)=ik\cdot\hat u_N(k)$ and $k\cdot P_kv_0=0$
for every $k\neq0$. Smoothness: a finite sum of exponentials. Finally,
among $e_1,e_2,e_3$ at least two are not parallel to $v_0$; pick such a
$k$ with $|k|=1$. Then $\chi(1/N)=1$ because $1/N\le\tfrac12$, and
$P_kv_0\neq0$, so $\hat u_N(k)\neq0$.
\end{proof}

The exact laws rest on a band-symmetry identity which, because $\chi$ is
not assumed monotone and the truncation is not a sharp ball, must be
available for a \emph{general} radial weight, with no positivity
assumption.

\begin{lemma}[band symmetry, lattice, general radial weight]
\label{app:bandsym}
Let $f:\Z^3\setminus\{0\}\to\R$ be radial, i.e.\ $f(k)=\varphi(|k|^2)$
for some $\varphi$, and suppose $\sum_{k\neq0}|k|^2|f(k)|<\infty$. Then
for all $i,j$
\[
\sum_{k\neq0}k_ik_j\,f(k)
=\delta_{ij}\,\frac13\sum_{k\neq0}|k|^2f(k).
\]
\end{lemma}

\begin{proof}
Absolute summability of $|k|^2f$ dominates every summand
$|k_ik_jf(k)|\le|k|^2|f(k)|$, so all rearrangements below are
legitimate. The index set $\Z^3\setminus\{0\}$ is invariant under each
sign flip $\sigma_i:k_i\mapsto-k_i$ and under all coordinate
permutations, and $f$, depending on $|k|^2$ only, is invariant under
both. For $i\neq j$ the bijection $\sigma_i$ negates the summand
$k_ik_jf(k)$, so that sum equals its own negative and vanishes. For
$i=j$, coordinate permutations carry the three sums $\sum_kk_i^2f(k)$
into one another, so they are equal; their sum is $\sum_k|k|^2f(k)$.
\end{proof}

\begin{lemma}[band symmetry, continuum]\label{app:bandsymcont}
Let $g:\R^3\to\R$ be radial with $\int_{\R^3}|\xi|^2|g(\xi)|\,d\xi<\infty$.
Then $\int\xi_i\xi_jg\,d\xi=\delta_{ij}\frac13\int|\xi|^2g\,d\xi$.
\end{lemma}

\begin{proof}
Identical: Lebesgue measure on $\R^3$ and $g$ are both invariant under
sign flips and permutations of the coordinates, and the hypothesis
dominates the integrand.
\end{proof}

\begin{theorem}[exact family laws]\label{app:exactlaws}
Set
\[
S_N^\chi:=\sum_{k\neq0}\frac{\chi(|k|/N)^2}{|k|^2},\qquad
T_N^\chi:=\sum_{k\neq0}\frac{\chi(|k|/N)^2}{|k|^4},\qquad
\Sigma_N^\chi:=\sum_{k\neq0}\frac{\chi(|k|/N)}{|k|^2},
\]
all three finite sums. Then
\[
H_0(u_N)=\tfrac23\|v_0\|^2\,T_N^\chi,\qquad
H_1(u_N)=\tfrac23\|v_0\|^2\,S_N^\chi,\qquad
N_0^2(u_N)=\frac{S_N^\chi}{T_N^\chi},\qquad
u_N(0)=\tfrac23\,\Sigma_N^\chi\,v_0 ,
\]
and consequently
$\|u_N\|_\infty^2/H_1(u_N)\ge 2(\Sigma_N^\chi)^2/(3S_N^\chi)$.
\end{theorem}

\begin{proof}
Since $P_k$ is an orthogonal projection,
$|P_kv_0|^2=\|v_0\|^2-(k\cdot v_0)^2/|k|^2$. Hence
\[
H_0=\sum_{k\neq0}\frac{\chi(|k|/N)^2}{|k|^4}\,|P_kv_0|^2
=\|v_0\|^2\sum_{k\neq0}\frac{\chi^2}{|k|^4}
-\sum_{k\neq0}\frac{\chi^2\,(k\cdot v_0)^2}{|k|^6}.
\]
Apply Lemma~\ref{app:bandsym} with the radial weight
$f(k)=\chi(|k|/N)^2|k|^{-6}$ (a finite sum, so the hypothesis is
trivially met):
\[
\sum_{k\neq0}\frac{\chi^2(k\cdot v_0)^2}{|k|^6}
=v_{0i}v_{0j}\sum_{k\neq0}k_ik_jf(k)
=\frac{\|v_0\|^2}{3}\sum_{k\neq0}\frac{\chi^2}{|k|^4}
=\frac{\|v_0\|^2}{3}T_N^\chi .
\]
Therefore $H_0=\|v_0\|^2T_N^\chi-\frac13\|v_0\|^2T_N^\chi
=\frac23\|v_0\|^2T_N^\chi$. The same computation with
$f=\chi^2|k|^{-4}$ applied to
$H_1=\sum_{k}|k|^2|\hat u_N(k)|^2$ gives
$H_1=\frac23\|v_0\|^2S_N^\chi$, and $N_0^2=H_1/H_0$ follows. For the
point value,
\[
u_N(0)=\sum_{k\neq0}\hat u_N(k)
=v_0\,\Sigma_N^\chi-\sum_{k\neq0}\frac{\chi\,k\,(k\cdot v_0)}{|k|^4},
\]
and Lemma~\ref{app:bandsym} with $f=\chi|k|^{-4}$ turns the last sum
into $\frac13v_0\Sigma_N^\chi$; hence
$u_N(0)=\frac23\Sigma_N^\chi v_0$. Finally
$\|u_N\|_\infty\ge|u_N(0)|=\frac23\Sigma_N^\chi\|v_0\|$, so
$\|u_N\|_\infty^2/H_1\ge\frac49(\Sigma_N^\chi)^2\|v_0\|^2
\big/\big(\frac23\|v_0\|^2S_N^\chi\big)
=\frac23(\Sigma_N^\chi)^2/S_N^\chi$.
\end{proof}

\subsection{Two-sided lattice bounds for the bandwidth}
\label{app:lattice}

The constants below are crude and explicit; no sharpness is attempted.

\begin{lemma}[dyadic shell count]\label{app:shell}
For every integer $j\ge0$,
$\#\{k\in\Z^3:2^j\le|k|<2^{j+1}\}\le64\cdot8^{j}$.
\end{lemma}

\begin{proof}
Such $k$ satisfy $|k|_\infty\le|k|<2^{j+1}$, hence
$|k_i|\le2^{j+1}-1$ for each $i$, so the count is at most
$(2^{j+2}-1)^3<2^{3j+6}=64\cdot8^{j}$.
\end{proof}

\begin{lemma}[upper bounds]\label{app:upper}
With $S_M:=\sum_{1\le|k|\le M}|k|^{-2}$ and
$T_\infty:=\sum_{k\neq0}|k|^{-4}$ one has
$S_M\le128M$ for every real $M\ge1$, and $T_\infty\le128$.
\end{lemma}

\begin{proof}
Let $J=\lfloor\log_2M\rfloor$. Every $k$ with $1\le|k|\le M$ lies in a
shell $2^j\le|k|<2^{j+1}$ with $0\le j\le J$, where $|k|^{-2}\le4^{-j}$.
By Lemma~\ref{app:shell},
$S_M\le\sum_{j=0}^{J}64\cdot8^{j}4^{-j}=64(2^{J+1}-1)<128\cdot2^{J}\le128M$.
Similarly $T_\infty\le\sum_{j\ge0}64\cdot8^{j}16^{-j}=64\cdot2=128$.
\end{proof}

\begin{lemma}[lower bound]\label{app:lower}
For every real $M\ge16$ one has $S_M\ge M/87$.
\end{lemma}

\begin{proof}
Consider
$\mathcal B:=\{k\in\Z^3:\ \tfrac{M}{2\sqrt3}<k_i\le\tfrac{M}{\sqrt3}\
\text{for }i=1,2,3\}$. For $k\in\mathcal B$ we have
$|k|^2\le3(M/\sqrt3)^2=M^2$ and $|k|^2>3M^2/12=M^2/4$, so every
$k\in\mathcal B$ is counted in $S_M$ (in particular $k\neq0$) and
contributes $|k|^{-2}\ge M^{-2}$. Each coordinate ranges over a
half-open interval of length $M/(2\sqrt3)$, which contains at least
$\frac{M}{2\sqrt3}-1$ integers; for $M\ge16$,
$\frac{M}{2\sqrt3}-1\ge0.28868M-0.0625M\ge0.226M$. Hence
$\#\mathcal B\ge(0.226M)^3>0.01154M^3$ and
$S_M\ge0.01154M^3\cdot M^{-2}=0.01154M\ge M/87$.
\end{proof}

\begin{proposition}[two-sided bandwidth bounds]\label{app:twosided}
For every admissible $\chi$, every $v_0\neq0$ and every integer
$N\ge32$,
\[
\frac{N}{348}\ \le\ S_N^\chi\ \le\ 128N,\qquad
6\ \le\ T_N^\chi\ \le\ 128,\qquad
\frac{N}{44544}\ \le\ N_0^2(u_N)\ \le\ \frac{64}{3}N .
\]
\end{proposition}

\begin{proof}
Upper bounds: $0\le\chi\le1$ and $\chi(|k|/N)=0$ for $|k|>N$, so
$S_N^\chi\le S_N\le128N$ and $T_N^\chi\le T_\infty\le128$ by
Lemma~\ref{app:upper}. Lower bound for $T_N^\chi$: for $N\ge2$ one has
$\chi(1/N)=1$, and the six lattice points with $|k|=1$ each contribute
$1$, so $T_N^\chi\ge6$. Lower bound for $S_N^\chi$: $\chi(|k|/N)=1$
whenever $|k|\le N/2$, hence $S_N^\chi\ge S_{N/2}$; for $N\ge32$ we have
$N/2\ge16$, so Lemma~\ref{app:lower} gives
$S_N^\chi\ge (N/2)/87\ge N/348$. Combining with
$N_0^2=S_N^\chi/T_N^\chi$ (Theorem~\ref{app:exactlaws}) yields
$N_0^2\ge\frac{N/348}{128}=\frac{N}{44544}$ and
$N_0^2\le\frac{128N}{6}=\frac{64}{3}N$.
\end{proof}

\subsection{The profile $V$}
\label{app:profile}

\begin{definition}\label{app:Vdef}
On $\R^3$ put
\[
h(\xi):=\frac{P_\xi v_0}{|\xi|^2}\quad(\xi\neq0),\qquad
F(\xi):=\chi(|\xi|)\,h(\xi)\quad(\xi\neq0),\quad F(0):=0,
\]
\[
V(y):=\widetilde F(y)=\int_{\R^3}F(\xi)e^{i\xi\cdot y}\,d\xi .
\]
\end{definition}

Because $P_\xi$ is homogeneous of degree $0$ we have
$P_{k/N}=P_k$ and therefore
$F(k/N)=\chi(|k|/N)\,N^2P_kv_0/|k|^2$, i.e.\ the exact scaling identity
\begin{equation}\label{app:scaling}
\hat u_N(k)=N^{-2}F(k/N)\qquad\text{for every }k\in\Z^3
\end{equation}
(the case $k=0$ holds by the conventions $\hat u_N(0)=0=F(0)$).
Identity~\eqref{app:scaling} is the engine of the whole argument: it
says that the family is a single fixed continuum profile sampled on the
lattice $N^{-1}\Z^3$.

\begin{theorem}[elementary properties of $V$]\label{app:Vprops}
\begin{enumerate}
\item[(1)] $|F(\xi)|\le\|v_0\|\,|\xi|^{-2}\mathbf 1_{\{|\xi|\le1\}}$;
hence $F\in L^1(\R^3)\cap L^p(\R^3)$ for every $p<3/2$, with
$\|F\|_{L^1}\le4\pi\|v_0\|$.
\item[(2)] $V$ is real, even, bounded and continuous, with
$\|V\|_\infty\le4\pi\|v_0\|$.
\item[(3)] $V\in C^\infty(\R^3)$, and for every multi-index $\beta$
\[
\partial^\beta V(y)=\int_{\R^3}(i\xi)^\beta F(\xi)e^{i\xi\cdot y}\,d\xi,
\qquad
\|\partial^\beta V\|_\infty\le\|v_0\|\int_{|\xi|\le1}|\xi|^{|\beta|-2}d\xi
=\frac{4\pi\|v_0\|}{|\beta|+1} .
\]
In particular each $\|\partial_iV_j\|_\infty\le2\pi\|v_0\|$, so
$\|\nabla V\|_\infty\le6\pi\|v_0\|$ for the Frobenius norm. Moreover $V$
extends to an entire function on $\mathbb C^3$ of exponential type
$\le1$.
\item[(4)] $\nabla\cdot V\equiv0$ on $\R^3$.
\item[(5)] $V(0)=\frac{8\pi}{3}\big(\int_0^\infty\chi(r)\,dr\big)v_0\neq0$.
\end{enumerate}
\end{theorem}

\begin{proof}
(1) $|P_\xi v_0|\le\|v_0\|$, $0\le\chi\le1$ and
$\operatorname{supp}\chi\subset[0,1]$ give the pointwise bound;
$\int_{|\xi|\le1}|\xi|^{-2}d\xi=4\pi\int_0^1dr=4\pi$, and
$\int_{|\xi|\le1}|\xi|^{-2p}d\xi<\infty$ exactly when $2p<3$.

(2) $F\in L^1$ gives, by dominated convergence, continuity of $V$ and
$\|V\|_\infty\le\|F\|_{L^1}\le4\pi\|v_0\|$. $F$ is real-valued and even
($P_{-\xi}=P_\xi$), so
$V(y)=\int F(\xi)\cos(\xi\cdot y)\,d\xi$ is real and even.

(3) By (1), $\xi^\beta F\in L^1$ for every $\beta$, since the exponent
$|\beta|-2>-3$ is locally integrable in $\R^3$ and the support is
bounded. Differentiation under the integral sign is then justified by
dominated convergence, giving the displayed formula and, by (1) again,
the bound
$\|v_0\|\int_{|\xi|\le1}|\xi|^{|\beta|-2}d\xi
=4\pi\|v_0\|\int_0^1r^{|\beta|}dr=4\pi\|v_0\|/(|\beta|+1)$. Since $F$
is an $L^1$ function supported in $\{|\xi|\le1\}$, the integral
$\int F(\xi)e^{i\xi\cdot y}d\xi$ converges for all complex $y$ and is
entire in $y$, with $|V(y)|\le\|F\|_{L^1}e^{|\operatorname{Im}y|}$;
this is exponential type $\le1$.

(4) $\nabla\cdot V(y)=\int i\,\xi\cdot F(\xi)e^{i\xi\cdot y}d\xi$ and
$\xi\cdot F(\xi)=\chi(|\xi|)\,\xi\cdot P_\xi v_0/|\xi|^2=0$ pointwise.

(5) $V(0)=\int F
=v_0\int\frac{\chi(|\xi|)}{|\xi|^2}d\xi
-\int\frac{\chi(|\xi|)\,\xi\,(\xi\cdot v_0)}{|\xi|^4}d\xi$.
Lemma~\ref{app:bandsymcont} with the radial
$g(\xi)=\chi(|\xi|)|\xi|^{-4}$ (integrable against $|\xi|^2$ by (1))
gives $\int\xi_i\xi_jg=\frac{\delta_{ij}}3\int|\xi|^2g
=\frac{\delta_{ij}}3\int\chi(|\xi|)|\xi|^{-2}d\xi$, so the second
integral equals $\frac13v_0\int\chi(|\xi|)|\xi|^{-2}d\xi$. Hence
\[
V(0)=\frac23v_0\int_{\R^3}\frac{\chi(|\xi|)}{|\xi|^2}\,d\xi
=\frac23v_0\cdot4\pi\int_0^\infty\chi(r)\,dr
=\frac{8\pi}{3}\Big(\int_0^\infty\chi\Big)v_0 ,
\]
and $\int_0^\infty\chi\ge\int_0^{1/2}1=\frac12>0$ because $\chi\ge0$ and
$\chi\equiv1$ on $[0,\tfrac12]$.
\end{proof}

\begin{remark}[structural form]\label{app:oseenform}
$F=P_\xi\big(v_0\chi(|\xi|)|\xi|^{-2}\big)$, so
$V=\PP\big(v_0\,g_\chi\big)$ where $g_\chi$ is the inverse transform of
the smoothed Newtonian symbol $\chi(|\cdot|)|\cdot|^{-2}$; solving
$\Delta\phi=g_\chi$ with $\phi$ radial exhibits $V$ in the Oseen normal
form $V(y)=a(r)v_0+b(r)(v_0\cdot\hat y)\hat y$ with
$a=\phi''+\phi'/r$, $b=\phi'/r-\phi''$. This form is not used below; it
explains why the far field of $V$ is exactly the Oseen profile
(Theorem~\ref{app:farfield}).
\end{remark}

We record here the two classical transform facts used later.

\begin{remark}[classical input: Riesz-potential transforms]
\label{app:classical-riesz}
For $0<\alpha<3$, and by analytic continuation for every $\alpha$ at
which the Gamma factors have no pole, the homogeneous function
$|y|^{-\alpha}$ on $\R^3$ is a tempered distribution whose Fourier
transform is the homogeneous distribution
\[
\mathcal F\big[|y|^{-\alpha}\big](\zeta)=c_{3,\alpha}|\zeta|^{\alpha-3},
\qquad
c_{3,\alpha}=\frac{2^{3-\alpha}\pi^{3/2}\,\Gamma\!\big(\frac{3-\alpha}{2}\big)}
{\Gamma\!\big(\frac{\alpha}{2}\big)} .
\]
See, e.g., \cite{SteinWeiss,GelfandShilov,Grafakos} (we give no section
locators, not having verified them). We use
$\Gamma(\tfrac12)=\sqrt\pi$, $\Gamma(-\tfrac12)=-2\sqrt\pi$,
$\Gamma(-\tfrac32)=\tfrac43\sqrt\pi$, which give
\[
c_{3,-1}=-8\pi,\quad c_{3,1}=4\pi,\quad c_{3,2}=2\pi^2,\quad
c_{3,4}=-\pi^2,\quad c_{3,6}=\tfrac{\pi^2}{12},
\]
together with the elementary rules
$\mathcal F[y_ay_bf]=-\partial_a\partial_b\mathcal F[f]$ and
$\mathcal F[y_ay_by_cy_df]=\partial_a\partial_b\partial_c\partial_d\mathcal F[f]$,
valid for tempered $f$.
\end{remark}

\begin{remark}[classical input: products and convolutions of homogeneous
distributions]\label{app:classical-homog}
Let $A,B\in L^1_{\rm loc}(\R^3)$ be homogeneous of degrees $-a,-b$ with
$0<a,b<3$ and $a+b<3$, and smooth away from the origin. Then $AB$ is
locally integrable and homogeneous of degree $-(a+b)$, hence tempered;
$\mathcal FA,\mathcal FB$ are locally integrable, homogeneous of degrees
$a-3,b-3$ and smooth away from the origin; the convolution
$\mathcal FA*\mathcal FB$ converges absolutely at every $\zeta\neq0$
(because $(3-a)+(3-b)>3$) and defines a locally integrable homogeneous
function of degree $3-a-b$; and the exchange formula
\[
\mathcal F[AB]=(2\pi)^{-3}\,\mathcal FA*\mathcal FB
\]
holds as tempered distributions. See, e.g.,
\cite{GelfandShilov,Grafakos}. We use this only in
Theorem~\ref{app:tauformula}, with $a=b=1$.
\end{remark}

\subsection{The far field, and why periodisation is unavailable}
\label{app:farfield-sec}

\begin{theorem}[far-field asymptotics]\label{app:farfield}
Let
\[
W(y):=\frac{v_0}{|y|}+\frac{(v_0\cdot y)\,y}{|y|^3}\qquad(y\neq0).
\]
Then $\widetilde h=\pi^2W$ as tempered distributions, and
\[
V(y)=\pi^2W(y)+\rho(y),\qquad
|\rho(y)|\le C_M(\chi,v_0)|y|^{-M}\quad\text{for }|y|\ge1
\text{ and every }2\le M\le4 .
\]
(Four derivatives are all that Definition~\ref{app:admissible}
supplies, and $M=4$ is all that Remark~\ref{app:noperiod} consumes; if
$\chi\in C^m$ with $m\ge2$ the same proof gives every
$2\le M\le m$.)
\end{theorem}

\begin{proof}
\emph{(a) The identity $\widetilde h=\pi^2W$.} Let
$\mathcal O_{ij}(y)=\frac1{8\pi}\big(\frac{\delta_{ij}}{|y|}
+\frac{y_iy_j}{|y|^3}\big)$ be the Oseen tensor, i.e.\ the fundamental
solution of the steady Stokes system
\cite{LadyzhenskayaBook,Galdi2011}, and
$h_{ij}(\xi)=\frac{\delta_{ij}}{|\xi|^2}-\frac{\xi_i\xi_j}{|\xi|^4}$, so
that $h_i=h_{ij}v_{0j}$. We verify $\mathcal F[\mathcal O_{ij}]=h_{ij}$
from Remark~\ref{app:classical-riesz}. First,
$\mathcal F[|y|^{-1}]=c_{3,1}|\zeta|^{-2}=4\pi|\zeta|^{-2}$, so
$\mathcal F\big[\tfrac{\delta_{ij}}{8\pi|y|}\big]
=\tfrac{\delta_{ij}}{2|\zeta|^{2}}$. Second, differentiating
$\partial_i|y|=y_i/|y|$ once more gives the pointwise identity
\[
\partial_i\partial_j|y|=\frac{\delta_{ij}}{|y|}-\frac{y_iy_j}{|y|^3},
\qquad\text{i.e.}\qquad
\frac{y_iy_j}{|y|^3}=\frac{\delta_{ij}}{|y|}-\partial_i\partial_j|y| .
\]
With $\mathcal F[|y|]=c_{3,-1}|\zeta|^{-4}=-8\pi|\zeta|^{-4}$ we get
$\mathcal F[\partial_i\partial_j|y|]
=(i\zeta_i)(i\zeta_j)\mathcal F[|y|]=8\pi\zeta_i\zeta_j|\zeta|^{-4}$, hence
\[
\mathcal F\Big[\frac{y_iy_j}{|y|^3}\Big]
=\frac{4\pi\delta_{ij}}{|\zeta|^2}-\frac{8\pi\zeta_i\zeta_j}{|\zeta|^4}.
\]
Therefore
$\mathcal F[\mathcal O_{ij}]
=\frac1{8\pi}\big(\frac{4\pi\delta_{ij}}{|\zeta|^2}
+\frac{4\pi\delta_{ij}}{|\zeta|^2}
-\frac{8\pi\zeta_i\zeta_j}{|\zeta|^4}\big)
=\frac{\delta_{ij}}{|\zeta|^2}-\frac{\zeta_i\zeta_j}{|\zeta|^4}=h_{ij}$.
All these are identities between homogeneous tempered distributions;
possible ambiguities are supported at the origin only and are excluded
because both sides are homogeneous of degree $-2>-3$, hence locally
integrable, with no room for a delta term. Fourier inversion
($\mathcal F\widetilde G=(2\pi)^3G$) now gives
$\widetilde h=(2\pi)^3\mathcal O\,v_0=\frac{(2\pi)^3}{8\pi}W=\pi^2W$,
where we used $8\pi\,\mathcal O_{ij}v_{0j}=W_i$.

\emph{(b) The remainder.} Write $G:=(1-\chi(|\cdot|))h$, so that
$F-h=-G$ and $\rho:=-\widetilde G$. Then $G$ vanishes on
$\{|\xi|<\tfrac12\}$ and is $C^4$ on all of $\R^3$; since $h$ is
homogeneous of degree $-2$ and smooth off the origin, and
$1-\chi(|\cdot|)$ is $C^4$ with its first four derivatives bounded, we
get
\[
|\partial^\beta G(\xi)|\le C_\beta(\chi)\|v_0\|\,\langle\xi\rangle^{-2-|\beta|}
\qquad(\xi\in\R^3,\ |\beta|\le4),
\]
which for $2\le|\beta|\le4$ places $\partial^\beta G$ in $L^1(\R^3)$
(the exponent $2+|\beta|>3$). From
the distributional identity
$y^\beta\widetilde G=i^{|\beta|}\,\widetilde{\partial^\beta G}$, the
right-hand side for $2\le|\beta|\le4$ is a bounded continuous function
with supremum at most $\|\partial^\beta G\|_{L^1}$. Choosing
$\beta=(M,0,0),(0,M,0),(0,0,M)$ for a given $2\le M\le4$ and using
$|y|^M\le3^{M/2}\max_i|y_i|^M$ shows that $\widetilde G$ agrees on
$\{|y|\ge1\}$ with a continuous function obeying
$|\widetilde G(y)|\le C_M(\chi,v_0)|y|^{-M}$. Since
$V=\widetilde F=\widetilde h-\widetilde G=\pi^2W+\rho$, the claim
follows.
\end{proof}

Thus the decay of $V$ is \emph{exactly} $|y|^{-1}$, with a
$\chi$-independent leading profile $\pi^2W$ satisfying
\begin{equation}\label{app:Wpos}
v_0\cdot W(y)=\frac{\|v_0\|^2}{|y|}+\frac{(v_0\cdot y)^2}{|y|^3}
\ \ge\ \frac{\|v_0\|^2}{|y|}\ >\ 0 .
\end{equation}

\begin{remark}[the naive periodisation argument is unavailable]
\label{app:noperiod}
It is tempting to reconstruct $u_N$ from $V$ by Poisson summation,
$u_N(x)\stackrel{?}{=}N\sum_{m\in\Z^3}V(N(x+2\pi m))$ with the $m\neq0$
terms uniformly $O(1/N)$. \textbf{The naive absolutely convergent
periodisation argument is unavailable}: the series of translates
diverges, because $V$ decays only like $|y|^{-1}$ and its leading
profile has a sign-definite spherical mean in the $v_0$ direction, so
there is no cancellation to exploit and no absolute convergence to
appeal to. Indeed, by Theorem~\ref{app:farfield} with $M=4$, for
$m\neq0$ and $x\in[-\pi,\pi)^3$ one has $|N(x+2\pi m)|\ge N\pi\ge1$ and
therefore, using \eqref{app:Wpos},
\[
v_0\cdot V\big(N(x+2\pi m)\big)
\ \ge\ \frac{\pi^2\|v_0\|^2}{N|x+2\pi m|}
-\frac{C_4\|v_0\|}{N^4|x+2\pi m|^4}.
\]
The shell $|m|_\infty=n$ contains $O(n^2)$ points with
$|x+2\pi m|\asymp n$, so $\sum_{m\neq0}|x+2\pi m|^{-4}<\infty$ while
$\sum_{m\neq0}|x+2\pi m|^{-1}=\infty$; hence
$\sum_{m\neq0}v_0\cdot V(N(x+2\pi m))=+\infty$ for every $x$ and every
$N$. The terms are moreover eventually positive, so the divergence is
not of a kind that a regular linear summation method would remove: for
$a_m\ge0$ with $\sum a_m=\infty$, monotone convergence gives
$\lim_{\varepsilon\downarrow0}\sum_me^{-\varepsilon|m|}a_m=+\infty$, so
Abel summation in particular does not converge here. We make no
assertion beyond this, and none is needed: the point is only that the
naive absolutely convergent periodisation argument is not available,
because $V$ decays only like $|y|^{-1}$ with a non-oscillating leading
profile whose spherical mean, $\tfrac43v_0/r$, is sign-definite in the
$v_0$ direction. If a periodisation formula
is wanted purely for orientation, the correct \emph{renormalised} one,
\[
u_N(x)=N\Big[V(Nx)+\sum_{m\neq0}\big(V(N(x+2\pi m))-V(2\pi Nm)\big)\Big]
+\text{const},
\]
does converge absolutely, the bracketed differences being $O(|m|^{-2})$
by Theorem~\ref{app:farfield}. \textbf{Nothing in this appendix uses any
periodisation identity}: \S\ref{app:riemann} replaces it by an exact
lattice-Riemann-sum identity, which is elementary, quantitative, and
additionally delivers the derivative statement. This remark, and with it
Theorem~\ref{app:farfield}, is inert for the proof of
Theorem~\ref{capacity}.
\end{remark}

\subsection{The inner rescaling: an exact lattice-Riemann identity with
an explicit rate}
\label{app:riemann}

For $\Phi:\R^3\to\mathbb C^d$ put
\[
\mathcal R_N[\Phi]:=N^{-3}\sum_{k\in\Z^3}\Phi(k/N)
-\int_{\R^3}\Phi(\xi)\,d\xi ,
\]
whenever both terms converge absolutely.

\begin{lemma}[exact inner rescaling]\label{app:innerexact}
For every $y\in\R^3$ and every integer $N\ge2$,
\[
u_N(y/N)=N\big(V(y)+E_N(y)\big),\qquad
\partial_iu_{N,j}(y/N)=N^2\big(\partial_iV_j(y)+E'_{N,ij}(y)\big),
\]
where, with $\Phi_y(\xi):=F(\xi)e^{i\xi\cdot y}$ and
$\Phi'_{y,ij}(\xi):=i\xi_iF_j(\xi)e^{i\xi\cdot y}$ (both set to $0$ at
$\xi=0$),
\[
E_N(y)=\mathcal R_N[\Phi_y],\qquad
E'_{N,ij}(y)=\mathcal R_N[\Phi'_{y,ij}] .
\]
These are identities, not approximations.
\end{lemma}

\begin{proof}
Both series are finite (support $|k|\le N$) and both integrals converge
absolutely by Theorem~\ref{app:Vprops}(1),(3). By \eqref{app:scaling},
\[
u_N(y/N)=\sum_{k\in\Z^3}\hat u_N(k)e^{i(k/N)\cdot y}
=\sum_{k}N^{-2}F(k/N)e^{i(k/N)\cdot y}
=N\cdot N^{-3}\sum_k\Phi_y(k/N),
\]
which equals $N\big(\int\Phi_y+\mathcal R_N[\Phi_y]\big)
=N\big(V(y)+E_N(y)\big)$, using $V(y)=\int\Phi_y$
(Definition~\ref{app:Vdef}) and $\Phi_y(0)=F(0)=0$. For the gradient,
$\partial_iu_{N,j}(x)=\sum_kik_i\hat u_{N,j}(k)e^{ik\cdot x}$; at
$x=y/N$ this is
\[
\sum_kiN^{-2}k_iF_j(k/N)e^{i(k/N)\cdot y}
=N^{-1}\sum_ki(k_i/N)F_j(k/N)e^{i(k/N)\cdot y}
=N^{2}\cdot N^{-3}\sum_k\Phi'_{y,ij}(k/N),
\]
and $\int\Phi'_{y,ij}=\partial_iV_j(y)$ by
Theorem~\ref{app:Vprops}(3).
\end{proof}

Everything now rests on bounding a Riemann-sum defect for an explicit,
compactly supported integrand with a single integrable algebraic
singularity.

\begin{theorem}[$C^1_{\rm loc}$ convergence with explicit rate]
\label{app:rate}
There is an absolute constant $A$ (the proof gives $A=3000$) such that
for every real $R\ge1$ and every integer $N\ge8R$,
\[
\varepsilon_N:=\sup_{|y|\le R}|E_N(y)|
\ \le\ A\,\|v_0\|\,\frac{c_\chi(1+\log_2N)+R}{N},
\qquad
\varepsilon'_N:=\sup_{|y|\le R}|E'_N(y)|
\ \le\ A\,\|v_0\|\,\frac{c_\chi+R}{N} .
\]
In particular $\varepsilon_N+\varepsilon'_N=O(N^{-1}\log N)\to0$ as
$N\to\infty$, for each fixed $R$.
\end{theorem}

\begin{proof}
Partition $\R^3$ into the half-open cubes
$Q_k=\frac kN+[0,\frac1N)^3$, $k\in\Z^3$, of side $1/N$, volume
$N^{-3}$ and diameter $\sqrt3/N$. For any $\Phi$ for which both the sum
and the integral converge absolutely,
\begin{equation}\label{app:cellwise}
\mathcal R_N[\Phi]=\sum_{k\in\Z^3}\int_{Q_k}
\big[\Phi(k/N)-\Phi(\xi)\big]\,d\xi ,
\end{equation}
since $|Q_k|=N^{-3}$. Absolute convergence holds for
$\Phi=\Phi_y$ and $\Phi=\Phi'_{y,ij}$: both are supported in
$\{|\xi|\le1\}$ and dominated there by $\|v_0\||\xi|^{-2}$ respectively
$\|v_0\||\xi|^{-1}$, which are integrable on $\R^3$; and the lattice
sums are finite sums, bounded by
$N^{-3}\sum_{0<|k|\le N}\|v_0\|N^{2}|k|^{-2}
=\|v_0\|N^{-1}S_N\le128\|v_0\|$ and by
$N^{-3}\sum_{0<|k|\le N}\|v_0\|N|k|^{-1}
\le\|v_0\|N^{-2}\cdot N\,S_N\le128\|v_0\|$ respectively, using
$|k|^{-1}\le N|k|^{-2}$ for $|k|\le N$ and Lemma~\ref{app:upper}.

\emph{Pointwise bounds.} On $0<|\xi|\le1$ we have
$|F(\xi)|\le\|v_0\||\xi|^{-2}$, and by the product rule applied to
$F=\chi(|\xi|)|\xi|^{-2}P_\xi v_0$,
\[
|\nabla F(\xi)|\ \le\
\underbrace{c_\chi\|v_0\||\xi|^{-2}}_{\nabla\chi(|\xi|)}
+\underbrace{2\|v_0\||\xi|^{-3}}_{\chi\,\nabla(|\xi|^{-2})}
+\underbrace{4\|v_0\||\xi|^{-3}}_{\chi|\xi|^{-2}\nabla(P_\xi v_0)}
\ \le\ 7c_\chi\|v_0\|\,|\xi|^{-3},
\]
where we used $c_\chi\ge1$, $|\xi|\le1$, $|\nabla(|\xi|^{-2})|
=2|\xi|^{-3}$, and, from
$\partial_a\big(\xi_i\xi_j/|\xi|^2\big)
=(\delta_{ai}\xi_j+\delta_{aj}\xi_i)/|\xi|^2
-2\xi_i\xi_j\xi_a/|\xi|^4$, the honest bound
$|\nabla(P_\xi v_0)|\le4\|v_0\|/|\xi|$. Consequently, for $|y|\le R$,
\begin{equation}\label{app:phibounds}
|\Phi_y|\le\frac{\|v_0\|}{|\xi|^{2}},\quad
|\nabla\Phi_y|\le\frac{7c_\chi\|v_0\|}{|\xi|^{3}}+\frac{R\|v_0\|}{|\xi|^{2}},
\qquad
|\Phi'_y|\le\frac{\|v_0\|}{|\xi|},\quad
|\nabla\Phi'_y|\le\frac{8c_\chi\|v_0\|}{|\xi|^{2}}+\frac{R\|v_0\|}{|\xi|} .
\end{equation}
(For the last one, $\nabla(i\xi_iF_je^{i\xi y})$ has the three terms
$|F|\le\|v_0\||\xi|^{-2}$, $|\xi||\nabla F|\le7c_\chi\|v_0\||\xi|^{-2}$
and $R|\xi||F|\le R\|v_0\||\xi|^{-1}$.)

\emph{The core.} Let
$\mathcal K_0:=\{k:\ Q_k\cap\{|\xi|\le4\sqrt3/N\}\neq\emptyset\}$. Any
such $Q_k$ lies in $\{|\xi|\le5\sqrt3/N\}$ and has $|k|\le5\sqrt3<9$,
i.e.\ $|k|\le8$. Note that $k=0$ is in the core, and that although
$\Phi_y(0)=0$ by convention while $\int_{Q_0}\Phi_y\neq0$, both are
bounded by the two terms below, so no separate treatment is needed. The
core contributes to \eqref{app:cellwise} at most
\[
\int_{|\xi|\le5\sqrt3/N}|\Phi_y|\,d\xi
+N^{-3}\!\!\sum_{0<|k|\le8}\!|\Phi_y(k/N)|
\ \le\ \|v_0\|\Big(\frac{4\pi\cdot5\sqrt3}{N}
+\frac1N\!\!\sum_{0<|k|\le8}\!\frac1{|k|^{2}}\Big)
\ \le\ \frac{1133\,\|v_0\|}{N},
\]
using $\int_{|\xi|\le a}|\xi|^{-2}d\xi=4\pi a$, then
$|\Phi_y(k/N)|\le\|v_0\|N^2|k|^{-2}$, and finally
$\sum_{0<|k|\le8}|k|^{-2}\le S_8\le1024$ by Lemma~\ref{app:upper}
($4\pi\cdot5\sqrt3\le109$, and $109+1024=1133$).

\emph{Dyadic annuli.} Put $J:=\lfloor\log_2\big(N/(4\sqrt3)\big)\rfloor$,
which is $\ge0$ because $N\ge8R\ge8>4\sqrt3$. Let $k\notin\mathcal K_0$
and set $r_k:=\inf_{\xi\in \overline{Q_k}}|\xi|$, so $r_k>4\sqrt3/N>0$.
If $r_k>1$ then $\Phi_y$ vanishes identically on $\overline{Q_k}$ and
that cell contributes $0$; otherwise choose the unique integer $j(k)$
with $2^{-j(k)-1}<r_k\le2^{-j(k)}$. Then $j(k)\ge0$, and
$2^{-j(k)}\ge r_k>4\sqrt3/N$ forces $2^{j(k)}<N/(4\sqrt3)$, i.e.\
$0\le j(k)\le J$; so the annuli
$A_j:=\{2^{-j-1}<|\xi|\le2^{-j}\}$, $0\le j\le J$, together with the
core, do cover $\operatorname{supp}\Phi_y$.

Fix $j$ and consider the cells with $j(k)=j$. They are pairwise
disjoint. Each satisfies $\overline{Q_k}\subset\{|\xi|\ge2^{-j-1}\}$ and
$\overline{Q_k}\subset\{|\xi|\le r_k+\sqrt3/N\le\tfrac54 2^{-j}\}$,
because $\sqrt3/N\le\tfrac14 2^{-j}$ (which follows from
$2^{-j}\ge2^{-J}\ge4\sqrt3/N$). Hence all of them lie in the fixed
annulus $B_j:=\{2^{-j-1}\le|\xi|\le\tfrac54 2^{-j}\}$, whose volume is
$\frac{4\pi}3\big((\tfrac54)^3-(\tfrac12)^3\big)2^{-3j}\le9\cdot2^{-3j}$;
therefore $\sum_{j(k)=j}|Q_k|\le9\cdot2^{-3j}$.

Since $\overline{Q_k}$ is convex, avoids the origin, and $\Phi_y$ is
$C^1$ there, the mean value inequality gives
$\big|\int_{Q_k}[\Phi_y(k/N)-\Phi_y(\xi)]d\xi\big|
\le\frac{\sqrt3}{N}\sup_{\overline{Q_k}}|\nabla\Phi_y|\cdot|Q_k|$. By
\eqref{app:phibounds} and $|\xi|\ge2^{-j-1}$ on $\overline{Q_k}$,
\[
\sup_{\overline{Q_k}}|\nabla\Phi_y|
\le 7c_\chi\|v_0\|2^{3(j+1)}+R\|v_0\|2^{2(j+1)}
=56c_\chi\|v_0\|8^{j}+4R\|v_0\|4^{j}.
\]
Summing over the cells with $j(k)=j$,
\[
\sum_{j(k)=j}\Big|\int_{Q_k}[\Phi_y(k/N)-\Phi_y]\Big|
\le\frac{\sqrt3}{N}\big(56c_\chi\|v_0\|8^{j}+4R\|v_0\|4^{j}\big)
\cdot9\cdot2^{-3j}
\le\frac{\|v_0\|}{N}\big(873\,c_\chi+63\,R\,2^{-j}\big).
\]
Summing over $0\le j\le J\le\log_2N$ and adding the core bound,
\[
|E_N(y)|\le\frac{\|v_0\|}{N}
\Big(1133+873\,c_\chi(1+\log_2N)+126R\Big)
\le\frac{3000\,\|v_0\|}{N}\big(c_\chi(1+\log_2N)+R\big),
\]
using $c_\chi\ge1$ and $1+\log_2N\ge1$. This is the first display.

\emph{The gradient.} Repeat verbatim with $\Phi'_{y,ij}$. The core
contributes at most
\[
\int_{|\xi|\le5\sqrt3/N}\frac{\|v_0\|}{|\xi|}\,d\xi
+N^{-3}\!\!\sum_{0<|k|\le8}\!\frac{\|v_0\|N}{|k|}
=\|v_0\|\Big(\frac{2\pi\cdot75}{N^{2}}
+\frac1{N^{2}}\!\!\sum_{0<|k|\le8}\!\frac1{|k|}\Big)
\le\frac{1816\,\|v_0\|}{N^{2}}\le\frac{1816\,\|v_0\|}{N},
\]
where $\int_{|\xi|\le a}|\xi|^{-1}d\xi=2\pi a^2$ and, by
Lemma~\ref{app:shell},
$\sum_{0<|k|\le8}|k|^{-1}\le\sum_{j=0}^{2}64\cdot8^{j}2^{-j}=1344$
(and $2\pi\cdot75\le472$, $472+1344=1816$). On the annuli,
$\sup_{\overline{Q_k}}|\nabla\Phi'_y|
\le8c_\chi\|v_0\|2^{2(j+1)}+R\|v_0\|2^{j+1}
=32c_\chi\|v_0\|4^{j}+2R\|v_0\|2^{j}$, so the level-$j$ contribution is
at most
\[
\frac{\sqrt3}{N}\big(32c_\chi\|v_0\|4^{j}+2R\|v_0\|2^{j}\big)
\cdot9\cdot2^{-3j}
\le\frac{\|v_0\|}{N}\big(499\,c_\chi2^{-j}+32\,R\,4^{-j}\big),
\]
and now the geometric series converge \emph{without} a logarithm,
summing over all $j\ge0$ to at most
$\frac{\|v_0\|}{N}(998c_\chi+43R)$. Adding the core,
$|E'_N(y)|\le\frac{\|v_0\|}{N}(1816+998c_\chi+43R)
\le\frac{3000\|v_0\|}{N}(c_\chi+R)$. The absence of a logarithm is
structural: the singularity of $\Phi'_y$ is only $|\xi|^{-1}$, the extra
factor $\xi$ taming the origin. The derivative estimate is thus not a
consequence of the $C^0$ one, but it is the easier of the two.
\end{proof}

\subsection{The concentrated test field and the duality lower bound}
\label{app:duality}

\begin{definition}\label{app:testfield}
Let $\Psi\in C_c^\infty(\R^3;\R^3)$ be \emph{real}, with
$\nabla\cdot\Psi=0$, $\operatorname{supp}\Psi\subset\{|y|\le R\}$ for
some $R\ge1$, and $\Psi\not\equiv0$. For $N>R/\pi$ define on $\T^3$
\[
\psi_N(x):=N^{3/2}\sum_{m\in\Z^3}\Psi\big(N(x+2\pi m)\big).
\]
\end{definition}

\begin{lemma}[properties of $\psi_N$]\label{app:psiprops}
Let $N>R/\pi$. Then $\psi_N$ is a well-defined real $C^\infty$ vector
field on $\T^3$; it is divergence-free and has zero mean; on the
fundamental domain $[-\pi,\pi)^3$ exactly the term $m=0$ is nonzero, so
$\operatorname{supp}\psi_N\subset\{|x|\le R/N\}$ there; and
\[
\|\psi_N\|_{L^2(\T^3)}=(2\pi)^{-3/2}\|\Psi\|_{L^2(\R^3)},
\]
which is independent of $N$.
\end{lemma}

\begin{proof}
Local finiteness: $\Psi(N(x+2\pi m))\neq0$ forces
$|x+2\pi m|\le R/N<\pi$; but for $x\in[-\pi,\pi)^3$ and $m\neq0$,
$|x+2\pi m|\ge2\pi|m|_\infty-\pi\ge\pi$, a contradiction. Hence on a
neighbourhood of any point the sum has at most one nonzero term, so
$\psi_N$ is smooth, real, and $2\pi\Z^3$-periodic by construction.
Divergence-free: locally
$\nabla\cdot\psi_N(x)=N^{5/2}(\nabla\cdot\Psi)(N(x+2\pi m))=0$.
Zero mean: with $d\mu=(2\pi)^{-3}dx$ and the substitution $y=Nx$,
\[
\int_{\T^3}\psi_N\,d\mu
=(2\pi)^{-3}N^{3/2}\int_{\R^3}\Psi(Nx)\,dx
=(2\pi)^{-3}N^{-3/2}\int_{\R^3}\Psi ,
\]
and $\int_{\R^3}\Psi_i=\int_{\R^3}\nabla\cdot(y_i\Psi)=0$ since
$\nabla\cdot\Psi=0$ and $y_i\Psi\in C^\infty_c$. (No curl
representation of $\Psi$ is needed: zero mean is automatic for a
compactly supported divergence-free field.) Norm:
\[
\|\psi_N\|_{L^2(\T^3)}^2
=(2\pi)^{-3}\int_{|x|\le R/N}N^{3}|\Psi(Nx)|^2dx
=(2\pi)^{-3}N^{3}N^{-3}\|\Psi\|_{L^2(\R^3)}^2 ,
\]
which is the stated identity.
\end{proof}

\begin{theorem}[duality identity with controlled error]\label{app:pairing}
Let $\Psi$ be as in Definition~\ref{app:testfield} and put
\[
I_\Psi:=\int_{\R^3}\big(V\cdot\nabla V\big)\cdot\Psi\,dy .
\]
Then for every integer $N\ge\max(8R,32)$,
\[
\big\langle\PP(u_N\cdot\nabla u_N),\psi_N\big\rangle_{L^2(\T^3)}
=(2\pi)^{-3}N^{3/2}\big(I_\Psi+\mathcal E_N\big),
\]
where
\[
|\mathcal E_N|\ \le\ \|\Psi\|_{L^1(\R^3)}
\Big[\varepsilon_N\big(6\pi\|v_0\|+\varepsilon'_N\big)
+4\pi\|v_0\|\,\varepsilon'_N\Big]
\ =\ O\Big(\frac{\log N}{N}\Big)
\]
with $\varepsilon_N,\varepsilon'_N$ as in Theorem~\ref{app:rate}.
\end{theorem}

\begin{proof}
By Lemma~\ref{app:psiprops}, $\psi_N$ is zero-mean and divergence-free,
so $\hat\psi_N(0)=0$ and $k\cdot\hat\psi_N(k)=0$ for all $k$, whence
$\PP\psi_N=\psi_N$. Since $\PP$ is an orthogonal projection on
$L^2(\T^3)$, hence self-adjoint,
\[
\big\langle\PP(u_N\cdot\nabla u_N),\psi_N\big\rangle
=\big\langle u_N\cdot\nabla u_N,\PP\psi_N\big\rangle
=\big\langle u_N\cdot\nabla u_N,\psi_N\big\rangle .
\]
Both fields are real. Because $N\ge8R$ and $R\ge1$ we have
$R/N\le\tfrac18<\pi$, so Lemma~\ref{app:psiprops} applies and
$\psi_N(x)=N^{3/2}\Psi(Nx)$ on the fundamental domain, supported in
$\{|x|\le R/N\}$. Using $d\mu=(2\pi)^{-3}dx$ and substituting $x=y/N$
($dx=N^{-3}dy$),
\[
\big\langle u_N\cdot\nabla u_N,\psi_N\big\rangle
=(2\pi)^{-3}N^{3/2}N^{-3}\int_{|y|\le R}
(u_N\cdot\nabla u_N)(y/N)\cdot\Psi(y)\,dy .
\]
By Lemma~\ref{app:innerexact}, for $|y|\le R$,
\[
(u_N\cdot\nabla u_N)_j(y/N)
=u_{N,i}(y/N)\,\partial_iu_{N,j}(y/N)
=N^{3}\big(V_i(y)+E_{N,i}(y)\big)\big(\partial_iV_j(y)+E'_{N,ij}(y)\big),
\]
so the factor $N^{-3}N^{3}$ cancels and
\[
\big\langle u_N\cdot\nabla u_N,\psi_N\big\rangle
=(2\pi)^{-3}N^{3/2}\big(I_\Psi+\mathcal E_N\big),
\qquad
\mathcal E_N=\int_{|y|\le R}
\big[E_{N,i}\partial_iV_j+V_iE'_{N,ij}+E_{N,i}E'_{N,ij}\big]\Psi_j\,dy ,
\]
where we used $\operatorname{supp}\Psi\subset\{|y|\le R\}$ to write the
main term as $I_\Psi$. Bounding each factor by its supremum and using
Theorem~\ref{app:Vprops}(2),(3) --- $\|V\|_\infty\le4\pi\|v_0\|$,
$\|\nabla V\|_\infty\le6\pi\|v_0\|$ --- gives the stated estimate;
$\varepsilon_N,\varepsilon'_N=O(N^{-1}\log N)$ by
Theorem~\ref{app:rate}, $R$ and $\Psi$ being fixed.
\end{proof}

\begin{proposition}[the $N^3$ lower bound, given a good test field]
\label{app:N3}
Suppose some $\Psi$ as in Definition~\ref{app:testfield} has
$I_\Psi\neq0$. Then
\[
\big\|\PP(u_N\cdot\nabla u_N)\big\|_{L^2(\T^3)}^2
\ \ge\ \frac{|I_\Psi+\mathcal E_N|^2}
{(2\pi)^{3}\,\|\Psi\|_{L^2(\R^3)}^{2}}\;N^{3}
\qquad(N\ge\max(8R,32)),
\]
and consequently
$\|\PP(u_N\cdot\nabla u_N)\|_{2}^2\ge c_0N^3$ for all $N\ge N_*$, where
\[
c_0=\frac{|I_\Psi|^2}{2\,(2\pi)^{3}\,\|\Psi\|_{L^2(\R^3)}^{2}},
\qquad
N_*=\min\Big\{N_0\ge\max(8R,32):\ |\mathcal E_N|\le(1-2^{-1/2})|I_\Psi|
\ \ \forall N\ge N_0\Big\},
\]
which is finite because $\mathcal E_N\to0$.
\end{proposition}

\begin{proof}
By Cauchy--Schwarz in $L^2(\T^3)$,
\[
\big\|\PP(u_N\cdot\nabla u_N)\big\|_{2}
\ \ge\ \frac{\big|\langle\PP(u_N\cdot\nabla u_N),\psi_N\rangle\big|}
{\|\psi_N\|_{2}} .
\]
Insert Theorem~\ref{app:pairing} and Lemma~\ref{app:psiprops}: the
right-hand side equals
\[
\frac{(2\pi)^{-3}N^{3/2}|I_\Psi+\mathcal E_N|}
{(2\pi)^{-3/2}\|\Psi\|_{L^2(\R^3)}}
=(2\pi)^{-3/2}N^{3/2}\frac{|I_\Psi+\mathcal E_N|}{\|\Psi\|_{L^2(\R^3)}} .
\]
Squaring gives the display. For $N\ge N_*$ we have
$|I_\Psi+\mathcal E_N|\ge2^{-1/2}|I_\Psi|$, whence
$|I_\Psi+\mathcal E_N|^2\ge\tfrac12|I_\Psi|^2$ and the constant $c_0$ as
stated.
\end{proof}

Everything is now reduced to the single question of whether some
admissible $\Psi$ has $I_\Psi\neq0$.

\subsection{Non-degeneracy of the continuum nonlinearity}
\label{app:nondeg}

\subsubsection*{Reduction to a symbol statement}

\begin{lemma}[symbol of the nonlinearity]\label{app:symbol}
$V\cdot\nabla V=\nabla\cdot(V\otimes V)$ has the representation
\[
(V\cdot\nabla V)(y)=\int_{\R^3}M(\zeta)e^{i\zeta\cdot y}\,d\zeta,
\qquad
M_i(\zeta)=i\,\zeta_j\,T_{ij}(\zeta),\qquad T_{ij}:=F_i*F_j ,
\]
where $T$ is real, symmetric, supported in $\{|\zeta|\le2\}$, and
$T,M\in L^1(\R^3)\cap L^2(\R^3)$. Consequently
$\PP(V\cdot\nabla V)\in L^2(\R^3)$ with symbol $P_\zeta M(\zeta)$, and
\[
\PP(V\cdot\nabla V)\equiv0
\quad\Longleftrightarrow\quad
\zeta\times\big(T(\zeta)\zeta\big)=0\ \text{ for a.e. }\zeta .
\]
\end{lemma}

\begin{proof}
$V=\widetilde F$ with $F\in L^1$, and $\partial_jV_i=\widetilde{i\xi_jF_i}$
with $\xi F\in L^1$ (Theorem~\ref{app:Vprops}(1),(3)). The pointwise
product of inverse transforms of $L^1$ functions is the inverse
transform of the convolution, so
$V_j\partial_jV_i=\widetilde{F_j*(i\,\cdot_j F_i)}$, i.e.\ the symbol of
$V\cdot\nabla V$ is
$i\int F_j(\zeta-\eta)\,\eta_j\,F_i(\eta)\,d\eta$. Since
$(\zeta-\eta)\cdot F(\zeta-\eta)=0$ identically (Definition~\ref{app:Vdef}),
we may replace $\eta_jF_j(\zeta-\eta)$ by $\zeta_jF_j(\zeta-\eta)$
inside the integral, obtaining
$M_i(\zeta)=i\zeta_j(F_i*F_j)(\zeta)=i\zeta_jT_{ij}(\zeta)$.

$T$ is real because $F$ is real, symmetric because convolution is
commutative, and supported in $\{|\zeta|\le2\}$ because
$\operatorname{supp}F\subset\{|\xi|\le1\}$. Integrability: by
Theorem~\ref{app:Vprops}(1), $F\in L^{4/3}$; Young's inequality with
$\frac1r=\frac2{4/3}-1=\frac12$ gives $T=F*F\in L^{2}$, and since $T$
has compact support, $T\in L^1\cap L^2$; the same holds for
$M=i\zeta T$, as $|\zeta|\le2$ on the support. By Plancherel,
$V\cdot\nabla V\in L^2(\R^3)$ and, since $|P_\zeta M|\le|M|$,
$\PP(V\cdot\nabla V)\in L^2(\R^3)$ with symbol $P_\zeta M$. Finally
$\PP(V\cdot\nabla V)=0$ in $L^2$ iff $P_\zeta M(\zeta)=0$ for a.e.
$\zeta$, and for a vector $w$ one has $P_\zeta w=0$ iff $w\parallel\zeta$
iff $\zeta\times w=0$; here $w=i\,T(\zeta)\zeta$.
\end{proof}

\begin{remark}[classical input: density of smooth compactly supported
solenoidal fields]\label{app:classical-density}
The space $C^\infty_{c,\sigma}(\R^3;\R^3)$ of real, smooth, compactly
supported, divergence-free vector fields is dense in
$L^2_\sigma(\R^3;\R^3)$, the closed subspace of real $L^2$ fields that
are distributionally divergence-free; indeed $L^2_\sigma$ is by the
standard definition the $L^2$-closure of $C^\infty_{c,\sigma}$, and on
$\R^3$ the two descriptions agree. See, e.g.,
\cite{Temam,LadyzhenskayaBook,ConstantinFoias}.
\end{remark}

\begin{remark}[classical input: Helmholtz--Leray decomposition on $\R^3$]
\label{app:classical-leray}
For $f\in L^2(\R^3;\R^3)$ one has $f-\PP f=\nabla p$ for some
$p\in L^2_{\rm loc}(\R^3)$ with $\nabla p\in L^2$; equivalently, the
symbol of $f-\PP f$ is $\zeta(\zeta\cdot\hat f)/|\zeta|^2$. Same
references.
\end{remark}

\begin{lemma}[from non-degeneracy to a test field]\label{app:testexists}
If $\PP(V\cdot\nabla V)\not\equiv0$, then there exists a real
$\Psi\in C^\infty_c(\R^3;\R^3)$ with $\nabla\cdot\Psi=0$ and
$I_\Psi=\int_{\R^3}(V\cdot\nabla V)\cdot\Psi\neq0$; after enlarging $R$
if necessary, $\Psi$ satisfies Definition~\ref{app:testfield}.
\end{lemma}

\begin{proof}
By Lemma~\ref{app:symbol}, $\PP(V\cdot\nabla V)$ lies in $L^2(\R^3)$,
is real (its symbol $P_\zeta M$ satisfies the reality condition because
$F$ is real and even), is distributionally divergence-free (its symbol
is orthogonal to $\zeta$), and is nonzero by hypothesis. Hence it is a
nonzero element of the real Hilbert space $L^2_\sigma(\R^3)$. By
Remark~\ref{app:classical-density} there is a real
$\Psi\in C^\infty_{c,\sigma}(\R^3)$ with
$\langle\PP(V\cdot\nabla V),\Psi\rangle_{L^2(\R^3)}\neq0$. By
Remark~\ref{app:classical-leray},
$V\cdot\nabla V-\PP(V\cdot\nabla V)=\nabla p$ with $\nabla p\in L^2$, and
$\int\nabla p\cdot\Psi=-\int p\,\nabla\cdot\Psi=0$ because $\Psi$ is
compactly supported and divergence-free. Therefore
$I_\Psi=\langle\PP(V\cdot\nabla V),\Psi\rangle\neq0$. Choosing
$R\ge\max(1,\sup\{|y|:y\in\operatorname{supp}\Psi\})$ puts $\Psi$ in
Definition~\ref{app:testfield}.
\end{proof}

\subsubsection*{The $\chi$-independent leading singularity}

By rotational equivariance we may and do assume from now on that
$v_0=\|v_0\|e_3$: $F$ and $V$ are linear in $v_0$ and rotate with it,
$T$ and $\tau$ are quadratic, and the conclusion
$\PP(V\cdot\nabla V)\not\equiv0$ is rotation invariant.

\begin{definition}\label{app:tau}
$\tau_{ij}:=h_i*h_j$, with $h(\xi)=P_\xi v_0/|\xi|^2$ the uncut symbol
of Definition~\ref{app:Vdef}.
\end{definition}

\begin{lemma}[$\tau$ is well defined and homogeneous]\label{app:tauwd}
The integral $\tau_{ij}(\zeta)=\int_{\R^3}h_i(\eta)h_j(\zeta-\eta)\,d\eta$
converges absolutely for every $\zeta\neq0$, and $\tau$ is homogeneous of
degree $-1$ and smooth on $\R^3\setminus\{0\}$.
\end{lemma}

\begin{proof}
$|h(\xi)|\le\|v_0\||\xi|^{-2}$. Near $\eta=0$ the integrand is
$O(|\eta|^{-2})$ and near $\eta=\zeta$ it is $O(|\zeta-\eta|^{-2})$, both
locally integrable on $\R^3$; for $|\eta|\ge2|\zeta|$ it is
$O(|\eta|^{-4})$, integrable at infinity. Substituting
$\eta=\lambda\eta'$ and using $h(\lambda\eta)=\lambda^{-2}h(\eta)$ gives
$\tau(\lambda\zeta)=\lambda^{-1}\tau(\zeta)$ for $\lambda>0$. Smoothness
away from $0$ follows from differentiating under the integral sign on
the region where the two singularities are separated.
\end{proof}

\begin{theorem}[closed form of the leading obstruction]
\label{app:tauformula}
With $v_0=\|v_0\|e_3$, for every $\zeta\neq0$,
\[
\zeta\times\big(\tau(\zeta)\zeta\big)
=\frac{3\pi^3}{8}\,\|v_0\|^2\,\frac{\zeta_3}{|\zeta|}\,
\big(\zeta\times e_3\big).
\]
In particular $\zeta\times(\tau(\zeta)\zeta)\neq0$ whenever
$\zeta_3\neq0$ and $\zeta\not\parallel e_3$.
\end{theorem}

\begin{proof}
Both sides are quadratic in $\|v_0\|$, so we normalise $\|v_0\|=1$,
$v_0=e_3$. By Theorem~\ref{app:farfield}, $\widetilde h=\pi^2W$ with
\[
W_i(y)=\frac{\delta_{i3}}{r}+\frac{y_3y_i}{r^3},\qquad r=|y| .
\]
$W_i$ is locally integrable, homogeneous of degree $-1$ and smooth away
from the origin, so Remark~\ref{app:classical-homog} applies with
$a=b=1$ and gives
\[
\tau_{ij}=h_i*h_j=(2\pi)^{-3}\,\mathcal F\big[\widetilde h_i\widetilde h_j\big]
=(2\pi)^{-3}\pi^4\,\mathcal F\big[W_iW_j\big],
\]
where we used $\mathcal F\widetilde h_i=(2\pi)^3h_i$ and the evenness of
$h$. (Both sides are homogeneous of degree $-1$ and locally integrable,
so no distribution supported at the origin can be added.) Expanding,
\[
W_iW_j=\frac{\delta_{i3}\delta_{j3}}{r^{2}}
+\frac{\delta_{i3}y_3y_j+\delta_{j3}y_3y_i}{r^{4}}
+\frac{y_3^{2}y_iy_j}{r^{6}} .
\]
Each bracket is locally integrable and homogeneous of degree $-2$, so by
Remark~\ref{app:classical-riesz} its transform is an honest homogeneous
distribution of degree $-1$, given by
\[
\mathcal F\Big[\frac1{r^{2}}\Big]=\frac{2\pi^{2}}{s},
\qquad
\mathcal F\Big[\frac{y_iy_j}{r^{4}}\Big]
=\pi^{2}\Big(\frac{\delta_{ij}}{s}-\frac{\zeta_i\zeta_j}{s^{3}}\Big),
\qquad
\mathcal F\Big[\frac{y_3^{2}y_iy_j}{r^{6}}\Big]
=\frac{\pi^{2}}{12}\,\partial_3^{2}\partial_i\partial_j\,s^{3},
\]
where $s:=|\zeta|$. (Indeed $\mathcal F[r^{-2}]=c_{3,2}s^{-1}$;
$\mathcal F[y_iy_jr^{-4}]=-\partial_i\partial_j(c_{3,4}s)
=\pi^2\partial_i\partial_js=\pi^2(\delta_{ij}/s-\zeta_i\zeta_j/s^3)$;
$\mathcal F[y_3^2y_iy_jr^{-6}]
=\partial_3^2\partial_i\partial_j(c_{3,6}s^3)$. Trace check on the
middle formula: $\pi^2(3/s-1/s)=2\pi^2/s=\mathcal F[r^{-2}]$, as it must
be.) Direct differentiation of $s^3$ gives
\[
\partial_3^{2}\partial_i\partial_j s^{3}
=3\Big[\frac{\delta_{ij}+2\delta_{i3}\delta_{j3}}{s}
-\frac{\zeta_3^{2}\delta_{ij}
+2\zeta_3(\delta_{i3}\zeta_j+\delta_{j3}\zeta_i)+\zeta_i\zeta_j}{s^{3}}
+\frac{3\zeta_i\zeta_j\zeta_3^{2}}{s^{5}}\Big].
\]
Assembling the three pieces (the middle one contributing
$\pi^2[2\delta_{i3}\delta_{j3}/s-\zeta_3(\delta_{i3}\zeta_j
+\delta_{j3}\zeta_i)/s^3]$) yields
$\mathcal F[W_iW_j]=\pi^{2}G_{ij}$ with
\[
G_{ij}=\frac{\delta_{ij}}{4s}
+\frac{9\,\delta_{i3}\delta_{j3}}{2s}
-\frac{\zeta_3^{2}\delta_{ij}}{4s^{3}}
-\frac{3\zeta_3(\delta_{i3}\zeta_j+\delta_{j3}\zeta_i)}{2s^{3}}
-\frac{\zeta_i\zeta_j}{4s^{3}}
+\frac{3\zeta_i\zeta_j\zeta_3^{2}}{4s^{5}} ,
\]
so that $\tau_{ij}=(2\pi)^{-3}\pi^{4}\cdot\pi^{2}G_{ij}
=\frac{\pi^{3}}{8}G_{ij}$.

Contract with $\zeta_j$, using $\zeta_j\zeta_j=s^2$. The terms
$\frac{\delta_{ij}}{4s}\zeta_j=\frac{\zeta_i}{4s}$ and
$-\frac{\zeta_i\zeta_j}{4s^3}\zeta_j=-\frac{\zeta_i}{4s}$ cancel. The
$\delta_{i3}\zeta_3/s$ coefficients are $\frac92$ from the second term
and $-\frac32$ from the fourth, total $3$. The $\zeta_i\zeta_3^2/s^3$
coefficients are $-\frac14$ (third term), $-\frac32$ (fourth term) and
$+\frac34$ (sixth term), total $-1$. Hence
\[
G(\zeta)\zeta=\frac{3\zeta_3}{s}\,e_3-\frac{\zeta_3^{2}}{s^{3}}\,\zeta ,
\qquad\text{so}\qquad
\zeta\times\big(G(\zeta)\zeta\big)=\frac{3\zeta_3}{s}\,(\zeta\times e_3),
\]
the second term being parallel to $\zeta$. Multiplying by
$\frac{\pi^3}{8}$ and restoring $\|v_0\|^2$ gives the claim. Finally
$\zeta\times e_3=0$ exactly when $\zeta\parallel e_3$.
\end{proof}

\begin{remark}[physical-space cross-check, free of any regularisation]
\label{app:crosscheck}
Theorem~\ref{app:tauformula} passes through the finite-part
continuations of Remark~\ref{app:classical-riesz}. The following
elementary computation confirms it independently, with no
regularisation and no Fourier transform, and is worth recording because
it identifies what the obstruction \emph{is}: the statement that the
Oseen field is not a stationary Euler flow.

For $W(y)=\frac{e_3}{r}+\frac{y_3y}{r^{3}}$ on $\R^3\setminus\{0\}$,
\[
\nabla\cdot W=0,\qquad
\omega:=\nabla\times W=\frac{2\,(e_3\times y)}{r^{3}},\qquad
\nabla\times\big(W\times\omega\big)
=\frac{12\,y_3\,(e_3\times y)}{r^{6}}\ \not\equiv0 .
\]
\emph{Proof.} $\nabla\cdot(e_3/r)=e_3\cdot\nabla r^{-1}=-y_3/r^3$, and
$\partial_i(y_3y_ir^{-3})=y_3r^{-3}+3y_3r^{-3}-3y_3r^{-3}=y_3/r^3$, so
$\nabla\cdot W=0$. Next
$\nabla\times(r^{-1}e_3)=\nabla(r^{-1})\times e_3=-r^{-3}y\times e_3
=r^{-3}(e_3\times y)$, while
$[\nabla\times(y_3y\,r^{-3})]_a
=\epsilon_{abi}\partial_b(y_3y_ir^{-3})
=\epsilon_{abi}(\delta_{b3}y_i+y_3\delta_{bi})r^{-3}
-3\epsilon_{abi}y_3y_iy_br^{-5}
=\epsilon_{a3i}y_ir^{-3}=(e_3\times y)_ar^{-3}$, the last two terms
vanishing by antisymmetry. Adding gives $\omega$. Then, using
$e_3\times(e_3\times y)=y_3e_3-y$ and $y\times(e_3\times y)=r^2e_3-y_3y$,
\[
W\times\omega
=\frac{2(y_3e_3-y)}{r^{4}}+\frac{2y_3e_3}{r^{4}}-\frac{2y_3^{2}y}{r^{6}}
=\frac{4y_3e_3}{r^{4}}-\frac{2y}{r^{4}}-\frac{2y_3^{2}y}{r^{6}} .
\]
Now curl each piece. For any radial $f$, $\nabla\times(f(r)y)
=\nabla f\times y=f'(r)\hat y\times y=0$, so $-2y/r^4$ contributes
nothing. With $g=-2y_3^{2}/r^{6}$ one has
$\nabla g=-4y_3e_3/r^{6}+12y_3^{2}y/r^{8}$, so
$\nabla\times(gy)=\nabla g\times y=-4y_3(e_3\times y)/r^{6}$. With
$f=4y_3/r^{4}$ one has $\nabla f=4e_3/r^{4}-16y_3y/r^{6}$, so
$\nabla\times(fe_3)=\nabla f\times e_3
=-16y_3(y\times e_3)/r^{6}=16y_3(e_3\times y)/r^{6}$. The total is
$12y_3(e_3\times y)/r^{6}$. Since
$W\cdot\nabla W=\nabla\frac{|W|^2}{2}-W\times\omega$, it follows that
$\nabla\times(W\cdot\nabla W)\not\equiv0$, i.e.\ $W\cdot\nabla W$ is not
a gradient, i.e.\ $\PP(W\cdot\nabla W)\neq0$. $\square$

\emph{Consistency with Theorem~\ref{app:tauformula}.} The curl of
$\PP(\widetilde h\cdot\nabla\widetilde h)$ has symbol
$i\zeta\times\big(iP_\zeta\tau\zeta\big)=-\zeta\times(\tau\zeta)$, using
$\zeta\times P_\zeta A=\zeta\times A$. With $\|v_0\|=1$,
Theorem~\ref{app:tauformula} makes this
$-\frac{3\pi^{3}}{8}\zeta_3(\zeta\times e_3)/|\zeta|$. Writing its
$a$-component as $-\frac{3\pi^3}{8}\epsilon_{ab3}\zeta_b\zeta_3/|\zeta|$
and using $\widetilde{|\zeta|^{-1}}=c_{3,1}r^{-2}=4\pi r^{-2}$ together
with $\widetilde{\zeta_b\zeta_3G}=-\partial_b\partial_3\widetilde G$,
\[
\widetilde{\Big[\frac{\zeta_b\zeta_3}{|\zeta|}\Big]}
=-4\pi\,\partial_b\partial_3r^{-2}
=\frac{8\pi\delta_{b3}}{r^{4}}-\frac{32\pi y_3y_b}{r^{6}} ;
\]
since $\epsilon_{ab3}\delta_{b3}=0$, only the second term survives,
leaving $-32\pi y_3(y\times e_3)_a/r^{6}
=+32\pi y_3(e_3\times y)_a/r^{6}$. Multiplying by
$-\frac{3\pi^{3}}{8}$ gives $-12\pi^{4}y_3(e_3\times y)/r^{6}$. On the
other hand $\widetilde h=\pi^2W$ gives
$\widetilde h\cdot\nabla\widetilde h=\pi^{4}W\cdot\nabla W$ and hence
$\nabla\times(\widetilde h\cdot\nabla\widetilde h)
=-\pi^{4}\nabla\times(W\times\omega)
=-12\pi^{4}y_3(e_3\times y)/r^{6}$. The two derivations agree exactly,
constants included.
\end{remark}

\subsubsection*{Transfer to the cutoff family}

\begin{lemma}[the cutoff perturbation is bounded near $\zeta=0$]
\label{app:perturb}
For all $|\zeta|\le\frac18$ and all $i,j$,
\[
\big|T_{ij}(\zeta)-\tau_{ij}(\zeta)\big|\ \le\ C_1:=64\pi\|v_0\|^{2},
\]
uniformly over all admissible $\chi$.
\end{lemma}

\begin{proof}
Since $F=\chi(|\cdot|)h$,
\[
T_{ij}(\zeta)-\tau_{ij}(\zeta)
=\int_{\R^3}h_i(\eta)h_j(\zeta-\eta)
\big[\chi(|\eta|)\chi(|\zeta-\eta|)-1\big]\,d\eta .
\]
The bracket vanishes whenever $|\eta|\le\frac12$ \emph{and}
$|\zeta-\eta|\le\frac12$, because $\chi\equiv1$ on $[0,\frac12]$. For
$|\zeta|\le\frac18$, the condition $|\eta|\le\frac14$ implies
$|\zeta-\eta|\le\frac18+\frac14=\frac38<\frac12$; hence the integrand is
supported in $\{|\eta|\ge\frac14\}$. There
$|\zeta-\eta|\ge|\eta|-\frac18\ge|\eta|-\frac{|\eta|}{2}
=\frac{|\eta|}{2}$. With $|h|\le\|v_0\||\cdot|^{-2}$ and
$|\chi\chi-1|\le1$,
\[
|T_{ij}-\tau_{ij}|
\le\|v_0\|^{2}\int_{|\eta|\ge1/4}\frac{d\eta}{|\eta|^{2}|\zeta-\eta|^{2}}
\le4\|v_0\|^{2}\int_{|\eta|\ge1/4}\frac{d\eta}{|\eta|^{4}}
=4\|v_0\|^{2}\cdot4\pi\int_{1/4}^{\infty}\frac{dr}{r^{2}}
=64\pi\|v_0\|^{2},
\]
which is $C_1$.
\end{proof}

\begin{theorem}[non-degeneracy]\label{app:vnondeg}
For every $v_0\in\R^3\setminus\{0\}$ and every admissible $\chi$,
\[
\PP_{\R^3}\big(V\cdot\nabla V\big)\ \not\equiv\ 0 .
\]
Explicitly, with $v_0=\|v_0\|e_3$, one has
$\zeta\times(T(\zeta)\zeta)\neq0$ --- hence $P_\zeta M(\zeta)\neq0$ ---
at every point of the open set of positive measure
\[
\mathcal U:=\Big\{\zeta\in\R^3:\ 0<|\zeta|<\tfrac{\pi^{2}}{2048},\
|\zeta_3|\ge\tfrac{|\zeta|}{2},\
|\zeta\times e_3|\ge\tfrac{|\zeta|}{2}\Big\}.
\]
\end{theorem}

\begin{proof}
Let $\zeta\in\mathcal U$. By Theorem~\ref{app:tauformula},
\[
\big|\zeta\times(\tau(\zeta)\zeta)\big|
=\frac{3\pi^{3}}{8}\|v_0\|^{2}\,
\frac{|\zeta_3|\,|\zeta\times e_3|}{|\zeta|}
\ \ge\ \frac{3\pi^{3}}{8}\|v_0\|^{2}\cdot\frac{|\zeta|}{4}
=\frac{3\pi^{3}}{32}\|v_0\|^{2}|\zeta| .
\]
Since $\frac{\pi^{2}}{2048}\approx4.8\times10^{-3}<\frac18$,
Lemma~\ref{app:perturb} applies, and the Frobenius norm of
$T(\zeta)-\tau(\zeta)$ is at most $3C_1$ (nine entries each at most
$C_1$). Hence
\[
\big|\zeta\times\big((T(\zeta)-\tau(\zeta))\zeta\big)\big|
\le|\zeta|\cdot3C_1|\zeta|=192\pi\|v_0\|^{2}|\zeta|^{2},
\]
and therefore
\[
\big|\zeta\times(T(\zeta)\zeta)\big|
\ \ge\ \|v_0\|^{2}|\zeta|
\Big(\frac{3\pi^{3}}{32}-192\pi|\zeta|\Big)\ >\ 0
\qquad\text{whenever }
|\zeta|<\frac{3\pi^{3}}{32\cdot192\pi}=\frac{3\pi^{2}}{6144}
=\frac{\pi^{2}}{2048},
\]
which holds on $\mathcal U$. The set $\mathcal U$ is open and nonempty
(it contains, e.g., all small $\zeta$ in the direction
$(1,0,1)/\sqrt2$), hence of positive measure. By
Lemma~\ref{app:symbol}, $P_\zeta M\neq0$ on a set of positive measure,
so $\PP(V\cdot\nabla V)\neq0$ in $L^2(\R^3)$.
\end{proof}

\begin{remark}[uniformity, and where the obstruction lives]
\label{app:uniformity}
The proof of Theorem~\ref{app:vnondeg} is uniform in $\chi$ and in
$v_0$: the factor $\|v_0\|^2$ cancels between the main term and the
perturbation, and the only properties of $\chi$ used are $0\le\chi\le1$
and $\chi\equiv1$ on $[0,\frac12]$ --- no smoothness of $\chi$ enters
\S\ref{app:nondeg} at all. Neither a specific $\chi$ nor any numerically
evaluated convolution integral is required: the obstruction is carried
entirely by the $\chi$-independent leading singularity of $F*F$ at
$\zeta=0$, which is the same for every admissible cutoff. Moreover
Theorem~\ref{app:tauformula} shows the obstruction degenerates exactly
on the two directions $\zeta\parallel v_0$ and $\zeta\perp v_0$, so the
two cone conditions in $\mathcal U$ are essentially optimal.

This uniformity does \textbf{not} propagate to the constants of
Theorem~\ref{app:main}. It is the \emph{qualitative} conclusion
$\PP(V\cdot\nabla V)\not\equiv0$ that holds uniformly in $(\chi,v_0)$,
and that is what licenses the quantifier ``for every admissible $\chi$
and every $v_0\neq0$'' there. The quantitative constants $c_0$ and
$N_*$ are extracted only after a test field $\Psi$ has been chosen for
the particular profile $V=V^{\chi,v_0}$
(Lemma~\ref{app:testexists}), and they depend on $(\chi,v_0)$; no
$(\chi,v_0)$-independent value of either is claimed, and neither is
effective (Remark~\ref{app:noneffective}).
\end{remark}

\subsection{Proof of Theorem~\ref{capacity}}
\label{app:final}

\begin{theorem}[capacity bound; restatement of Theorem~\ref{capacity}]
\label{app:main}
Let $v_0\in\R^3\setminus\{0\}$, let $\chi$ be admissible
(Definition~\ref{app:admissible}), and let $u_N=u_N^{\chi,v_0}$ be the
family \eqref{app:famdef}. Then:
\begin{enumerate}
\item[(1)] $u_N$ is a real, zero-mean, exactly divergence-free
trigonometric polynomial banded at $|k|\le N$, and the exact laws of
Theorem~\ref{app:exactlaws} hold;
\item[(2)] $\dfrac{N}{44544}\le N_0^2(u_N)\le\dfrac{64}{3}N$ for every
$N\ge32$;
\item[(3)] there exist $c_0=c_0(\chi,v_0)>0$ and
$N_*=N_*(\chi,v_0)<\infty$ such that
\[
\big\|\PP(u_N\cdot\nabla u_N)\big\|_{L^2(\T^3)}^{2}\ \ge\ c_0\,N^{3}
\qquad\text{for all }N\ge N_* .
\]
\end{enumerate}
\end{theorem}

\begin{proof}
(1) is Lemma~\ref{app:structure} together with
Theorem~\ref{app:exactlaws}. (2) is Proposition~\ref{app:twosided}.

(3) By Theorem~\ref{app:vnondeg}, $\PP_{\R^3}(V\cdot\nabla V)\not\equiv0$,
where $V$ is the profile of Definition~\ref{app:Vdef} attached to the
given $(\chi,v_0)$. By Lemma~\ref{app:testexists} there is a real
$\Psi\in C^\infty_c(\R^3;\R^3)$, divergence-free, supported in
$\{|y|\le R\}$ for some $R\ge1$, with $I_\Psi\neq0$. Fix such a $\Psi$.
Lemma~\ref{app:innerexact} and Theorem~\ref{app:rate} give the exact
inner rescaling with $\varepsilon_N+\varepsilon'_N=O(N^{-1}\log N)$;
Lemma~\ref{app:psiprops} and Theorem~\ref{app:pairing} convert the
pairing of $\PP(u_N\cdot\nabla u_N)$ against the concentrated field
$\psi_N$ into $(2\pi)^{-3}N^{3/2}(I_\Psi+\mathcal E_N)$ with
$\mathcal E_N\to0$; and Proposition~\ref{app:N3} then yields the stated
bound with
\[
c_0=c_0(\chi,v_0)=\frac{|I_\Psi|^{2}}
{2\,(2\pi)^{3}\,\|\Psi\|_{L^2(\R^3)}^{2}},
\qquad
N_*=N_*(\chi,v_0)
=\min\Big\{N_0\ge\max(8R,32):|\mathcal E_N|\le(1-2^{-1/2})|I_\Psi|
\ \forall N\ge N_0\Big\}.
\]
Both depend on $(\chi,v_0)$ alone, through $V=V^{\chi,v_0}$ and through
the test field $\Psi$ selected for it.
\end{proof}

\begin{remark}[what the constants depend on, and that they are
non-effective]\label{app:noneffective}
The constants $c_0$ and $N_*$ produced above depend on the pair
$(\chi,v_0)$ only, through the profile $V=V^{\chi,v_0}$ and through the
test field $\Psi$ selected in Lemma~\ref{app:testexists}. They are
\textbf{not effective}. The reason is localised precisely at
Lemma~\ref{app:testexists}: $\Psi$ is obtained from the \emph{density}
of $C^\infty_{c,\sigma}(\R^3)$ in $L^2_\sigma(\R^3)$
(Remark~\ref{app:classical-density}), a pure existence argument, so
neither $\Psi$, nor its support radius $R$, nor $I_\Psi$, nor
$\|\Psi\|_{L^2}$, nor $c_0$, nor $N_*$ is exhibited. Every other
constant in this appendix is explicit, namely those of
Lemmas~\ref{app:shell}--\ref{app:lower},
Proposition~\ref{app:twosided}, Theorems~\ref{app:exactlaws},
\ref{app:Vprops}, \ref{app:farfield}, \ref{app:rate},
\ref{app:pairing}, \ref{app:tauformula}, \ref{app:vnondeg},
Lemmas~\ref{app:innerexact}, \ref{app:psiprops}, \ref{app:symbol},
\ref{app:tauwd}, \ref{app:perturb}, and Proposition~\ref{app:N3}
\emph{as a function of} $(\Psi,R,I_\Psi)$.

An effective variant is available in principle --- take $\Psi$ with
$\widehat\Psi$ a smooth bump supported in the explicit set $\mathcal U$
of Theorem~\ref{app:vnondeg}, mollified to compact support in $y$ ---
but it is expensive: non-vanishing is certified only for
$|\zeta|<\pi^2/2048\approx4.8\times10^{-3}$, so such a $\Psi$ has
Fourier support at that scale, hence spatial radius $R\gtrsim10^{3}$,
hence $N_*\ge8R\gtrsim10^{4}$, with a correspondingly small $c_0$. This
is the sole source of non-effectivity in the paper.
\end{remark}

\begin{remark}[the exponent $3$ is sharp, and the sharp cutoff is not
covered]\label{app:sharpness}
The exponent cannot be improved: by Theorem~\ref{main}(d) and
Theorem~\ref{app:exactlaws},
$\|\PP(u_N\cdot\nabla u_N)\|_2^2\le\|u_N\|_\infty^2H_1(u_N)$, and both
$\|u_N\|_\infty$ and $H_1(u_N)$ are $O(N)$ by
Proposition~\ref{app:twosided} and
$\|u_N\|_\infty\le\sum_{k\neq0}|\hat u_N(k)|\le\|v_0\|\Sigma_N^\chi
\le\|v_0\|S_N\le128\|v_0\|N$; so the left-hand side is $O(N^3)$.

Smoothness of $\chi$ is used in exactly one place, namely
Theorem~\ref{app:rate}, through $|\nabla F|\lesssim c_\chi|\xi|^{-3}$.
For the sharply truncated family $\chi=\mathbf 1_{[0,1]}$ this bound
fails, $F$ jumping across $|\xi|=1$; the Riemann-sum defect then carries
a lattice-point (Gauss-circle) term, and $V$ loses the rapidly decaying
tail correction of Theorem~\ref{app:farfield}, acquiring instead an
oscillatory $|y|^{-2}$ tail. Nothing in \S\ref{app:duality} or
\S\ref{app:nondeg} would change, both being insensitive to $\chi$; the
sharp-cutoff case is thus a technical gap in the inner-rescaling step,
not a structural one. It is left open because no statement in this paper
requires it (see Remark~\ref{lstarremark}).
\end{remark}
  inside the paper's \appendix.

\section{Complete proof of the capacity bound (Theorem~\ref{capacity})}
\label{app:capacity}

This appendix proves Theorem~\ref{capacity} in full. Nothing below is
quoted from elsewhere except four classical facts, each isolated and
labelled as such (Remarks~\ref{app:classical-riesz},
\ref{app:classical-homog}, \ref{app:classical-density} and
\ref{app:classical-leray}); no numerical statement is used anywhere.

The chain is: the family $u_N$ has an exact scaling
$\hat u_N(k)=N^{-2}F(k/N)$ (\S\ref{app:profile}), so that the lattice
sums reconstructing $u_N$ near the origin are \emph{literally} Riemann
sums, of spacing $1/N$, for the inverse Fourier transform $V$ of $F$
(\S\ref{app:riemann}); pairing $\PP(u_N\cdot\nabla u_N)$ against a test
field concentrated at scale $1/N$ therefore converts, with an explicit
error, into the fixed continuum pairing
$\int_{\R^3}(V\cdot\nabla V)\cdot\Psi$ times $N^{3/2}$
(\S\ref{app:duality}); and that continuum pairing is nonzero for some
$\Psi$ because $\PP_{\R^3}(V\cdot\nabla V)\not\equiv0$, which is the
Fourier form of the elementary fact that the Oseen field is not a
stationary Euler flow (\S\ref{app:nondeg}). The bookkeeping is assembled
in \S\ref{app:final}.

\subsection{Conventions}
\label{app:conventions}

On $\T^3=(\R/2\pi\Z)^3$ we use the normalised measure
$d\mu=(2\pi)^{-3}dx$, so that for $u:\T^3\to\R^3$
\[
u(x)=\sum_{k\in\Z^3}\hat u(k)e^{ik\cdot x},\qquad
\hat u(k)=\int_{\T^3}u(x)e^{-ik\cdot x}\,d\mu(x),
\]
\[
\langle f,g\rangle=\int_{\T^3}f\cdot\bar g\,d\mu
=\sum_{k\in\Z^3}\hat f(k)\cdot\overline{\hat g(k)},
\qquad
H_r(u)=\sum_{k\in\Z^3}|k|^{2r}|\hat u(k)|^2 .
\]
On $\R^3$ we use the forward transform and the inverse-transform
bracket
\[
\mathcal F f(\zeta)=\int_{\R^3}f(y)e^{-i\zeta\cdot y}\,dy,
\qquad
\widetilde G(y)=\int_{\R^3}G(\xi)e^{i\xi\cdot y}\,d\xi ,
\]
so that $\mathcal F\widetilde G=(2\pi)^3G$ and
$\|\widetilde G\|_{L^2(\R^3)}^2=(2\pi)^3\|G\|_{L^2(\R^3)}^2$.

For $\xi\in\R^3\setminus\{0\}$ put
$P_\xi=I-\dfrac{\xi\otimes\xi}{|\xi|^2}$. Thus $P_\xi$ is the orthogonal
projection onto $\xi^\perp$: it is symmetric, idempotent, $P_\xi\xi=0$,
$\|P_\xi\|\le1$, $P_{-\xi}=P_\xi$, it is homogeneous of degree $0$ and
smooth on $\R^3\setminus\{0\}$, and consequently
$P_{k}=P_{k/N}$ for every $k\neq0$ and every $N>0$. The Leray
projection $\PP$ acts on $\T^3$ by
$\widehat{\PP f}(k)=P_k\hat f(k)$ for $k\neq0$, leaving the mean
untouched, and on $\R^3$ by the same symbol $P_\zeta$. On both
$L^2(\T^3)$ and $L^2(\R^3)$, $\PP$ is an orthogonal projection, hence
self-adjoint.

Throughout, $\|v_0\|$ is the Euclidean norm of $v_0\in\R^3$, and
$e_1,e_2,e_3$ is the standard basis. Repeated Latin indices are summed.
One notational warning: in this appendix $V$ always denotes the
continuum profile of Definition~\ref{app:Vdef}; the unrelated use of
$V$ for the modal variance $\mathrm{Var}_p(|k|^2)$ in
Theorem~\ref{main}(a) does not occur here.

\subsection{The cutoff class, the family, and its exact laws}
\label{app:laws}

We restate the definitions of the main text for completeness.

\begin{definition}[admissible cutoff; restatement of
Definition~\ref{admissible}]\label{app:admissible}
$\chi\in C^4([0,\infty);[0,1])$ with $\chi\equiv1$ on
$[0,\tfrac12]$ and $\operatorname{supp}\chi\subset[0,1]$. No
monotonicity is assumed --- it is used nowhere. Write
$c_\chi:=1+\|\chi'\|_{L^\infty}\ (\ge1)$.
\end{definition}

Such $\chi$ exist in abundance. The standard mollifier construction
gives a $C^\infty$ example, $\chi(t)=\Theta\big(2t-1\big)$ with
$\Theta(s)=1$ for $s\le0$, $\Theta(s)=0$ for $s\ge1$ and
$\Theta(s)=\frac{e^{-1/(1-s)}}{e^{-1/s}+e^{-1/(1-s)}}$ on $(0,1)$; the
exactly rational cutoff of Appendix~\ref{app:certificates}, built from
the degree-$9$ smoothstep, lies in $C^4\setminus C^5$ and is likewise
admissible.

\begin{remark}[what smoothness is actually used]\label{app:c1}
Definition~\ref{app:admissible} is still stronger than the proof
requires. The only place any derivative of $\chi$ is used
load-bearingly is Theorem~\ref{app:rate}, and it enters solely through
$c_\chi=1+\|\chi'\|_\infty$; Sections~\ref{app:duality} and
\ref{app:nondeg} use only $0\le\chi\le1$, $\chi\equiv1$ on
$[0,\tfrac12]$ and $\operatorname{supp}\chi\subset[0,1]$, and
Sections~\ref{app:laws}, \ref{app:lattice}, \ref{app:profile} use no
derivative of $\chi$ at all. Every statement of this appendix on which
Theorem~\ref{app:main} depends therefore holds verbatim for
$\chi\in C^1([0,\infty);[0,1])$ --- indeed for Lipschitz $\chi$ ---
with those support properties. The four derivatives are demanded only
so that Theorem~\ref{app:farfield} holds as stated with $M=4$; that
theorem is consumed only by the inert Remark~\ref{app:noperiod}, and
nothing in the proof of Theorem~\ref{app:main} uses it. In particular
$V\in C^\infty$ (Theorem~\ref{app:Vprops}(3)) and the Paley--Wiener
statement there are bought by the \emph{compact support} of $F$, not by
smoothness of $\chi$, and are unaffected.
\end{remark}

\begin{definition}\label{app:family}
Fix $v_0\in\R^3\setminus\{0\}$, an admissible $\chi$, and an integer
$N\ge2$. Define $u_N=u_N^{\chi,v_0}$ on $\T^3$ by
\begin{equation}\label{app:famdef}
\hat u_N(k)=\chi\Big(\frac{|k|}{N}\Big)\frac{P_kv_0}{|k|^{2}}
\quad(k\neq0),\qquad \hat u_N(0)=0 .
\end{equation}
\end{definition}

Since $\hat u_N(k)=0$ for $|k|>N$, $u_N$ is a finite trigonometric
polynomial banded at $|k|\le N$; its coefficients are real vectors and
no integrality of $v_0$ is required.

\begin{lemma}[basic structure]\label{app:structure}
$u_N$ is real, zero-mean, exactly divergence-free, $C^\infty$, and
$u_N\not\equiv0$.
\end{lemma}

\begin{proof}
Zero mean is $\hat u_N(0)=0$. Reality: the coefficients are real
vectors and $P_{-k}=P_k$, so
$\hat u_N(-k)=\hat u_N(k)=\overline{\hat u_N(k)}$, which is exactly the
reality condition $\hat u(-k)=\overline{\hat u(k)}$. Divergence-free:
$\widehat{\nabla\cdot u_N}(k)=ik\cdot\hat u_N(k)$ and $k\cdot P_kv_0=0$
for every $k\neq0$. Smoothness: a finite sum of exponentials. Finally,
among $e_1,e_2,e_3$ at least two are not parallel to $v_0$; pick such a
$k$ with $|k|=1$. Then $\chi(1/N)=1$ because $1/N\le\tfrac12$, and
$P_kv_0\neq0$, so $\hat u_N(k)\neq0$.
\end{proof}

The exact laws rest on a band-symmetry identity which, because $\chi$ is
not assumed monotone and the truncation is not a sharp ball, must be
available for a \emph{general} radial weight, with no positivity
assumption.

\begin{lemma}[band symmetry, lattice, general radial weight]
\label{app:bandsym}
Let $f:\Z^3\setminus\{0\}\to\R$ be radial, i.e.\ $f(k)=\varphi(|k|^2)$
for some $\varphi$, and suppose $\sum_{k\neq0}|k|^2|f(k)|<\infty$. Then
for all $i,j$
\[
\sum_{k\neq0}k_ik_j\,f(k)
=\delta_{ij}\,\frac13\sum_{k\neq0}|k|^2f(k).
\]
\end{lemma}

\begin{proof}
Absolute summability of $|k|^2f$ dominates every summand
$|k_ik_jf(k)|\le|k|^2|f(k)|$, so all rearrangements below are
legitimate. The index set $\Z^3\setminus\{0\}$ is invariant under each
sign flip $\sigma_i:k_i\mapsto-k_i$ and under all coordinate
permutations, and $f$, depending on $|k|^2$ only, is invariant under
both. For $i\neq j$ the bijection $\sigma_i$ negates the summand
$k_ik_jf(k)$, so that sum equals its own negative and vanishes. For
$i=j$, coordinate permutations carry the three sums $\sum_kk_i^2f(k)$
into one another, so they are equal; their sum is $\sum_k|k|^2f(k)$.
\end{proof}

\begin{lemma}[band symmetry, continuum]\label{app:bandsymcont}
Let $g:\R^3\to\R$ be radial with $\int_{\R^3}|\xi|^2|g(\xi)|\,d\xi<\infty$.
Then $\int\xi_i\xi_jg\,d\xi=\delta_{ij}\frac13\int|\xi|^2g\,d\xi$.
\end{lemma}

\begin{proof}
Identical: Lebesgue measure on $\R^3$ and $g$ are both invariant under
sign flips and permutations of the coordinates, and the hypothesis
dominates the integrand.
\end{proof}

\begin{theorem}[exact family laws]\label{app:exactlaws}
Set
\[
S_N^\chi:=\sum_{k\neq0}\frac{\chi(|k|/N)^2}{|k|^2},\qquad
T_N^\chi:=\sum_{k\neq0}\frac{\chi(|k|/N)^2}{|k|^4},\qquad
\Sigma_N^\chi:=\sum_{k\neq0}\frac{\chi(|k|/N)}{|k|^2},
\]
all three finite sums. Then
\[
H_0(u_N)=\tfrac23\|v_0\|^2\,T_N^\chi,\qquad
H_1(u_N)=\tfrac23\|v_0\|^2\,S_N^\chi,\qquad
N_0^2(u_N)=\frac{S_N^\chi}{T_N^\chi},\qquad
u_N(0)=\tfrac23\,\Sigma_N^\chi\,v_0 ,
\]
and consequently
$\|u_N\|_\infty^2/H_1(u_N)\ge 2(\Sigma_N^\chi)^2/(3S_N^\chi)$.
\end{theorem}

\begin{proof}
Since $P_k$ is an orthogonal projection,
$|P_kv_0|^2=\|v_0\|^2-(k\cdot v_0)^2/|k|^2$. Hence
\[
H_0=\sum_{k\neq0}\frac{\chi(|k|/N)^2}{|k|^4}\,|P_kv_0|^2
=\|v_0\|^2\sum_{k\neq0}\frac{\chi^2}{|k|^4}
-\sum_{k\neq0}\frac{\chi^2\,(k\cdot v_0)^2}{|k|^6}.
\]
Apply Lemma~\ref{app:bandsym} with the radial weight
$f(k)=\chi(|k|/N)^2|k|^{-6}$ (a finite sum, so the hypothesis is
trivially met):
\[
\sum_{k\neq0}\frac{\chi^2(k\cdot v_0)^2}{|k|^6}
=v_{0i}v_{0j}\sum_{k\neq0}k_ik_jf(k)
=\frac{\|v_0\|^2}{3}\sum_{k\neq0}\frac{\chi^2}{|k|^4}
=\frac{\|v_0\|^2}{3}T_N^\chi .
\]
Therefore $H_0=\|v_0\|^2T_N^\chi-\frac13\|v_0\|^2T_N^\chi
=\frac23\|v_0\|^2T_N^\chi$. The same computation with
$f=\chi^2|k|^{-4}$ applied to
$H_1=\sum_{k}|k|^2|\hat u_N(k)|^2$ gives
$H_1=\frac23\|v_0\|^2S_N^\chi$, and $N_0^2=H_1/H_0$ follows. For the
point value,
\[
u_N(0)=\sum_{k\neq0}\hat u_N(k)
=v_0\,\Sigma_N^\chi-\sum_{k\neq0}\frac{\chi\,k\,(k\cdot v_0)}{|k|^4},
\]
and Lemma~\ref{app:bandsym} with $f=\chi|k|^{-4}$ turns the last sum
into $\frac13v_0\Sigma_N^\chi$; hence
$u_N(0)=\frac23\Sigma_N^\chi v_0$. Finally
$\|u_N\|_\infty\ge|u_N(0)|=\frac23\Sigma_N^\chi\|v_0\|$, so
$\|u_N\|_\infty^2/H_1\ge\frac49(\Sigma_N^\chi)^2\|v_0\|^2
\big/\big(\frac23\|v_0\|^2S_N^\chi\big)
=\frac23(\Sigma_N^\chi)^2/S_N^\chi$.
\end{proof}

\subsection{Two-sided lattice bounds for the bandwidth}
\label{app:lattice}

The constants below are crude and explicit; no sharpness is attempted.

\begin{lemma}[dyadic shell count]\label{app:shell}
For every integer $j\ge0$,
$\#\{k\in\Z^3:2^j\le|k|<2^{j+1}\}\le64\cdot8^{j}$.
\end{lemma}

\begin{proof}
Such $k$ satisfy $|k|_\infty\le|k|<2^{j+1}$, hence
$|k_i|\le2^{j+1}-1$ for each $i$, so the count is at most
$(2^{j+2}-1)^3<2^{3j+6}=64\cdot8^{j}$.
\end{proof}

\begin{lemma}[upper bounds]\label{app:upper}
With $S_M:=\sum_{1\le|k|\le M}|k|^{-2}$ and
$T_\infty:=\sum_{k\neq0}|k|^{-4}$ one has
$S_M\le128M$ for every real $M\ge1$, and $T_\infty\le128$.
\end{lemma}

\begin{proof}
Let $J=\lfloor\log_2M\rfloor$. Every $k$ with $1\le|k|\le M$ lies in a
shell $2^j\le|k|<2^{j+1}$ with $0\le j\le J$, where $|k|^{-2}\le4^{-j}$.
By Lemma~\ref{app:shell},
$S_M\le\sum_{j=0}^{J}64\cdot8^{j}4^{-j}=64(2^{J+1}-1)<128\cdot2^{J}\le128M$.
Similarly $T_\infty\le\sum_{j\ge0}64\cdot8^{j}16^{-j}=64\cdot2=128$.
\end{proof}

\begin{lemma}[lower bound]\label{app:lower}
For every real $M\ge16$ one has $S_M\ge M/87$.
\end{lemma}

\begin{proof}
Consider
$\mathcal B:=\{k\in\Z^3:\ \tfrac{M}{2\sqrt3}<k_i\le\tfrac{M}{\sqrt3}\
\text{for }i=1,2,3\}$. For $k\in\mathcal B$ we have
$|k|^2\le3(M/\sqrt3)^2=M^2$ and $|k|^2>3M^2/12=M^2/4$, so every
$k\in\mathcal B$ is counted in $S_M$ (in particular $k\neq0$) and
contributes $|k|^{-2}\ge M^{-2}$. Each coordinate ranges over a
half-open interval of length $M/(2\sqrt3)$, which contains at least
$\frac{M}{2\sqrt3}-1$ integers; for $M\ge16$,
$\frac{M}{2\sqrt3}-1\ge0.28868M-0.0625M\ge0.226M$. Hence
$\#\mathcal B\ge(0.226M)^3>0.01154M^3$ and
$S_M\ge0.01154M^3\cdot M^{-2}=0.01154M\ge M/87$.
\end{proof}

\begin{proposition}[two-sided bandwidth bounds]\label{app:twosided}
For every admissible $\chi$, every $v_0\neq0$ and every integer
$N\ge32$,
\[
\frac{N}{348}\ \le\ S_N^\chi\ \le\ 128N,\qquad
6\ \le\ T_N^\chi\ \le\ 128,\qquad
\frac{N}{44544}\ \le\ N_0^2(u_N)\ \le\ \frac{64}{3}N .
\]
\end{proposition}

\begin{proof}
Upper bounds: $0\le\chi\le1$ and $\chi(|k|/N)=0$ for $|k|>N$, so
$S_N^\chi\le S_N\le128N$ and $T_N^\chi\le T_\infty\le128$ by
Lemma~\ref{app:upper}. Lower bound for $T_N^\chi$: for $N\ge2$ one has
$\chi(1/N)=1$, and the six lattice points with $|k|=1$ each contribute
$1$, so $T_N^\chi\ge6$. Lower bound for $S_N^\chi$: $\chi(|k|/N)=1$
whenever $|k|\le N/2$, hence $S_N^\chi\ge S_{N/2}$; for $N\ge32$ we have
$N/2\ge16$, so Lemma~\ref{app:lower} gives
$S_N^\chi\ge (N/2)/87\ge N/348$. Combining with
$N_0^2=S_N^\chi/T_N^\chi$ (Theorem~\ref{app:exactlaws}) yields
$N_0^2\ge\frac{N/348}{128}=\frac{N}{44544}$ and
$N_0^2\le\frac{128N}{6}=\frac{64}{3}N$.
\end{proof}

\subsection{The profile $V$}
\label{app:profile}

\begin{definition}\label{app:Vdef}
On $\R^3$ put
\[
h(\xi):=\frac{P_\xi v_0}{|\xi|^2}\quad(\xi\neq0),\qquad
F(\xi):=\chi(|\xi|)\,h(\xi)\quad(\xi\neq0),\quad F(0):=0,
\]
\[
V(y):=\widetilde F(y)=\int_{\R^3}F(\xi)e^{i\xi\cdot y}\,d\xi .
\]
\end{definition}

Because $P_\xi$ is homogeneous of degree $0$ we have
$P_{k/N}=P_k$ and therefore
$F(k/N)=\chi(|k|/N)\,N^2P_kv_0/|k|^2$, i.e.\ the exact scaling identity
\begin{equation}\label{app:scaling}
\hat u_N(k)=N^{-2}F(k/N)\qquad\text{for every }k\in\Z^3
\end{equation}
(the case $k=0$ holds by the conventions $\hat u_N(0)=0=F(0)$).
Identity~\eqref{app:scaling} is the engine of the whole argument: it
says that the family is a single fixed continuum profile sampled on the
lattice $N^{-1}\Z^3$.

\begin{theorem}[elementary properties of $V$]\label{app:Vprops}
\begin{enumerate}
\item[(1)] $|F(\xi)|\le\|v_0\|\,|\xi|^{-2}\mathbf 1_{\{|\xi|\le1\}}$;
hence $F\in L^1(\R^3)\cap L^p(\R^3)$ for every $p<3/2$, with
$\|F\|_{L^1}\le4\pi\|v_0\|$.
\item[(2)] $V$ is real, even, bounded and continuous, with
$\|V\|_\infty\le4\pi\|v_0\|$.
\item[(3)] $V\in C^\infty(\R^3)$, and for every multi-index $\beta$
\[
\partial^\beta V(y)=\int_{\R^3}(i\xi)^\beta F(\xi)e^{i\xi\cdot y}\,d\xi,
\qquad
\|\partial^\beta V\|_\infty\le\|v_0\|\int_{|\xi|\le1}|\xi|^{|\beta|-2}d\xi
=\frac{4\pi\|v_0\|}{|\beta|+1} .
\]
In particular each $\|\partial_iV_j\|_\infty\le2\pi\|v_0\|$, so
$\|\nabla V\|_\infty\le6\pi\|v_0\|$ for the Frobenius norm. Moreover $V$
extends to an entire function on $\mathbb C^3$ of exponential type
$\le1$.
\item[(4)] $\nabla\cdot V\equiv0$ on $\R^3$.
\item[(5)] $V(0)=\frac{8\pi}{3}\big(\int_0^\infty\chi(r)\,dr\big)v_0\neq0$.
\end{enumerate}
\end{theorem}

\begin{proof}
(1) $|P_\xi v_0|\le\|v_0\|$, $0\le\chi\le1$ and
$\operatorname{supp}\chi\subset[0,1]$ give the pointwise bound;
$\int_{|\xi|\le1}|\xi|^{-2}d\xi=4\pi\int_0^1dr=4\pi$, and
$\int_{|\xi|\le1}|\xi|^{-2p}d\xi<\infty$ exactly when $2p<3$.

(2) $F\in L^1$ gives, by dominated convergence, continuity of $V$ and
$\|V\|_\infty\le\|F\|_{L^1}\le4\pi\|v_0\|$. $F$ is real-valued and even
($P_{-\xi}=P_\xi$), so
$V(y)=\int F(\xi)\cos(\xi\cdot y)\,d\xi$ is real and even.

(3) By (1), $\xi^\beta F\in L^1$ for every $\beta$, since the exponent
$|\beta|-2>-3$ is locally integrable in $\R^3$ and the support is
bounded. Differentiation under the integral sign is then justified by
dominated convergence, giving the displayed formula and, by (1) again,
the bound
$\|v_0\|\int_{|\xi|\le1}|\xi|^{|\beta|-2}d\xi
=4\pi\|v_0\|\int_0^1r^{|\beta|}dr=4\pi\|v_0\|/(|\beta|+1)$. Since $F$
is an $L^1$ function supported in $\{|\xi|\le1\}$, the integral
$\int F(\xi)e^{i\xi\cdot y}d\xi$ converges for all complex $y$ and is
entire in $y$, with $|V(y)|\le\|F\|_{L^1}e^{|\operatorname{Im}y|}$;
this is exponential type $\le1$.

(4) $\nabla\cdot V(y)=\int i\,\xi\cdot F(\xi)e^{i\xi\cdot y}d\xi$ and
$\xi\cdot F(\xi)=\chi(|\xi|)\,\xi\cdot P_\xi v_0/|\xi|^2=0$ pointwise.

(5) $V(0)=\int F
=v_0\int\frac{\chi(|\xi|)}{|\xi|^2}d\xi
-\int\frac{\chi(|\xi|)\,\xi\,(\xi\cdot v_0)}{|\xi|^4}d\xi$.
Lemma~\ref{app:bandsymcont} with the radial
$g(\xi)=\chi(|\xi|)|\xi|^{-4}$ (integrable against $|\xi|^2$ by (1))
gives $\int\xi_i\xi_jg=\frac{\delta_{ij}}3\int|\xi|^2g
=\frac{\delta_{ij}}3\int\chi(|\xi|)|\xi|^{-2}d\xi$, so the second
integral equals $\frac13v_0\int\chi(|\xi|)|\xi|^{-2}d\xi$. Hence
\[
V(0)=\frac23v_0\int_{\R^3}\frac{\chi(|\xi|)}{|\xi|^2}\,d\xi
=\frac23v_0\cdot4\pi\int_0^\infty\chi(r)\,dr
=\frac{8\pi}{3}\Big(\int_0^\infty\chi\Big)v_0 ,
\]
and $\int_0^\infty\chi\ge\int_0^{1/2}1=\frac12>0$ because $\chi\ge0$ and
$\chi\equiv1$ on $[0,\tfrac12]$.
\end{proof}

\begin{remark}[structural form]\label{app:oseenform}
$F=P_\xi\big(v_0\chi(|\xi|)|\xi|^{-2}\big)$, so
$V=\PP\big(v_0\,g_\chi\big)$ where $g_\chi$ is the inverse transform of
the smoothed Newtonian symbol $\chi(|\cdot|)|\cdot|^{-2}$; solving
$\Delta\phi=g_\chi$ with $\phi$ radial exhibits $V$ in the Oseen normal
form $V(y)=a(r)v_0+b(r)(v_0\cdot\hat y)\hat y$ with
$a=\phi''+\phi'/r$, $b=\phi'/r-\phi''$. This form is not used below; it
explains why the far field of $V$ is exactly the Oseen profile
(Theorem~\ref{app:farfield}).
\end{remark}

We record here the two classical transform facts used later.

\begin{remark}[classical input: Riesz-potential transforms]
\label{app:classical-riesz}
For $0<\alpha<3$, and by analytic continuation for every $\alpha$ at
which the Gamma factors have no pole, the homogeneous function
$|y|^{-\alpha}$ on $\R^3$ is a tempered distribution whose Fourier
transform is the homogeneous distribution
\[
\mathcal F\big[|y|^{-\alpha}\big](\zeta)=c_{3,\alpha}|\zeta|^{\alpha-3},
\qquad
c_{3,\alpha}=\frac{2^{3-\alpha}\pi^{3/2}\,\Gamma\!\big(\frac{3-\alpha}{2}\big)}
{\Gamma\!\big(\frac{\alpha}{2}\big)} .
\]
See, e.g., \cite{SteinWeiss,GelfandShilov,Grafakos} (we give no section
locators, not having verified them). We use
$\Gamma(\tfrac12)=\sqrt\pi$, $\Gamma(-\tfrac12)=-2\sqrt\pi$,
$\Gamma(-\tfrac32)=\tfrac43\sqrt\pi$, which give
\[
c_{3,-1}=-8\pi,\quad c_{3,1}=4\pi,\quad c_{3,2}=2\pi^2,\quad
c_{3,4}=-\pi^2,\quad c_{3,6}=\tfrac{\pi^2}{12},
\]
together with the elementary rules
$\mathcal F[y_ay_bf]=-\partial_a\partial_b\mathcal F[f]$ and
$\mathcal F[y_ay_by_cy_df]=\partial_a\partial_b\partial_c\partial_d\mathcal F[f]$,
valid for tempered $f$.
\end{remark}

\begin{remark}[classical input: products and convolutions of homogeneous
distributions]\label{app:classical-homog}
Let $A,B\in L^1_{\rm loc}(\R^3)$ be homogeneous of degrees $-a,-b$ with
$0<a,b<3$ and $a+b<3$, and smooth away from the origin. Then $AB$ is
locally integrable and homogeneous of degree $-(a+b)$, hence tempered;
$\mathcal FA,\mathcal FB$ are locally integrable, homogeneous of degrees
$a-3,b-3$ and smooth away from the origin; the convolution
$\mathcal FA*\mathcal FB$ converges absolutely at every $\zeta\neq0$
(because $(3-a)+(3-b)>3$) and defines a locally integrable homogeneous
function of degree $3-a-b$; and the exchange formula
\[
\mathcal F[AB]=(2\pi)^{-3}\,\mathcal FA*\mathcal FB
\]
holds as tempered distributions. See, e.g.,
\cite{GelfandShilov,Grafakos}. We use this only in
Theorem~\ref{app:tauformula}, with $a=b=1$.
\end{remark}

\subsection{The far field, and why periodisation is unavailable}
\label{app:farfield-sec}

\begin{theorem}[far-field asymptotics]\label{app:farfield}
Let
\[
W(y):=\frac{v_0}{|y|}+\frac{(v_0\cdot y)\,y}{|y|^3}\qquad(y\neq0).
\]
Then $\widetilde h=\pi^2W$ as tempered distributions, and
\[
V(y)=\pi^2W(y)+\rho(y),\qquad
|\rho(y)|\le C_M(\chi,v_0)|y|^{-M}\quad\text{for }|y|\ge1
\text{ and every }2\le M\le4 .
\]
(Four derivatives are all that Definition~\ref{app:admissible}
supplies, and $M=4$ is all that Remark~\ref{app:noperiod} consumes; if
$\chi\in C^m$ with $m\ge2$ the same proof gives every
$2\le M\le m$.)
\end{theorem}

\begin{proof}
\emph{(a) The identity $\widetilde h=\pi^2W$.} Let
$\mathcal O_{ij}(y)=\frac1{8\pi}\big(\frac{\delta_{ij}}{|y|}
+\frac{y_iy_j}{|y|^3}\big)$ be the Oseen tensor, i.e.\ the fundamental
solution of the steady Stokes system
\cite{LadyzhenskayaBook,Galdi2011}, and
$h_{ij}(\xi)=\frac{\delta_{ij}}{|\xi|^2}-\frac{\xi_i\xi_j}{|\xi|^4}$, so
that $h_i=h_{ij}v_{0j}$. We verify $\mathcal F[\mathcal O_{ij}]=h_{ij}$
from Remark~\ref{app:classical-riesz}. First,
$\mathcal F[|y|^{-1}]=c_{3,1}|\zeta|^{-2}=4\pi|\zeta|^{-2}$, so
$\mathcal F\big[\tfrac{\delta_{ij}}{8\pi|y|}\big]
=\tfrac{\delta_{ij}}{2|\zeta|^{2}}$. Second, differentiating
$\partial_i|y|=y_i/|y|$ once more gives the pointwise identity
\[
\partial_i\partial_j|y|=\frac{\delta_{ij}}{|y|}-\frac{y_iy_j}{|y|^3},
\qquad\text{i.e.}\qquad
\frac{y_iy_j}{|y|^3}=\frac{\delta_{ij}}{|y|}-\partial_i\partial_j|y| .
\]
With $\mathcal F[|y|]=c_{3,-1}|\zeta|^{-4}=-8\pi|\zeta|^{-4}$ we get
$\mathcal F[\partial_i\partial_j|y|]
=(i\zeta_i)(i\zeta_j)\mathcal F[|y|]=8\pi\zeta_i\zeta_j|\zeta|^{-4}$, hence
\[
\mathcal F\Big[\frac{y_iy_j}{|y|^3}\Big]
=\frac{4\pi\delta_{ij}}{|\zeta|^2}-\frac{8\pi\zeta_i\zeta_j}{|\zeta|^4}.
\]
Therefore
$\mathcal F[\mathcal O_{ij}]
=\frac1{8\pi}\big(\frac{4\pi\delta_{ij}}{|\zeta|^2}
+\frac{4\pi\delta_{ij}}{|\zeta|^2}
-\frac{8\pi\zeta_i\zeta_j}{|\zeta|^4}\big)
=\frac{\delta_{ij}}{|\zeta|^2}-\frac{\zeta_i\zeta_j}{|\zeta|^4}=h_{ij}$.
All these are identities between homogeneous tempered distributions;
possible ambiguities are supported at the origin only and are excluded
because both sides are homogeneous of degree $-2>-3$, hence locally
integrable, with no room for a delta term. Fourier inversion
($\mathcal F\widetilde G=(2\pi)^3G$) now gives
$\widetilde h=(2\pi)^3\mathcal O\,v_0=\frac{(2\pi)^3}{8\pi}W=\pi^2W$,
where we used $8\pi\,\mathcal O_{ij}v_{0j}=W_i$.

\emph{(b) The remainder.} Write $G:=(1-\chi(|\cdot|))h$, so that
$F-h=-G$ and $\rho:=-\widetilde G$. Then $G$ vanishes on
$\{|\xi|<\tfrac12\}$ and is $C^4$ on all of $\R^3$; since $h$ is
homogeneous of degree $-2$ and smooth off the origin, and
$1-\chi(|\cdot|)$ is $C^4$ with its first four derivatives bounded, we
get
\[
|\partial^\beta G(\xi)|\le C_\beta(\chi)\|v_0\|\,\langle\xi\rangle^{-2-|\beta|}
\qquad(\xi\in\R^3,\ |\beta|\le4),
\]
which for $2\le|\beta|\le4$ places $\partial^\beta G$ in $L^1(\R^3)$
(the exponent $2+|\beta|>3$). From
the distributional identity
$y^\beta\widetilde G=i^{|\beta|}\,\widetilde{\partial^\beta G}$, the
right-hand side for $2\le|\beta|\le4$ is a bounded continuous function
with supremum at most $\|\partial^\beta G\|_{L^1}$. Choosing
$\beta=(M,0,0),(0,M,0),(0,0,M)$ for a given $2\le M\le4$ and using
$|y|^M\le3^{M/2}\max_i|y_i|^M$ shows that $\widetilde G$ agrees on
$\{|y|\ge1\}$ with a continuous function obeying
$|\widetilde G(y)|\le C_M(\chi,v_0)|y|^{-M}$. Since
$V=\widetilde F=\widetilde h-\widetilde G=\pi^2W+\rho$, the claim
follows.
\end{proof}

Thus the decay of $V$ is \emph{exactly} $|y|^{-1}$, with a
$\chi$-independent leading profile $\pi^2W$ satisfying
\begin{equation}\label{app:Wpos}
v_0\cdot W(y)=\frac{\|v_0\|^2}{|y|}+\frac{(v_0\cdot y)^2}{|y|^3}
\ \ge\ \frac{\|v_0\|^2}{|y|}\ >\ 0 .
\end{equation}

\begin{remark}[the naive periodisation argument is unavailable]
\label{app:noperiod}
It is tempting to reconstruct $u_N$ from $V$ by Poisson summation,
$u_N(x)\stackrel{?}{=}N\sum_{m\in\Z^3}V(N(x+2\pi m))$ with the $m\neq0$
terms uniformly $O(1/N)$. \textbf{The naive absolutely convergent
periodisation argument is unavailable}: the series of translates
diverges, because $V$ decays only like $|y|^{-1}$ and its leading
profile has a sign-definite spherical mean in the $v_0$ direction, so
there is no cancellation to exploit and no absolute convergence to
appeal to. Indeed, by Theorem~\ref{app:farfield} with $M=4$, for
$m\neq0$ and $x\in[-\pi,\pi)^3$ one has $|N(x+2\pi m)|\ge N\pi\ge1$ and
therefore, using \eqref{app:Wpos},
\[
v_0\cdot V\big(N(x+2\pi m)\big)
\ \ge\ \frac{\pi^2\|v_0\|^2}{N|x+2\pi m|}
-\frac{C_4\|v_0\|}{N^4|x+2\pi m|^4}.
\]
The shell $|m|_\infty=n$ contains $O(n^2)$ points with
$|x+2\pi m|\asymp n$, so $\sum_{m\neq0}|x+2\pi m|^{-4}<\infty$ while
$\sum_{m\neq0}|x+2\pi m|^{-1}=\infty$; hence
$\sum_{m\neq0}v_0\cdot V(N(x+2\pi m))=+\infty$ for every $x$ and every
$N$. The terms are moreover eventually positive, so the divergence is
not of a kind that a regular linear summation method would remove: for
$a_m\ge0$ with $\sum a_m=\infty$, monotone convergence gives
$\lim_{\varepsilon\downarrow0}\sum_me^{-\varepsilon|m|}a_m=+\infty$, so
Abel summation in particular does not converge here. We make no
assertion beyond this, and none is needed: the point is only that the
naive absolutely convergent periodisation argument is not available,
because $V$ decays only like $|y|^{-1}$ with a non-oscillating leading
profile whose spherical mean, $\tfrac43v_0/r$, is sign-definite in the
$v_0$ direction. If a periodisation formula
is wanted purely for orientation, the correct \emph{renormalised} one,
\[
u_N(x)=N\Big[V(Nx)+\sum_{m\neq0}\big(V(N(x+2\pi m))-V(2\pi Nm)\big)\Big]
+\text{const},
\]
does converge absolutely, the bracketed differences being $O(|m|^{-2})$
by Theorem~\ref{app:farfield}. \textbf{Nothing in this appendix uses any
periodisation identity}: \S\ref{app:riemann} replaces it by an exact
lattice-Riemann-sum identity, which is elementary, quantitative, and
additionally delivers the derivative statement. This remark, and with it
Theorem~\ref{app:farfield}, is inert for the proof of
Theorem~\ref{capacity}.
\end{remark}

\subsection{The inner rescaling: an exact lattice-Riemann identity with
an explicit rate}
\label{app:riemann}

For $\Phi:\R^3\to\mathbb C^d$ put
\[
\mathcal R_N[\Phi]:=N^{-3}\sum_{k\in\Z^3}\Phi(k/N)
-\int_{\R^3}\Phi(\xi)\,d\xi ,
\]
whenever both terms converge absolutely.

\begin{lemma}[exact inner rescaling]\label{app:innerexact}
For every $y\in\R^3$ and every integer $N\ge2$,
\[
u_N(y/N)=N\big(V(y)+E_N(y)\big),\qquad
\partial_iu_{N,j}(y/N)=N^2\big(\partial_iV_j(y)+E'_{N,ij}(y)\big),
\]
where, with $\Phi_y(\xi):=F(\xi)e^{i\xi\cdot y}$ and
$\Phi'_{y,ij}(\xi):=i\xi_iF_j(\xi)e^{i\xi\cdot y}$ (both set to $0$ at
$\xi=0$),
\[
E_N(y)=\mathcal R_N[\Phi_y],\qquad
E'_{N,ij}(y)=\mathcal R_N[\Phi'_{y,ij}] .
\]
These are identities, not approximations.
\end{lemma}

\begin{proof}
Both series are finite (support $|k|\le N$) and both integrals converge
absolutely by Theorem~\ref{app:Vprops}(1),(3). By \eqref{app:scaling},
\[
u_N(y/N)=\sum_{k\in\Z^3}\hat u_N(k)e^{i(k/N)\cdot y}
=\sum_{k}N^{-2}F(k/N)e^{i(k/N)\cdot y}
=N\cdot N^{-3}\sum_k\Phi_y(k/N),
\]
which equals $N\big(\int\Phi_y+\mathcal R_N[\Phi_y]\big)
=N\big(V(y)+E_N(y)\big)$, using $V(y)=\int\Phi_y$
(Definition~\ref{app:Vdef}) and $\Phi_y(0)=F(0)=0$. For the gradient,
$\partial_iu_{N,j}(x)=\sum_kik_i\hat u_{N,j}(k)e^{ik\cdot x}$; at
$x=y/N$ this is
\[
\sum_kiN^{-2}k_iF_j(k/N)e^{i(k/N)\cdot y}
=N^{-1}\sum_ki(k_i/N)F_j(k/N)e^{i(k/N)\cdot y}
=N^{2}\cdot N^{-3}\sum_k\Phi'_{y,ij}(k/N),
\]
and $\int\Phi'_{y,ij}=\partial_iV_j(y)$ by
Theorem~\ref{app:Vprops}(3).
\end{proof}

Everything now rests on bounding a Riemann-sum defect for an explicit,
compactly supported integrand with a single integrable algebraic
singularity.

\begin{theorem}[$C^1_{\rm loc}$ convergence with explicit rate]
\label{app:rate}
There is an absolute constant $A$ (the proof gives $A=3000$) such that
for every real $R\ge1$ and every integer $N\ge8R$,
\[
\varepsilon_N:=\sup_{|y|\le R}|E_N(y)|
\ \le\ A\,\|v_0\|\,\frac{c_\chi(1+\log_2N)+R}{N},
\qquad
\varepsilon'_N:=\sup_{|y|\le R}|E'_N(y)|
\ \le\ A\,\|v_0\|\,\frac{c_\chi+R}{N} .
\]
In particular $\varepsilon_N+\varepsilon'_N=O(N^{-1}\log N)\to0$ as
$N\to\infty$, for each fixed $R$.
\end{theorem}

\begin{proof}
Partition $\R^3$ into the half-open cubes
$Q_k=\frac kN+[0,\frac1N)^3$, $k\in\Z^3$, of side $1/N$, volume
$N^{-3}$ and diameter $\sqrt3/N$. For any $\Phi$ for which both the sum
and the integral converge absolutely,
\begin{equation}\label{app:cellwise}
\mathcal R_N[\Phi]=\sum_{k\in\Z^3}\int_{Q_k}
\big[\Phi(k/N)-\Phi(\xi)\big]\,d\xi ,
\end{equation}
since $|Q_k|=N^{-3}$. Absolute convergence holds for
$\Phi=\Phi_y$ and $\Phi=\Phi'_{y,ij}$: both are supported in
$\{|\xi|\le1\}$ and dominated there by $\|v_0\||\xi|^{-2}$ respectively
$\|v_0\||\xi|^{-1}$, which are integrable on $\R^3$; and the lattice
sums are finite sums, bounded by
$N^{-3}\sum_{0<|k|\le N}\|v_0\|N^{2}|k|^{-2}
=\|v_0\|N^{-1}S_N\le128\|v_0\|$ and by
$N^{-3}\sum_{0<|k|\le N}\|v_0\|N|k|^{-1}
\le\|v_0\|N^{-2}\cdot N\,S_N\le128\|v_0\|$ respectively, using
$|k|^{-1}\le N|k|^{-2}$ for $|k|\le N$ and Lemma~\ref{app:upper}.

\emph{Pointwise bounds.} On $0<|\xi|\le1$ we have
$|F(\xi)|\le\|v_0\||\xi|^{-2}$, and by the product rule applied to
$F=\chi(|\xi|)|\xi|^{-2}P_\xi v_0$,
\[
|\nabla F(\xi)|\ \le\
\underbrace{c_\chi\|v_0\||\xi|^{-2}}_{\nabla\chi(|\xi|)}
+\underbrace{2\|v_0\||\xi|^{-3}}_{\chi\,\nabla(|\xi|^{-2})}
+\underbrace{4\|v_0\||\xi|^{-3}}_{\chi|\xi|^{-2}\nabla(P_\xi v_0)}
\ \le\ 7c_\chi\|v_0\|\,|\xi|^{-3},
\]
where we used $c_\chi\ge1$, $|\xi|\le1$, $|\nabla(|\xi|^{-2})|
=2|\xi|^{-3}$, and, from
$\partial_a\big(\xi_i\xi_j/|\xi|^2\big)
=(\delta_{ai}\xi_j+\delta_{aj}\xi_i)/|\xi|^2
-2\xi_i\xi_j\xi_a/|\xi|^4$, the honest bound
$|\nabla(P_\xi v_0)|\le4\|v_0\|/|\xi|$. Consequently, for $|y|\le R$,
\begin{equation}\label{app:phibounds}
|\Phi_y|\le\frac{\|v_0\|}{|\xi|^{2}},\quad
|\nabla\Phi_y|\le\frac{7c_\chi\|v_0\|}{|\xi|^{3}}+\frac{R\|v_0\|}{|\xi|^{2}},
\qquad
|\Phi'_y|\le\frac{\|v_0\|}{|\xi|},\quad
|\nabla\Phi'_y|\le\frac{8c_\chi\|v_0\|}{|\xi|^{2}}+\frac{R\|v_0\|}{|\xi|} .
\end{equation}
(For the last one, $\nabla(i\xi_iF_je^{i\xi y})$ has the three terms
$|F|\le\|v_0\||\xi|^{-2}$, $|\xi||\nabla F|\le7c_\chi\|v_0\||\xi|^{-2}$
and $R|\xi||F|\le R\|v_0\||\xi|^{-1}$.)

\emph{The core.} Let
$\mathcal K_0:=\{k:\ Q_k\cap\{|\xi|\le4\sqrt3/N\}\neq\emptyset\}$. Any
such $Q_k$ lies in $\{|\xi|\le5\sqrt3/N\}$ and has $|k|\le5\sqrt3<9$,
i.e.\ $|k|\le8$. Note that $k=0$ is in the core, and that although
$\Phi_y(0)=0$ by convention while $\int_{Q_0}\Phi_y\neq0$, both are
bounded by the two terms below, so no separate treatment is needed. The
core contributes to \eqref{app:cellwise} at most
\[
\int_{|\xi|\le5\sqrt3/N}|\Phi_y|\,d\xi
+N^{-3}\!\!\sum_{0<|k|\le8}\!|\Phi_y(k/N)|
\ \le\ \|v_0\|\Big(\frac{4\pi\cdot5\sqrt3}{N}
+\frac1N\!\!\sum_{0<|k|\le8}\!\frac1{|k|^{2}}\Big)
\ \le\ \frac{1133\,\|v_0\|}{N},
\]
using $\int_{|\xi|\le a}|\xi|^{-2}d\xi=4\pi a$, then
$|\Phi_y(k/N)|\le\|v_0\|N^2|k|^{-2}$, and finally
$\sum_{0<|k|\le8}|k|^{-2}\le S_8\le1024$ by Lemma~\ref{app:upper}
($4\pi\cdot5\sqrt3\le109$, and $109+1024=1133$).

\emph{Dyadic annuli.} Put $J:=\lfloor\log_2\big(N/(4\sqrt3)\big)\rfloor$,
which is $\ge0$ because $N\ge8R\ge8>4\sqrt3$. Let $k\notin\mathcal K_0$
and set $r_k:=\inf_{\xi\in \overline{Q_k}}|\xi|$, so $r_k>4\sqrt3/N>0$.
If $r_k>1$ then $\Phi_y$ vanishes identically on $\overline{Q_k}$ and
that cell contributes $0$; otherwise choose the unique integer $j(k)$
with $2^{-j(k)-1}<r_k\le2^{-j(k)}$. Then $j(k)\ge0$, and
$2^{-j(k)}\ge r_k>4\sqrt3/N$ forces $2^{j(k)}<N/(4\sqrt3)$, i.e.\
$0\le j(k)\le J$; so the annuli
$A_j:=\{2^{-j-1}<|\xi|\le2^{-j}\}$, $0\le j\le J$, together with the
core, do cover $\operatorname{supp}\Phi_y$.

Fix $j$ and consider the cells with $j(k)=j$. They are pairwise
disjoint. Each satisfies $\overline{Q_k}\subset\{|\xi|\ge2^{-j-1}\}$ and
$\overline{Q_k}\subset\{|\xi|\le r_k+\sqrt3/N\le\tfrac54 2^{-j}\}$,
because $\sqrt3/N\le\tfrac14 2^{-j}$ (which follows from
$2^{-j}\ge2^{-J}\ge4\sqrt3/N$). Hence all of them lie in the fixed
annulus $B_j:=\{2^{-j-1}\le|\xi|\le\tfrac54 2^{-j}\}$, whose volume is
$\frac{4\pi}3\big((\tfrac54)^3-(\tfrac12)^3\big)2^{-3j}\le9\cdot2^{-3j}$;
therefore $\sum_{j(k)=j}|Q_k|\le9\cdot2^{-3j}$.

Since $\overline{Q_k}$ is convex, avoids the origin, and $\Phi_y$ is
$C^1$ there, the mean value inequality gives
$\big|\int_{Q_k}[\Phi_y(k/N)-\Phi_y(\xi)]d\xi\big|
\le\frac{\sqrt3}{N}\sup_{\overline{Q_k}}|\nabla\Phi_y|\cdot|Q_k|$. By
\eqref{app:phibounds} and $|\xi|\ge2^{-j-1}$ on $\overline{Q_k}$,
\[
\sup_{\overline{Q_k}}|\nabla\Phi_y|
\le 7c_\chi\|v_0\|2^{3(j+1)}+R\|v_0\|2^{2(j+1)}
=56c_\chi\|v_0\|8^{j}+4R\|v_0\|4^{j}.
\]
Summing over the cells with $j(k)=j$,
\[
\sum_{j(k)=j}\Big|\int_{Q_k}[\Phi_y(k/N)-\Phi_y]\Big|
\le\frac{\sqrt3}{N}\big(56c_\chi\|v_0\|8^{j}+4R\|v_0\|4^{j}\big)
\cdot9\cdot2^{-3j}
\le\frac{\|v_0\|}{N}\big(873\,c_\chi+63\,R\,2^{-j}\big).
\]
Summing over $0\le j\le J\le\log_2N$ and adding the core bound,
\[
|E_N(y)|\le\frac{\|v_0\|}{N}
\Big(1133+873\,c_\chi(1+\log_2N)+126R\Big)
\le\frac{3000\,\|v_0\|}{N}\big(c_\chi(1+\log_2N)+R\big),
\]
using $c_\chi\ge1$ and $1+\log_2N\ge1$. This is the first display.

\emph{The gradient.} Repeat verbatim with $\Phi'_{y,ij}$. The core
contributes at most
\[
\int_{|\xi|\le5\sqrt3/N}\frac{\|v_0\|}{|\xi|}\,d\xi
+N^{-3}\!\!\sum_{0<|k|\le8}\!\frac{\|v_0\|N}{|k|}
=\|v_0\|\Big(\frac{2\pi\cdot75}{N^{2}}
+\frac1{N^{2}}\!\!\sum_{0<|k|\le8}\!\frac1{|k|}\Big)
\le\frac{1816\,\|v_0\|}{N^{2}}\le\frac{1816\,\|v_0\|}{N},
\]
where $\int_{|\xi|\le a}|\xi|^{-1}d\xi=2\pi a^2$ and, by
Lemma~\ref{app:shell},
$\sum_{0<|k|\le8}|k|^{-1}\le\sum_{j=0}^{2}64\cdot8^{j}2^{-j}=1344$
(and $2\pi\cdot75\le472$, $472+1344=1816$). On the annuli,
$\sup_{\overline{Q_k}}|\nabla\Phi'_y|
\le8c_\chi\|v_0\|2^{2(j+1)}+R\|v_0\|2^{j+1}
=32c_\chi\|v_0\|4^{j}+2R\|v_0\|2^{j}$, so the level-$j$ contribution is
at most
\[
\frac{\sqrt3}{N}\big(32c_\chi\|v_0\|4^{j}+2R\|v_0\|2^{j}\big)
\cdot9\cdot2^{-3j}
\le\frac{\|v_0\|}{N}\big(499\,c_\chi2^{-j}+32\,R\,4^{-j}\big),
\]
and now the geometric series converge \emph{without} a logarithm,
summing over all $j\ge0$ to at most
$\frac{\|v_0\|}{N}(998c_\chi+43R)$. Adding the core,
$|E'_N(y)|\le\frac{\|v_0\|}{N}(1816+998c_\chi+43R)
\le\frac{3000\|v_0\|}{N}(c_\chi+R)$. The absence of a logarithm is
structural: the singularity of $\Phi'_y$ is only $|\xi|^{-1}$, the extra
factor $\xi$ taming the origin. The derivative estimate is thus not a
consequence of the $C^0$ one, but it is the easier of the two.
\end{proof}

\subsection{The concentrated test field and the duality lower bound}
\label{app:duality}

\begin{definition}\label{app:testfield}
Let $\Psi\in C_c^\infty(\R^3;\R^3)$ be \emph{real}, with
$\nabla\cdot\Psi=0$, $\operatorname{supp}\Psi\subset\{|y|\le R\}$ for
some $R\ge1$, and $\Psi\not\equiv0$. For $N>R/\pi$ define on $\T^3$
\[
\psi_N(x):=N^{3/2}\sum_{m\in\Z^3}\Psi\big(N(x+2\pi m)\big).
\]
\end{definition}

\begin{lemma}[properties of $\psi_N$]\label{app:psiprops}
Let $N>R/\pi$. Then $\psi_N$ is a well-defined real $C^\infty$ vector
field on $\T^3$; it is divergence-free and has zero mean; on the
fundamental domain $[-\pi,\pi)^3$ exactly the term $m=0$ is nonzero, so
$\operatorname{supp}\psi_N\subset\{|x|\le R/N\}$ there; and
\[
\|\psi_N\|_{L^2(\T^3)}=(2\pi)^{-3/2}\|\Psi\|_{L^2(\R^3)},
\]
which is independent of $N$.
\end{lemma}

\begin{proof}
Local finiteness: $\Psi(N(x+2\pi m))\neq0$ forces
$|x+2\pi m|\le R/N<\pi$; but for $x\in[-\pi,\pi)^3$ and $m\neq0$,
$|x+2\pi m|\ge2\pi|m|_\infty-\pi\ge\pi$, a contradiction. Hence on a
neighbourhood of any point the sum has at most one nonzero term, so
$\psi_N$ is smooth, real, and $2\pi\Z^3$-periodic by construction.
Divergence-free: locally
$\nabla\cdot\psi_N(x)=N^{5/2}(\nabla\cdot\Psi)(N(x+2\pi m))=0$.
Zero mean: with $d\mu=(2\pi)^{-3}dx$ and the substitution $y=Nx$,
\[
\int_{\T^3}\psi_N\,d\mu
=(2\pi)^{-3}N^{3/2}\int_{\R^3}\Psi(Nx)\,dx
=(2\pi)^{-3}N^{-3/2}\int_{\R^3}\Psi ,
\]
and $\int_{\R^3}\Psi_i=\int_{\R^3}\nabla\cdot(y_i\Psi)=0$ since
$\nabla\cdot\Psi=0$ and $y_i\Psi\in C^\infty_c$. (No curl
representation of $\Psi$ is needed: zero mean is automatic for a
compactly supported divergence-free field.) Norm:
\[
\|\psi_N\|_{L^2(\T^3)}^2
=(2\pi)^{-3}\int_{|x|\le R/N}N^{3}|\Psi(Nx)|^2dx
=(2\pi)^{-3}N^{3}N^{-3}\|\Psi\|_{L^2(\R^3)}^2 ,
\]
which is the stated identity.
\end{proof}

\begin{theorem}[duality identity with controlled error]\label{app:pairing}
Let $\Psi$ be as in Definition~\ref{app:testfield} and put
\[
I_\Psi:=\int_{\R^3}\big(V\cdot\nabla V\big)\cdot\Psi\,dy .
\]
Then for every integer $N\ge\max(8R,32)$,
\[
\big\langle\PP(u_N\cdot\nabla u_N),\psi_N\big\rangle_{L^2(\T^3)}
=(2\pi)^{-3}N^{3/2}\big(I_\Psi+\mathcal E_N\big),
\]
where
\[
|\mathcal E_N|\ \le\ \|\Psi\|_{L^1(\R^3)}
\Big[\varepsilon_N\big(6\pi\|v_0\|+\varepsilon'_N\big)
+4\pi\|v_0\|\,\varepsilon'_N\Big]
\ =\ O\Big(\frac{\log N}{N}\Big)
\]
with $\varepsilon_N,\varepsilon'_N$ as in Theorem~\ref{app:rate}.
\end{theorem}

\begin{proof}
By Lemma~\ref{app:psiprops}, $\psi_N$ is zero-mean and divergence-free,
so $\hat\psi_N(0)=0$ and $k\cdot\hat\psi_N(k)=0$ for all $k$, whence
$\PP\psi_N=\psi_N$. Since $\PP$ is an orthogonal projection on
$L^2(\T^3)$, hence self-adjoint,
\[
\big\langle\PP(u_N\cdot\nabla u_N),\psi_N\big\rangle
=\big\langle u_N\cdot\nabla u_N,\PP\psi_N\big\rangle
=\big\langle u_N\cdot\nabla u_N,\psi_N\big\rangle .
\]
Both fields are real. Because $N\ge8R$ and $R\ge1$ we have
$R/N\le\tfrac18<\pi$, so Lemma~\ref{app:psiprops} applies and
$\psi_N(x)=N^{3/2}\Psi(Nx)$ on the fundamental domain, supported in
$\{|x|\le R/N\}$. Using $d\mu=(2\pi)^{-3}dx$ and substituting $x=y/N$
($dx=N^{-3}dy$),
\[
\big\langle u_N\cdot\nabla u_N,\psi_N\big\rangle
=(2\pi)^{-3}N^{3/2}N^{-3}\int_{|y|\le R}
(u_N\cdot\nabla u_N)(y/N)\cdot\Psi(y)\,dy .
\]
By Lemma~\ref{app:innerexact}, for $|y|\le R$,
\[
(u_N\cdot\nabla u_N)_j(y/N)
=u_{N,i}(y/N)\,\partial_iu_{N,j}(y/N)
=N^{3}\big(V_i(y)+E_{N,i}(y)\big)\big(\partial_iV_j(y)+E'_{N,ij}(y)\big),
\]
so the factor $N^{-3}N^{3}$ cancels and
\[
\big\langle u_N\cdot\nabla u_N,\psi_N\big\rangle
=(2\pi)^{-3}N^{3/2}\big(I_\Psi+\mathcal E_N\big),
\qquad
\mathcal E_N=\int_{|y|\le R}
\big[E_{N,i}\partial_iV_j+V_iE'_{N,ij}+E_{N,i}E'_{N,ij}\big]\Psi_j\,dy ,
\]
where we used $\operatorname{supp}\Psi\subset\{|y|\le R\}$ to write the
main term as $I_\Psi$. Bounding each factor by its supremum and using
Theorem~\ref{app:Vprops}(2),(3) --- $\|V\|_\infty\le4\pi\|v_0\|$,
$\|\nabla V\|_\infty\le6\pi\|v_0\|$ --- gives the stated estimate;
$\varepsilon_N,\varepsilon'_N=O(N^{-1}\log N)$ by
Theorem~\ref{app:rate}, $R$ and $\Psi$ being fixed.
\end{proof}

\begin{proposition}[the $N^3$ lower bound, given a good test field]
\label{app:N3}
Suppose some $\Psi$ as in Definition~\ref{app:testfield} has
$I_\Psi\neq0$. Then
\[
\big\|\PP(u_N\cdot\nabla u_N)\big\|_{L^2(\T^3)}^2
\ \ge\ \frac{|I_\Psi+\mathcal E_N|^2}
{(2\pi)^{3}\,\|\Psi\|_{L^2(\R^3)}^{2}}\;N^{3}
\qquad(N\ge\max(8R,32)),
\]
and consequently
$\|\PP(u_N\cdot\nabla u_N)\|_{2}^2\ge c_0N^3$ for all $N\ge N_*$, where
\[
c_0=\frac{|I_\Psi|^2}{2\,(2\pi)^{3}\,\|\Psi\|_{L^2(\R^3)}^{2}},
\qquad
N_*=\min\Big\{N_0\ge\max(8R,32):\ |\mathcal E_N|\le(1-2^{-1/2})|I_\Psi|
\ \ \forall N\ge N_0\Big\},
\]
which is finite because $\mathcal E_N\to0$.
\end{proposition}

\begin{proof}
By Cauchy--Schwarz in $L^2(\T^3)$,
\[
\big\|\PP(u_N\cdot\nabla u_N)\big\|_{2}
\ \ge\ \frac{\big|\langle\PP(u_N\cdot\nabla u_N),\psi_N\rangle\big|}
{\|\psi_N\|_{2}} .
\]
Insert Theorem~\ref{app:pairing} and Lemma~\ref{app:psiprops}: the
right-hand side equals
\[
\frac{(2\pi)^{-3}N^{3/2}|I_\Psi+\mathcal E_N|}
{(2\pi)^{-3/2}\|\Psi\|_{L^2(\R^3)}}
=(2\pi)^{-3/2}N^{3/2}\frac{|I_\Psi+\mathcal E_N|}{\|\Psi\|_{L^2(\R^3)}} .
\]
Squaring gives the display. For $N\ge N_*$ we have
$|I_\Psi+\mathcal E_N|\ge2^{-1/2}|I_\Psi|$, whence
$|I_\Psi+\mathcal E_N|^2\ge\tfrac12|I_\Psi|^2$ and the constant $c_0$ as
stated.
\end{proof}

Everything is now reduced to the single question of whether some
admissible $\Psi$ has $I_\Psi\neq0$.

\subsection{Non-degeneracy of the continuum nonlinearity}
\label{app:nondeg}

\subsubsection*{Reduction to a symbol statement}

\begin{lemma}[symbol of the nonlinearity]\label{app:symbol}
$V\cdot\nabla V=\nabla\cdot(V\otimes V)$ has the representation
\[
(V\cdot\nabla V)(y)=\int_{\R^3}M(\zeta)e^{i\zeta\cdot y}\,d\zeta,
\qquad
M_i(\zeta)=i\,\zeta_j\,T_{ij}(\zeta),\qquad T_{ij}:=F_i*F_j ,
\]
where $T$ is real, symmetric, supported in $\{|\zeta|\le2\}$, and
$T,M\in L^1(\R^3)\cap L^2(\R^3)$. Consequently
$\PP(V\cdot\nabla V)\in L^2(\R^3)$ with symbol $P_\zeta M(\zeta)$, and
\[
\PP(V\cdot\nabla V)\equiv0
\quad\Longleftrightarrow\quad
\zeta\times\big(T(\zeta)\zeta\big)=0\ \text{ for a.e. }\zeta .
\]
\end{lemma}

\begin{proof}
$V=\widetilde F$ with $F\in L^1$, and $\partial_jV_i=\widetilde{i\xi_jF_i}$
with $\xi F\in L^1$ (Theorem~\ref{app:Vprops}(1),(3)). The pointwise
product of inverse transforms of $L^1$ functions is the inverse
transform of the convolution, so
$V_j\partial_jV_i=\widetilde{F_j*(i\,\cdot_j F_i)}$, i.e.\ the symbol of
$V\cdot\nabla V$ is
$i\int F_j(\zeta-\eta)\,\eta_j\,F_i(\eta)\,d\eta$. Since
$(\zeta-\eta)\cdot F(\zeta-\eta)=0$ identically (Definition~\ref{app:Vdef}),
we may replace $\eta_jF_j(\zeta-\eta)$ by $\zeta_jF_j(\zeta-\eta)$
inside the integral, obtaining
$M_i(\zeta)=i\zeta_j(F_i*F_j)(\zeta)=i\zeta_jT_{ij}(\zeta)$.

$T$ is real because $F$ is real, symmetric because convolution is
commutative, and supported in $\{|\zeta|\le2\}$ because
$\operatorname{supp}F\subset\{|\xi|\le1\}$. Integrability: by
Theorem~\ref{app:Vprops}(1), $F\in L^{4/3}$; Young's inequality with
$\frac1r=\frac2{4/3}-1=\frac12$ gives $T=F*F\in L^{2}$, and since $T$
has compact support, $T\in L^1\cap L^2$; the same holds for
$M=i\zeta T$, as $|\zeta|\le2$ on the support. By Plancherel,
$V\cdot\nabla V\in L^2(\R^3)$ and, since $|P_\zeta M|\le|M|$,
$\PP(V\cdot\nabla V)\in L^2(\R^3)$ with symbol $P_\zeta M$. Finally
$\PP(V\cdot\nabla V)=0$ in $L^2$ iff $P_\zeta M(\zeta)=0$ for a.e.
$\zeta$, and for a vector $w$ one has $P_\zeta w=0$ iff $w\parallel\zeta$
iff $\zeta\times w=0$; here $w=i\,T(\zeta)\zeta$.
\end{proof}

\begin{remark}[classical input: density of smooth compactly supported
solenoidal fields]\label{app:classical-density}
The space $C^\infty_{c,\sigma}(\R^3;\R^3)$ of real, smooth, compactly
supported, divergence-free vector fields is dense in
$L^2_\sigma(\R^3;\R^3)$, the closed subspace of real $L^2$ fields that
are distributionally divergence-free; indeed $L^2_\sigma$ is by the
standard definition the $L^2$-closure of $C^\infty_{c,\sigma}$, and on
$\R^3$ the two descriptions agree. See, e.g.,
\cite{Temam,LadyzhenskayaBook,ConstantinFoias}.
\end{remark}

\begin{remark}[classical input: Helmholtz--Leray decomposition on $\R^3$]
\label{app:classical-leray}
For $f\in L^2(\R^3;\R^3)$ one has $f-\PP f=\nabla p$ for some
$p\in L^2_{\rm loc}(\R^3)$ with $\nabla p\in L^2$; equivalently, the
symbol of $f-\PP f$ is $\zeta(\zeta\cdot\hat f)/|\zeta|^2$. Same
references.
\end{remark}

\begin{lemma}[from non-degeneracy to a test field]\label{app:testexists}
If $\PP(V\cdot\nabla V)\not\equiv0$, then there exists a real
$\Psi\in C^\infty_c(\R^3;\R^3)$ with $\nabla\cdot\Psi=0$ and
$I_\Psi=\int_{\R^3}(V\cdot\nabla V)\cdot\Psi\neq0$; after enlarging $R$
if necessary, $\Psi$ satisfies Definition~\ref{app:testfield}.
\end{lemma}

\begin{proof}
By Lemma~\ref{app:symbol}, $\PP(V\cdot\nabla V)$ lies in $L^2(\R^3)$,
is real (its symbol $P_\zeta M$ satisfies the reality condition because
$F$ is real and even), is distributionally divergence-free (its symbol
is orthogonal to $\zeta$), and is nonzero by hypothesis. Hence it is a
nonzero element of the real Hilbert space $L^2_\sigma(\R^3)$. By
Remark~\ref{app:classical-density} there is a real
$\Psi\in C^\infty_{c,\sigma}(\R^3)$ with
$\langle\PP(V\cdot\nabla V),\Psi\rangle_{L^2(\R^3)}\neq0$. By
Remark~\ref{app:classical-leray},
$V\cdot\nabla V-\PP(V\cdot\nabla V)=\nabla p$ with $\nabla p\in L^2$, and
$\int\nabla p\cdot\Psi=-\int p\,\nabla\cdot\Psi=0$ because $\Psi$ is
compactly supported and divergence-free. Therefore
$I_\Psi=\langle\PP(V\cdot\nabla V),\Psi\rangle\neq0$. Choosing
$R\ge\max(1,\sup\{|y|:y\in\operatorname{supp}\Psi\})$ puts $\Psi$ in
Definition~\ref{app:testfield}.
\end{proof}

\subsubsection*{The $\chi$-independent leading singularity}

By rotational equivariance we may and do assume from now on that
$v_0=\|v_0\|e_3$: $F$ and $V$ are linear in $v_0$ and rotate with it,
$T$ and $\tau$ are quadratic, and the conclusion
$\PP(V\cdot\nabla V)\not\equiv0$ is rotation invariant.

\begin{definition}\label{app:tau}
$\tau_{ij}:=h_i*h_j$, with $h(\xi)=P_\xi v_0/|\xi|^2$ the uncut symbol
of Definition~\ref{app:Vdef}.
\end{definition}

\begin{lemma}[$\tau$ is well defined and homogeneous]\label{app:tauwd}
The integral $\tau_{ij}(\zeta)=\int_{\R^3}h_i(\eta)h_j(\zeta-\eta)\,d\eta$
converges absolutely for every $\zeta\neq0$, and $\tau$ is homogeneous of
degree $-1$ and smooth on $\R^3\setminus\{0\}$.
\end{lemma}

\begin{proof}
$|h(\xi)|\le\|v_0\||\xi|^{-2}$. Near $\eta=0$ the integrand is
$O(|\eta|^{-2})$ and near $\eta=\zeta$ it is $O(|\zeta-\eta|^{-2})$, both
locally integrable on $\R^3$; for $|\eta|\ge2|\zeta|$ it is
$O(|\eta|^{-4})$, integrable at infinity. Substituting
$\eta=\lambda\eta'$ and using $h(\lambda\eta)=\lambda^{-2}h(\eta)$ gives
$\tau(\lambda\zeta)=\lambda^{-1}\tau(\zeta)$ for $\lambda>0$. Smoothness
away from $0$ follows from differentiating under the integral sign on
the region where the two singularities are separated.
\end{proof}

\begin{theorem}[closed form of the leading obstruction]
\label{app:tauformula}
With $v_0=\|v_0\|e_3$, for every $\zeta\neq0$,
\[
\zeta\times\big(\tau(\zeta)\zeta\big)
=\frac{3\pi^3}{8}\,\|v_0\|^2\,\frac{\zeta_3}{|\zeta|}\,
\big(\zeta\times e_3\big).
\]
In particular $\zeta\times(\tau(\zeta)\zeta)\neq0$ whenever
$\zeta_3\neq0$ and $\zeta\not\parallel e_3$.
\end{theorem}

\begin{proof}
Both sides are quadratic in $\|v_0\|$, so we normalise $\|v_0\|=1$,
$v_0=e_3$. By Theorem~\ref{app:farfield}, $\widetilde h=\pi^2W$ with
\[
W_i(y)=\frac{\delta_{i3}}{r}+\frac{y_3y_i}{r^3},\qquad r=|y| .
\]
$W_i$ is locally integrable, homogeneous of degree $-1$ and smooth away
from the origin, so Remark~\ref{app:classical-homog} applies with
$a=b=1$ and gives
\[
\tau_{ij}=h_i*h_j=(2\pi)^{-3}\,\mathcal F\big[\widetilde h_i\widetilde h_j\big]
=(2\pi)^{-3}\pi^4\,\mathcal F\big[W_iW_j\big],
\]
where we used $\mathcal F\widetilde h_i=(2\pi)^3h_i$ and the evenness of
$h$. (Both sides are homogeneous of degree $-1$ and locally integrable,
so no distribution supported at the origin can be added.) Expanding,
\[
W_iW_j=\frac{\delta_{i3}\delta_{j3}}{r^{2}}
+\frac{\delta_{i3}y_3y_j+\delta_{j3}y_3y_i}{r^{4}}
+\frac{y_3^{2}y_iy_j}{r^{6}} .
\]
Each bracket is locally integrable and homogeneous of degree $-2$, so by
Remark~\ref{app:classical-riesz} its transform is an honest homogeneous
distribution of degree $-1$, given by
\[
\mathcal F\Big[\frac1{r^{2}}\Big]=\frac{2\pi^{2}}{s},
\qquad
\mathcal F\Big[\frac{y_iy_j}{r^{4}}\Big]
=\pi^{2}\Big(\frac{\delta_{ij}}{s}-\frac{\zeta_i\zeta_j}{s^{3}}\Big),
\qquad
\mathcal F\Big[\frac{y_3^{2}y_iy_j}{r^{6}}\Big]
=\frac{\pi^{2}}{12}\,\partial_3^{2}\partial_i\partial_j\,s^{3},
\]
where $s:=|\zeta|$. (Indeed $\mathcal F[r^{-2}]=c_{3,2}s^{-1}$;
$\mathcal F[y_iy_jr^{-4}]=-\partial_i\partial_j(c_{3,4}s)
=\pi^2\partial_i\partial_js=\pi^2(\delta_{ij}/s-\zeta_i\zeta_j/s^3)$;
$\mathcal F[y_3^2y_iy_jr^{-6}]
=\partial_3^2\partial_i\partial_j(c_{3,6}s^3)$. Trace check on the
middle formula: $\pi^2(3/s-1/s)=2\pi^2/s=\mathcal F[r^{-2}]$, as it must
be.) Direct differentiation of $s^3$ gives
\[
\partial_3^{2}\partial_i\partial_j s^{3}
=3\Big[\frac{\delta_{ij}+2\delta_{i3}\delta_{j3}}{s}
-\frac{\zeta_3^{2}\delta_{ij}
+2\zeta_3(\delta_{i3}\zeta_j+\delta_{j3}\zeta_i)+\zeta_i\zeta_j}{s^{3}}
+\frac{3\zeta_i\zeta_j\zeta_3^{2}}{s^{5}}\Big].
\]
Assembling the three pieces (the middle one contributing
$\pi^2[2\delta_{i3}\delta_{j3}/s-\zeta_3(\delta_{i3}\zeta_j
+\delta_{j3}\zeta_i)/s^3]$) yields
$\mathcal F[W_iW_j]=\pi^{2}G_{ij}$ with
\[
G_{ij}=\frac{\delta_{ij}}{4s}
+\frac{9\,\delta_{i3}\delta_{j3}}{2s}
-\frac{\zeta_3^{2}\delta_{ij}}{4s^{3}}
-\frac{3\zeta_3(\delta_{i3}\zeta_j+\delta_{j3}\zeta_i)}{2s^{3}}
-\frac{\zeta_i\zeta_j}{4s^{3}}
+\frac{3\zeta_i\zeta_j\zeta_3^{2}}{4s^{5}} ,
\]
so that $\tau_{ij}=(2\pi)^{-3}\pi^{4}\cdot\pi^{2}G_{ij}
=\frac{\pi^{3}}{8}G_{ij}$.

Contract with $\zeta_j$, using $\zeta_j\zeta_j=s^2$. The terms
$\frac{\delta_{ij}}{4s}\zeta_j=\frac{\zeta_i}{4s}$ and
$-\frac{\zeta_i\zeta_j}{4s^3}\zeta_j=-\frac{\zeta_i}{4s}$ cancel. The
$\delta_{i3}\zeta_3/s$ coefficients are $\frac92$ from the second term
and $-\frac32$ from the fourth, total $3$. The $\zeta_i\zeta_3^2/s^3$
coefficients are $-\frac14$ (third term), $-\frac32$ (fourth term) and
$+\frac34$ (sixth term), total $-1$. Hence
\[
G(\zeta)\zeta=\frac{3\zeta_3}{s}\,e_3-\frac{\zeta_3^{2}}{s^{3}}\,\zeta ,
\qquad\text{so}\qquad
\zeta\times\big(G(\zeta)\zeta\big)=\frac{3\zeta_3}{s}\,(\zeta\times e_3),
\]
the second term being parallel to $\zeta$. Multiplying by
$\frac{\pi^3}{8}$ and restoring $\|v_0\|^2$ gives the claim. Finally
$\zeta\times e_3=0$ exactly when $\zeta\parallel e_3$.
\end{proof}

\begin{remark}[physical-space cross-check, free of any regularisation]
\label{app:crosscheck}
Theorem~\ref{app:tauformula} passes through the finite-part
continuations of Remark~\ref{app:classical-riesz}. The following
elementary computation confirms it independently, with no
regularisation and no Fourier transform, and is worth recording because
it identifies what the obstruction \emph{is}: the statement that the
Oseen field is not a stationary Euler flow.

For $W(y)=\frac{e_3}{r}+\frac{y_3y}{r^{3}}$ on $\R^3\setminus\{0\}$,
\[
\nabla\cdot W=0,\qquad
\omega:=\nabla\times W=\frac{2\,(e_3\times y)}{r^{3}},\qquad
\nabla\times\big(W\times\omega\big)
=\frac{12\,y_3\,(e_3\times y)}{r^{6}}\ \not\equiv0 .
\]
\emph{Proof.} $\nabla\cdot(e_3/r)=e_3\cdot\nabla r^{-1}=-y_3/r^3$, and
$\partial_i(y_3y_ir^{-3})=y_3r^{-3}+3y_3r^{-3}-3y_3r^{-3}=y_3/r^3$, so
$\nabla\cdot W=0$. Next
$\nabla\times(r^{-1}e_3)=\nabla(r^{-1})\times e_3=-r^{-3}y\times e_3
=r^{-3}(e_3\times y)$, while
$[\nabla\times(y_3y\,r^{-3})]_a
=\epsilon_{abi}\partial_b(y_3y_ir^{-3})
=\epsilon_{abi}(\delta_{b3}y_i+y_3\delta_{bi})r^{-3}
-3\epsilon_{abi}y_3y_iy_br^{-5}
=\epsilon_{a3i}y_ir^{-3}=(e_3\times y)_ar^{-3}$, the last two terms
vanishing by antisymmetry. Adding gives $\omega$. Then, using
$e_3\times(e_3\times y)=y_3e_3-y$ and $y\times(e_3\times y)=r^2e_3-y_3y$,
\[
W\times\omega
=\frac{2(y_3e_3-y)}{r^{4}}+\frac{2y_3e_3}{r^{4}}-\frac{2y_3^{2}y}{r^{6}}
=\frac{4y_3e_3}{r^{4}}-\frac{2y}{r^{4}}-\frac{2y_3^{2}y}{r^{6}} .
\]
Now curl each piece. For any radial $f$, $\nabla\times(f(r)y)
=\nabla f\times y=f'(r)\hat y\times y=0$, so $-2y/r^4$ contributes
nothing. With $g=-2y_3^{2}/r^{6}$ one has
$\nabla g=-4y_3e_3/r^{6}+12y_3^{2}y/r^{8}$, so
$\nabla\times(gy)=\nabla g\times y=-4y_3(e_3\times y)/r^{6}$. With
$f=4y_3/r^{4}$ one has $\nabla f=4e_3/r^{4}-16y_3y/r^{6}$, so
$\nabla\times(fe_3)=\nabla f\times e_3
=-16y_3(y\times e_3)/r^{6}=16y_3(e_3\times y)/r^{6}$. The total is
$12y_3(e_3\times y)/r^{6}$. Since
$W\cdot\nabla W=\nabla\frac{|W|^2}{2}-W\times\omega$, it follows that
$\nabla\times(W\cdot\nabla W)\not\equiv0$, i.e.\ $W\cdot\nabla W$ is not
a gradient, i.e.\ $\PP(W\cdot\nabla W)\neq0$. $\square$

\emph{Consistency with Theorem~\ref{app:tauformula}.} The curl of
$\PP(\widetilde h\cdot\nabla\widetilde h)$ has symbol
$i\zeta\times\big(iP_\zeta\tau\zeta\big)=-\zeta\times(\tau\zeta)$, using
$\zeta\times P_\zeta A=\zeta\times A$. With $\|v_0\|=1$,
Theorem~\ref{app:tauformula} makes this
$-\frac{3\pi^{3}}{8}\zeta_3(\zeta\times e_3)/|\zeta|$. Writing its
$a$-component as $-\frac{3\pi^3}{8}\epsilon_{ab3}\zeta_b\zeta_3/|\zeta|$
and using $\widetilde{|\zeta|^{-1}}=c_{3,1}r^{-2}=4\pi r^{-2}$ together
with $\widetilde{\zeta_b\zeta_3G}=-\partial_b\partial_3\widetilde G$,
\[
\widetilde{\Big[\frac{\zeta_b\zeta_3}{|\zeta|}\Big]}
=-4\pi\,\partial_b\partial_3r^{-2}
=\frac{8\pi\delta_{b3}}{r^{4}}-\frac{32\pi y_3y_b}{r^{6}} ;
\]
since $\epsilon_{ab3}\delta_{b3}=0$, only the second term survives,
leaving $-32\pi y_3(y\times e_3)_a/r^{6}
=+32\pi y_3(e_3\times y)_a/r^{6}$. Multiplying by
$-\frac{3\pi^{3}}{8}$ gives $-12\pi^{4}y_3(e_3\times y)/r^{6}$. On the
other hand $\widetilde h=\pi^2W$ gives
$\widetilde h\cdot\nabla\widetilde h=\pi^{4}W\cdot\nabla W$ and hence
$\nabla\times(\widetilde h\cdot\nabla\widetilde h)
=-\pi^{4}\nabla\times(W\times\omega)
=-12\pi^{4}y_3(e_3\times y)/r^{6}$. The two derivations agree exactly,
constants included.
\end{remark}

\subsubsection*{Transfer to the cutoff family}

\begin{lemma}[the cutoff perturbation is bounded near $\zeta=0$]
\label{app:perturb}
For all $|\zeta|\le\frac18$ and all $i,j$,
\[
\big|T_{ij}(\zeta)-\tau_{ij}(\zeta)\big|\ \le\ C_1:=64\pi\|v_0\|^{2},
\]
uniformly over all admissible $\chi$.
\end{lemma}

\begin{proof}
Since $F=\chi(|\cdot|)h$,
\[
T_{ij}(\zeta)-\tau_{ij}(\zeta)
=\int_{\R^3}h_i(\eta)h_j(\zeta-\eta)
\big[\chi(|\eta|)\chi(|\zeta-\eta|)-1\big]\,d\eta .
\]
The bracket vanishes whenever $|\eta|\le\frac12$ \emph{and}
$|\zeta-\eta|\le\frac12$, because $\chi\equiv1$ on $[0,\frac12]$. For
$|\zeta|\le\frac18$, the condition $|\eta|\le\frac14$ implies
$|\zeta-\eta|\le\frac18+\frac14=\frac38<\frac12$; hence the integrand is
supported in $\{|\eta|\ge\frac14\}$. There
$|\zeta-\eta|\ge|\eta|-\frac18\ge|\eta|-\frac{|\eta|}{2}
=\frac{|\eta|}{2}$. With $|h|\le\|v_0\||\cdot|^{-2}$ and
$|\chi\chi-1|\le1$,
\[
|T_{ij}-\tau_{ij}|
\le\|v_0\|^{2}\int_{|\eta|\ge1/4}\frac{d\eta}{|\eta|^{2}|\zeta-\eta|^{2}}
\le4\|v_0\|^{2}\int_{|\eta|\ge1/4}\frac{d\eta}{|\eta|^{4}}
=4\|v_0\|^{2}\cdot4\pi\int_{1/4}^{\infty}\frac{dr}{r^{2}}
=64\pi\|v_0\|^{2},
\]
which is $C_1$.
\end{proof}

\begin{theorem}[non-degeneracy]\label{app:vnondeg}
For every $v_0\in\R^3\setminus\{0\}$ and every admissible $\chi$,
\[
\PP_{\R^3}\big(V\cdot\nabla V\big)\ \not\equiv\ 0 .
\]
Explicitly, with $v_0=\|v_0\|e_3$, one has
$\zeta\times(T(\zeta)\zeta)\neq0$ --- hence $P_\zeta M(\zeta)\neq0$ ---
at every point of the open set of positive measure
\[
\mathcal U:=\Big\{\zeta\in\R^3:\ 0<|\zeta|<\tfrac{\pi^{2}}{2048},\
|\zeta_3|\ge\tfrac{|\zeta|}{2},\
|\zeta\times e_3|\ge\tfrac{|\zeta|}{2}\Big\}.
\]
\end{theorem}

\begin{proof}
Let $\zeta\in\mathcal U$. By Theorem~\ref{app:tauformula},
\[
\big|\zeta\times(\tau(\zeta)\zeta)\big|
=\frac{3\pi^{3}}{8}\|v_0\|^{2}\,
\frac{|\zeta_3|\,|\zeta\times e_3|}{|\zeta|}
\ \ge\ \frac{3\pi^{3}}{8}\|v_0\|^{2}\cdot\frac{|\zeta|}{4}
=\frac{3\pi^{3}}{32}\|v_0\|^{2}|\zeta| .
\]
Since $\frac{\pi^{2}}{2048}\approx4.8\times10^{-3}<\frac18$,
Lemma~\ref{app:perturb} applies, and the Frobenius norm of
$T(\zeta)-\tau(\zeta)$ is at most $3C_1$ (nine entries each at most
$C_1$). Hence
\[
\big|\zeta\times\big((T(\zeta)-\tau(\zeta))\zeta\big)\big|
\le|\zeta|\cdot3C_1|\zeta|=192\pi\|v_0\|^{2}|\zeta|^{2},
\]
and therefore
\[
\big|\zeta\times(T(\zeta)\zeta)\big|
\ \ge\ \|v_0\|^{2}|\zeta|
\Big(\frac{3\pi^{3}}{32}-192\pi|\zeta|\Big)\ >\ 0
\qquad\text{whenever }
|\zeta|<\frac{3\pi^{3}}{32\cdot192\pi}=\frac{3\pi^{2}}{6144}
=\frac{\pi^{2}}{2048},
\]
which holds on $\mathcal U$. The set $\mathcal U$ is open and nonempty
(it contains, e.g., all small $\zeta$ in the direction
$(1,0,1)/\sqrt2$), hence of positive measure. By
Lemma~\ref{app:symbol}, $P_\zeta M\neq0$ on a set of positive measure,
so $\PP(V\cdot\nabla V)\neq0$ in $L^2(\R^3)$.
\end{proof}

\begin{remark}[uniformity, and where the obstruction lives]
\label{app:uniformity}
The proof of Theorem~\ref{app:vnondeg} is uniform in $\chi$ and in
$v_0$: the factor $\|v_0\|^2$ cancels between the main term and the
perturbation, and the only properties of $\chi$ used are $0\le\chi\le1$
and $\chi\equiv1$ on $[0,\frac12]$ --- no smoothness of $\chi$ enters
\S\ref{app:nondeg} at all. Neither a specific $\chi$ nor any numerically
evaluated convolution integral is required: the obstruction is carried
entirely by the $\chi$-independent leading singularity of $F*F$ at
$\zeta=0$, which is the same for every admissible cutoff. Moreover
Theorem~\ref{app:tauformula} shows the obstruction degenerates exactly
on the two directions $\zeta\parallel v_0$ and $\zeta\perp v_0$, so the
two cone conditions in $\mathcal U$ are essentially optimal.

This uniformity does \textbf{not} propagate to the constants of
Theorem~\ref{app:main}. It is the \emph{qualitative} conclusion
$\PP(V\cdot\nabla V)\not\equiv0$ that holds uniformly in $(\chi,v_0)$,
and that is what licenses the quantifier ``for every admissible $\chi$
and every $v_0\neq0$'' there. The quantitative constants $c_0$ and
$N_*$ are extracted only after a test field $\Psi$ has been chosen for
the particular profile $V=V^{\chi,v_0}$
(Lemma~\ref{app:testexists}), and they depend on $(\chi,v_0)$; no
$(\chi,v_0)$-independent value of either is claimed, and neither is
effective (Remark~\ref{app:noneffective}).
\end{remark}

\subsection{Proof of Theorem~\ref{capacity}}
\label{app:final}

\begin{theorem}[capacity bound; restatement of Theorem~\ref{capacity}]
\label{app:main}
Let $v_0\in\R^3\setminus\{0\}$, let $\chi$ be admissible
(Definition~\ref{app:admissible}), and let $u_N=u_N^{\chi,v_0}$ be the
family \eqref{app:famdef}. Then:
\begin{enumerate}
\item[(1)] $u_N$ is a real, zero-mean, exactly divergence-free
trigonometric polynomial banded at $|k|\le N$, and the exact laws of
Theorem~\ref{app:exactlaws} hold;
\item[(2)] $\dfrac{N}{44544}\le N_0^2(u_N)\le\dfrac{64}{3}N$ for every
$N\ge32$;
\item[(3)] there exist $c_0=c_0(\chi,v_0)>0$ and
$N_*=N_*(\chi,v_0)<\infty$ such that
\[
\big\|\PP(u_N\cdot\nabla u_N)\big\|_{L^2(\T^3)}^{2}\ \ge\ c_0\,N^{3}
\qquad\text{for all }N\ge N_* .
\]
\end{enumerate}
\end{theorem}

\begin{proof}
(1) is Lemma~\ref{app:structure} together with
Theorem~\ref{app:exactlaws}. (2) is Proposition~\ref{app:twosided}.

(3) By Theorem~\ref{app:vnondeg}, $\PP_{\R^3}(V\cdot\nabla V)\not\equiv0$,
where $V$ is the profile of Definition~\ref{app:Vdef} attached to the
given $(\chi,v_0)$. By Lemma~\ref{app:testexists} there is a real
$\Psi\in C^\infty_c(\R^3;\R^3)$, divergence-free, supported in
$\{|y|\le R\}$ for some $R\ge1$, with $I_\Psi\neq0$. Fix such a $\Psi$.
Lemma~\ref{app:innerexact} and Theorem~\ref{app:rate} give the exact
inner rescaling with $\varepsilon_N+\varepsilon'_N=O(N^{-1}\log N)$;
Lemma~\ref{app:psiprops} and Theorem~\ref{app:pairing} convert the
pairing of $\PP(u_N\cdot\nabla u_N)$ against the concentrated field
$\psi_N$ into $(2\pi)^{-3}N^{3/2}(I_\Psi+\mathcal E_N)$ with
$\mathcal E_N\to0$; and Proposition~\ref{app:N3} then yields the stated
bound with
\[
c_0=c_0(\chi,v_0)=\frac{|I_\Psi|^{2}}
{2\,(2\pi)^{3}\,\|\Psi\|_{L^2(\R^3)}^{2}},
\qquad
N_*=N_*(\chi,v_0)
=\min\Big\{N_0\ge\max(8R,32):|\mathcal E_N|\le(1-2^{-1/2})|I_\Psi|
\ \forall N\ge N_0\Big\}.
\]
Both depend on $(\chi,v_0)$ alone, through $V=V^{\chi,v_0}$ and through
the test field $\Psi$ selected for it.
\end{proof}

\begin{remark}[what the constants depend on, and that they are
non-effective]\label{app:noneffective}
The constants $c_0$ and $N_*$ produced above depend on the pair
$(\chi,v_0)$ only, through the profile $V=V^{\chi,v_0}$ and through the
test field $\Psi$ selected in Lemma~\ref{app:testexists}. They are
\textbf{not effective}. The reason is localised precisely at
Lemma~\ref{app:testexists}: $\Psi$ is obtained from the \emph{density}
of $C^\infty_{c,\sigma}(\R^3)$ in $L^2_\sigma(\R^3)$
(Remark~\ref{app:classical-density}), a pure existence argument, so
neither $\Psi$, nor its support radius $R$, nor $I_\Psi$, nor
$\|\Psi\|_{L^2}$, nor $c_0$, nor $N_*$ is exhibited. Every other
constant in this appendix is explicit, namely those of
Lemmas~\ref{app:shell}--\ref{app:lower},
Proposition~\ref{app:twosided}, Theorems~\ref{app:exactlaws},
\ref{app:Vprops}, \ref{app:farfield}, \ref{app:rate},
\ref{app:pairing}, \ref{app:tauformula}, \ref{app:vnondeg},
Lemmas~\ref{app:innerexact}, \ref{app:psiprops}, \ref{app:symbol},
\ref{app:tauwd}, \ref{app:perturb}, and Proposition~\ref{app:N3}
\emph{as a function of} $(\Psi,R,I_\Psi)$.

An effective variant is available in principle --- take $\Psi$ with
$\widehat\Psi$ a smooth bump supported in the explicit set $\mathcal U$
of Theorem~\ref{app:vnondeg}, mollified to compact support in $y$ ---
but it is expensive: non-vanishing is certified only for
$|\zeta|<\pi^2/2048\approx4.8\times10^{-3}$, so such a $\Psi$ has
Fourier support at that scale, hence spatial radius $R\gtrsim10^{3}$,
hence $N_*\ge8R\gtrsim10^{4}$, with a correspondingly small $c_0$. This
is the sole source of non-effectivity in the paper.
\end{remark}

\begin{remark}[the exponent $3$ is sharp, and the sharp cutoff is not
covered]\label{app:sharpness}
The exponent cannot be improved: by Theorem~\ref{main}(d) and
Theorem~\ref{app:exactlaws},
$\|\PP(u_N\cdot\nabla u_N)\|_2^2\le\|u_N\|_\infty^2H_1(u_N)$, and both
$\|u_N\|_\infty$ and $H_1(u_N)$ are $O(N)$ by
Proposition~\ref{app:twosided} and
$\|u_N\|_\infty\le\sum_{k\neq0}|\hat u_N(k)|\le\|v_0\|\Sigma_N^\chi
\le\|v_0\|S_N\le128\|v_0\|N$; so the left-hand side is $O(N^3)$.

Smoothness of $\chi$ is used in exactly one place, namely
Theorem~\ref{app:rate}, through $|\nabla F|\lesssim c_\chi|\xi|^{-3}$.
For the sharply truncated family $\chi=\mathbf 1_{[0,1]}$ this bound
fails, $F$ jumping across $|\xi|=1$; the Riemann-sum defect then carries
a lattice-point (Gauss-circle) term, and $V$ loses the rapidly decaying
tail correction of Theorem~\ref{app:farfield}, acquiring instead an
oscillatory $|y|^{-2}$ tail. Nothing in \S\ref{app:duality} or
\S\ref{app:nondeg} would change, both being insensitive to $\chi$; the
sharp-cutoff case is thus a technical gap in the inner-rescaling step,
not a structural one. It is left open because no statement in this paper
requires it (see Remark~\ref{lstarremark}).
\end{remark}
  inside the paper's \appendix.

\section{Complete proof of the capacity bound (Theorem~\ref{capacity})}
\label{app:capacity}

This appendix proves Theorem~\ref{capacity} in full. Nothing below is
quoted from elsewhere except four classical facts, each isolated and
labelled as such (Remarks~\ref{app:classical-riesz},
\ref{app:classical-homog}, \ref{app:classical-density} and
\ref{app:classical-leray}); no numerical statement is used anywhere.

The chain is: the family $u_N$ has an exact scaling
$\hat u_N(k)=N^{-2}F(k/N)$ (\S\ref{app:profile}), so that the lattice
sums reconstructing $u_N$ near the origin are \emph{literally} Riemann
sums, of spacing $1/N$, for the inverse Fourier transform $V$ of $F$
(\S\ref{app:riemann}); pairing $\PP(u_N\cdot\nabla u_N)$ against a test
field concentrated at scale $1/N$ therefore converts, with an explicit
error, into the fixed continuum pairing
$\int_{\R^3}(V\cdot\nabla V)\cdot\Psi$ times $N^{3/2}$
(\S\ref{app:duality}); and that continuum pairing is nonzero for some
$\Psi$ because $\PP_{\R^3}(V\cdot\nabla V)\not\equiv0$, which is the
Fourier form of the elementary fact that the Oseen field is not a
stationary Euler flow (\S\ref{app:nondeg}). The bookkeeping is assembled
in \S\ref{app:final}.

\subsection{Conventions}
\label{app:conventions}

On $\T^3=(\R/2\pi\Z)^3$ we use the normalised measure
$d\mu=(2\pi)^{-3}dx$, so that for $u:\T^3\to\R^3$
\[
u(x)=\sum_{k\in\Z^3}\hat u(k)e^{ik\cdot x},\qquad
\hat u(k)=\int_{\T^3}u(x)e^{-ik\cdot x}\,d\mu(x),
\]
\[
\langle f,g\rangle=\int_{\T^3}f\cdot\bar g\,d\mu
=\sum_{k\in\Z^3}\hat f(k)\cdot\overline{\hat g(k)},
\qquad
H_r(u)=\sum_{k\in\Z^3}|k|^{2r}|\hat u(k)|^2 .
\]
On $\R^3$ we use the forward transform and the inverse-transform
bracket
\[
\mathcal F f(\zeta)=\int_{\R^3}f(y)e^{-i\zeta\cdot y}\,dy,
\qquad
\widetilde G(y)=\int_{\R^3}G(\xi)e^{i\xi\cdot y}\,d\xi ,
\]
so that $\mathcal F\widetilde G=(2\pi)^3G$ and
$\|\widetilde G\|_{L^2(\R^3)}^2=(2\pi)^3\|G\|_{L^2(\R^3)}^2$.

For $\xi\in\R^3\setminus\{0\}$ put
$P_\xi=I-\dfrac{\xi\otimes\xi}{|\xi|^2}$. Thus $P_\xi$ is the orthogonal
projection onto $\xi^\perp$: it is symmetric, idempotent, $P_\xi\xi=0$,
$\|P_\xi\|\le1$, $P_{-\xi}=P_\xi$, it is homogeneous of degree $0$ and
smooth on $\R^3\setminus\{0\}$, and consequently
$P_{k}=P_{k/N}$ for every $k\neq0$ and every $N>0$. The Leray
projection $\PP$ acts on $\T^3$ by
$\widehat{\PP f}(k)=P_k\hat f(k)$ for $k\neq0$, leaving the mean
untouched, and on $\R^3$ by the same symbol $P_\zeta$. On both
$L^2(\T^3)$ and $L^2(\R^3)$, $\PP$ is an orthogonal projection, hence
self-adjoint.

Throughout, $\|v_0\|$ is the Euclidean norm of $v_0\in\R^3$, and
$e_1,e_2,e_3$ is the standard basis. Repeated Latin indices are summed.
One notational warning: in this appendix $V$ always denotes the
continuum profile of Definition~\ref{app:Vdef}; the unrelated use of
$V$ for the modal variance $\mathrm{Var}_p(|k|^2)$ in
Theorem~\ref{main}(a) does not occur here.

\subsection{The cutoff class, the family, and its exact laws}
\label{app:laws}

We restate the definitions of the main text for completeness.

\begin{definition}[admissible cutoff; restatement of
Definition~\ref{admissible}]\label{app:admissible}
$\chi\in C^4([0,\infty);[0,1])$ with $\chi\equiv1$ on
$[0,\tfrac12]$ and $\operatorname{supp}\chi\subset[0,1]$. No
monotonicity is assumed --- it is used nowhere. Write
$c_\chi:=1+\|\chi'\|_{L^\infty}\ (\ge1)$.
\end{definition}

Such $\chi$ exist in abundance. The standard mollifier construction
gives a $C^\infty$ example, $\chi(t)=\Theta\big(2t-1\big)$ with
$\Theta(s)=1$ for $s\le0$, $\Theta(s)=0$ for $s\ge1$ and
$\Theta(s)=\frac{e^{-1/(1-s)}}{e^{-1/s}+e^{-1/(1-s)}}$ on $(0,1)$; the
exactly rational cutoff of Appendix~\ref{app:certificates}, built from
the degree-$9$ smoothstep, lies in $C^4\setminus C^5$ and is likewise
admissible.

\begin{remark}[what smoothness is actually used]\label{app:c1}
Definition~\ref{app:admissible} is still stronger than the proof
requires. The only place any derivative of $\chi$ is used
load-bearingly is Theorem~\ref{app:rate}, and it enters solely through
$c_\chi=1+\|\chi'\|_\infty$; Sections~\ref{app:duality} and
\ref{app:nondeg} use only $0\le\chi\le1$, $\chi\equiv1$ on
$[0,\tfrac12]$ and $\operatorname{supp}\chi\subset[0,1]$, and
Sections~\ref{app:laws}, \ref{app:lattice}, \ref{app:profile} use no
derivative of $\chi$ at all. Every statement of this appendix on which
Theorem~\ref{app:main} depends therefore holds verbatim for
$\chi\in C^1([0,\infty);[0,1])$ --- indeed for Lipschitz $\chi$ ---
with those support properties. The four derivatives are demanded only
so that Theorem~\ref{app:farfield} holds as stated with $M=4$; that
theorem is consumed only by the inert Remark~\ref{app:noperiod}, and
nothing in the proof of Theorem~\ref{app:main} uses it. In particular
$V\in C^\infty$ (Theorem~\ref{app:Vprops}(3)) and the Paley--Wiener
statement there are bought by the \emph{compact support} of $F$, not by
smoothness of $\chi$, and are unaffected.
\end{remark}

\begin{definition}\label{app:family}
Fix $v_0\in\R^3\setminus\{0\}$, an admissible $\chi$, and an integer
$N\ge2$. Define $u_N=u_N^{\chi,v_0}$ on $\T^3$ by
\begin{equation}\label{app:famdef}
\hat u_N(k)=\chi\Big(\frac{|k|}{N}\Big)\frac{P_kv_0}{|k|^{2}}
\quad(k\neq0),\qquad \hat u_N(0)=0 .
\end{equation}
\end{definition}

Since $\hat u_N(k)=0$ for $|k|>N$, $u_N$ is a finite trigonometric
polynomial banded at $|k|\le N$; its coefficients are real vectors and
no integrality of $v_0$ is required.

\begin{lemma}[basic structure]\label{app:structure}
$u_N$ is real, zero-mean, exactly divergence-free, $C^\infty$, and
$u_N\not\equiv0$.
\end{lemma}

\begin{proof}
Zero mean is $\hat u_N(0)=0$. Reality: the coefficients are real
vectors and $P_{-k}=P_k$, so
$\hat u_N(-k)=\hat u_N(k)=\overline{\hat u_N(k)}$, which is exactly the
reality condition $\hat u(-k)=\overline{\hat u(k)}$. Divergence-free:
$\widehat{\nabla\cdot u_N}(k)=ik\cdot\hat u_N(k)$ and $k\cdot P_kv_0=0$
for every $k\neq0$. Smoothness: a finite sum of exponentials. Finally,
among $e_1,e_2,e_3$ at least two are not parallel to $v_0$; pick such a
$k$ with $|k|=1$. Then $\chi(1/N)=1$ because $1/N\le\tfrac12$, and
$P_kv_0\neq0$, so $\hat u_N(k)\neq0$.
\end{proof}

The exact laws rest on a band-symmetry identity which, because $\chi$ is
not assumed monotone and the truncation is not a sharp ball, must be
available for a \emph{general} radial weight, with no positivity
assumption.

\begin{lemma}[band symmetry, lattice, general radial weight]
\label{app:bandsym}
Let $f:\Z^3\setminus\{0\}\to\R$ be radial, i.e.\ $f(k)=\varphi(|k|^2)$
for some $\varphi$, and suppose $\sum_{k\neq0}|k|^2|f(k)|<\infty$. Then
for all $i,j$
\[
\sum_{k\neq0}k_ik_j\,f(k)
=\delta_{ij}\,\frac13\sum_{k\neq0}|k|^2f(k).
\]
\end{lemma}

\begin{proof}
Absolute summability of $|k|^2f$ dominates every summand
$|k_ik_jf(k)|\le|k|^2|f(k)|$, so all rearrangements below are
legitimate. The index set $\Z^3\setminus\{0\}$ is invariant under each
sign flip $\sigma_i:k_i\mapsto-k_i$ and under all coordinate
permutations, and $f$, depending on $|k|^2$ only, is invariant under
both. For $i\neq j$ the bijection $\sigma_i$ negates the summand
$k_ik_jf(k)$, so that sum equals its own negative and vanishes. For
$i=j$, coordinate permutations carry the three sums $\sum_kk_i^2f(k)$
into one another, so they are equal; their sum is $\sum_k|k|^2f(k)$.
\end{proof}

\begin{lemma}[band symmetry, continuum]\label{app:bandsymcont}
Let $g:\R^3\to\R$ be radial with $\int_{\R^3}|\xi|^2|g(\xi)|\,d\xi<\infty$.
Then $\int\xi_i\xi_jg\,d\xi=\delta_{ij}\frac13\int|\xi|^2g\,d\xi$.
\end{lemma}

\begin{proof}
Identical: Lebesgue measure on $\R^3$ and $g$ are both invariant under
sign flips and permutations of the coordinates, and the hypothesis
dominates the integrand.
\end{proof}

\begin{theorem}[exact family laws]\label{app:exactlaws}
Set
\[
S_N^\chi:=\sum_{k\neq0}\frac{\chi(|k|/N)^2}{|k|^2},\qquad
T_N^\chi:=\sum_{k\neq0}\frac{\chi(|k|/N)^2}{|k|^4},\qquad
\Sigma_N^\chi:=\sum_{k\neq0}\frac{\chi(|k|/N)}{|k|^2},
\]
all three finite sums. Then
\[
H_0(u_N)=\tfrac23\|v_0\|^2\,T_N^\chi,\qquad
H_1(u_N)=\tfrac23\|v_0\|^2\,S_N^\chi,\qquad
N_0^2(u_N)=\frac{S_N^\chi}{T_N^\chi},\qquad
u_N(0)=\tfrac23\,\Sigma_N^\chi\,v_0 ,
\]
and consequently
$\|u_N\|_\infty^2/H_1(u_N)\ge 2(\Sigma_N^\chi)^2/(3S_N^\chi)$.
\end{theorem}

\begin{proof}
Since $P_k$ is an orthogonal projection,
$|P_kv_0|^2=\|v_0\|^2-(k\cdot v_0)^2/|k|^2$. Hence
\[
H_0=\sum_{k\neq0}\frac{\chi(|k|/N)^2}{|k|^4}\,|P_kv_0|^2
=\|v_0\|^2\sum_{k\neq0}\frac{\chi^2}{|k|^4}
-\sum_{k\neq0}\frac{\chi^2\,(k\cdot v_0)^2}{|k|^6}.
\]
Apply Lemma~\ref{app:bandsym} with the radial weight
$f(k)=\chi(|k|/N)^2|k|^{-6}$ (a finite sum, so the hypothesis is
trivially met):
\[
\sum_{k\neq0}\frac{\chi^2(k\cdot v_0)^2}{|k|^6}
=v_{0i}v_{0j}\sum_{k\neq0}k_ik_jf(k)
=\frac{\|v_0\|^2}{3}\sum_{k\neq0}\frac{\chi^2}{|k|^4}
=\frac{\|v_0\|^2}{3}T_N^\chi .
\]
Therefore $H_0=\|v_0\|^2T_N^\chi-\frac13\|v_0\|^2T_N^\chi
=\frac23\|v_0\|^2T_N^\chi$. The same computation with
$f=\chi^2|k|^{-4}$ applied to
$H_1=\sum_{k}|k|^2|\hat u_N(k)|^2$ gives
$H_1=\frac23\|v_0\|^2S_N^\chi$, and $N_0^2=H_1/H_0$ follows. For the
point value,
\[
u_N(0)=\sum_{k\neq0}\hat u_N(k)
=v_0\,\Sigma_N^\chi-\sum_{k\neq0}\frac{\chi\,k\,(k\cdot v_0)}{|k|^4},
\]
and Lemma~\ref{app:bandsym} with $f=\chi|k|^{-4}$ turns the last sum
into $\frac13v_0\Sigma_N^\chi$; hence
$u_N(0)=\frac23\Sigma_N^\chi v_0$. Finally
$\|u_N\|_\infty\ge|u_N(0)|=\frac23\Sigma_N^\chi\|v_0\|$, so
$\|u_N\|_\infty^2/H_1\ge\frac49(\Sigma_N^\chi)^2\|v_0\|^2
\big/\big(\frac23\|v_0\|^2S_N^\chi\big)
=\frac23(\Sigma_N^\chi)^2/S_N^\chi$.
\end{proof}

\subsection{Two-sided lattice bounds for the bandwidth}
\label{app:lattice}

The constants below are crude and explicit; no sharpness is attempted.

\begin{lemma}[dyadic shell count]\label{app:shell}
For every integer $j\ge0$,
$\#\{k\in\Z^3:2^j\le|k|<2^{j+1}\}\le64\cdot8^{j}$.
\end{lemma}

\begin{proof}
Such $k$ satisfy $|k|_\infty\le|k|<2^{j+1}$, hence
$|k_i|\le2^{j+1}-1$ for each $i$, so the count is at most
$(2^{j+2}-1)^3<2^{3j+6}=64\cdot8^{j}$.
\end{proof}

\begin{lemma}[upper bounds]\label{app:upper}
With $S_M:=\sum_{1\le|k|\le M}|k|^{-2}$ and
$T_\infty:=\sum_{k\neq0}|k|^{-4}$ one has
$S_M\le128M$ for every real $M\ge1$, and $T_\infty\le128$.
\end{lemma}

\begin{proof}
Let $J=\lfloor\log_2M\rfloor$. Every $k$ with $1\le|k|\le M$ lies in a
shell $2^j\le|k|<2^{j+1}$ with $0\le j\le J$, where $|k|^{-2}\le4^{-j}$.
By Lemma~\ref{app:shell},
$S_M\le\sum_{j=0}^{J}64\cdot8^{j}4^{-j}=64(2^{J+1}-1)<128\cdot2^{J}\le128M$.
Similarly $T_\infty\le\sum_{j\ge0}64\cdot8^{j}16^{-j}=64\cdot2=128$.
\end{proof}

\begin{lemma}[lower bound]\label{app:lower}
For every real $M\ge16$ one has $S_M\ge M/87$.
\end{lemma}

\begin{proof}
Consider
$\mathcal B:=\{k\in\Z^3:\ \tfrac{M}{2\sqrt3}<k_i\le\tfrac{M}{\sqrt3}\
\text{for }i=1,2,3\}$. For $k\in\mathcal B$ we have
$|k|^2\le3(M/\sqrt3)^2=M^2$ and $|k|^2>3M^2/12=M^2/4$, so every
$k\in\mathcal B$ is counted in $S_M$ (in particular $k\neq0$) and
contributes $|k|^{-2}\ge M^{-2}$. Each coordinate ranges over a
half-open interval of length $M/(2\sqrt3)$, which contains at least
$\frac{M}{2\sqrt3}-1$ integers; for $M\ge16$,
$\frac{M}{2\sqrt3}-1\ge0.28868M-0.0625M\ge0.226M$. Hence
$\#\mathcal B\ge(0.226M)^3>0.01154M^3$ and
$S_M\ge0.01154M^3\cdot M^{-2}=0.01154M\ge M/87$.
\end{proof}

\begin{proposition}[two-sided bandwidth bounds]\label{app:twosided}
For every admissible $\chi$, every $v_0\neq0$ and every integer
$N\ge32$,
\[
\frac{N}{348}\ \le\ S_N^\chi\ \le\ 128N,\qquad
6\ \le\ T_N^\chi\ \le\ 128,\qquad
\frac{N}{44544}\ \le\ N_0^2(u_N)\ \le\ \frac{64}{3}N .
\]
\end{proposition}

\begin{proof}
Upper bounds: $0\le\chi\le1$ and $\chi(|k|/N)=0$ for $|k|>N$, so
$S_N^\chi\le S_N\le128N$ and $T_N^\chi\le T_\infty\le128$ by
Lemma~\ref{app:upper}. Lower bound for $T_N^\chi$: for $N\ge2$ one has
$\chi(1/N)=1$, and the six lattice points with $|k|=1$ each contribute
$1$, so $T_N^\chi\ge6$. Lower bound for $S_N^\chi$: $\chi(|k|/N)=1$
whenever $|k|\le N/2$, hence $S_N^\chi\ge S_{N/2}$; for $N\ge32$ we have
$N/2\ge16$, so Lemma~\ref{app:lower} gives
$S_N^\chi\ge (N/2)/87\ge N/348$. Combining with
$N_0^2=S_N^\chi/T_N^\chi$ (Theorem~\ref{app:exactlaws}) yields
$N_0^2\ge\frac{N/348}{128}=\frac{N}{44544}$ and
$N_0^2\le\frac{128N}{6}=\frac{64}{3}N$.
\end{proof}

\subsection{The profile $V$}
\label{app:profile}

\begin{definition}\label{app:Vdef}
On $\R^3$ put
\[
h(\xi):=\frac{P_\xi v_0}{|\xi|^2}\quad(\xi\neq0),\qquad
F(\xi):=\chi(|\xi|)\,h(\xi)\quad(\xi\neq0),\quad F(0):=0,
\]
\[
V(y):=\widetilde F(y)=\int_{\R^3}F(\xi)e^{i\xi\cdot y}\,d\xi .
\]
\end{definition}

Because $P_\xi$ is homogeneous of degree $0$ we have
$P_{k/N}=P_k$ and therefore
$F(k/N)=\chi(|k|/N)\,N^2P_kv_0/|k|^2$, i.e.\ the exact scaling identity
\begin{equation}\label{app:scaling}
\hat u_N(k)=N^{-2}F(k/N)\qquad\text{for every }k\in\Z^3
\end{equation}
(the case $k=0$ holds by the conventions $\hat u_N(0)=0=F(0)$).
Identity~\eqref{app:scaling} is the engine of the whole argument: it
says that the family is a single fixed continuum profile sampled on the
lattice $N^{-1}\Z^3$.

\begin{theorem}[elementary properties of $V$]\label{app:Vprops}
\begin{enumerate}
\item[(1)] $|F(\xi)|\le\|v_0\|\,|\xi|^{-2}\mathbf 1_{\{|\xi|\le1\}}$;
hence $F\in L^1(\R^3)\cap L^p(\R^3)$ for every $p<3/2$, with
$\|F\|_{L^1}\le4\pi\|v_0\|$.
\item[(2)] $V$ is real, even, bounded and continuous, with
$\|V\|_\infty\le4\pi\|v_0\|$.
\item[(3)] $V\in C^\infty(\R^3)$, and for every multi-index $\beta$
\[
\partial^\beta V(y)=\int_{\R^3}(i\xi)^\beta F(\xi)e^{i\xi\cdot y}\,d\xi,
\qquad
\|\partial^\beta V\|_\infty\le\|v_0\|\int_{|\xi|\le1}|\xi|^{|\beta|-2}d\xi
=\frac{4\pi\|v_0\|}{|\beta|+1} .
\]
In particular each $\|\partial_iV_j\|_\infty\le2\pi\|v_0\|$, so
$\|\nabla V\|_\infty\le6\pi\|v_0\|$ for the Frobenius norm. Moreover $V$
extends to an entire function on $\mathbb C^3$ of exponential type
$\le1$.
\item[(4)] $\nabla\cdot V\equiv0$ on $\R^3$.
\item[(5)] $V(0)=\frac{8\pi}{3}\big(\int_0^\infty\chi(r)\,dr\big)v_0\neq0$.
\end{enumerate}
\end{theorem}

\begin{proof}
(1) $|P_\xi v_0|\le\|v_0\|$, $0\le\chi\le1$ and
$\operatorname{supp}\chi\subset[0,1]$ give the pointwise bound;
$\int_{|\xi|\le1}|\xi|^{-2}d\xi=4\pi\int_0^1dr=4\pi$, and
$\int_{|\xi|\le1}|\xi|^{-2p}d\xi<\infty$ exactly when $2p<3$.

(2) $F\in L^1$ gives, by dominated convergence, continuity of $V$ and
$\|V\|_\infty\le\|F\|_{L^1}\le4\pi\|v_0\|$. $F$ is real-valued and even
($P_{-\xi}=P_\xi$), so
$V(y)=\int F(\xi)\cos(\xi\cdot y)\,d\xi$ is real and even.

(3) By (1), $\xi^\beta F\in L^1$ for every $\beta$, since the exponent
$|\beta|-2>-3$ is locally integrable in $\R^3$ and the support is
bounded. Differentiation under the integral sign is then justified by
dominated convergence, giving the displayed formula and, by (1) again,
the bound
$\|v_0\|\int_{|\xi|\le1}|\xi|^{|\beta|-2}d\xi
=4\pi\|v_0\|\int_0^1r^{|\beta|}dr=4\pi\|v_0\|/(|\beta|+1)$. Since $F$
is an $L^1$ function supported in $\{|\xi|\le1\}$, the integral
$\int F(\xi)e^{i\xi\cdot y}d\xi$ converges for all complex $y$ and is
entire in $y$, with $|V(y)|\le\|F\|_{L^1}e^{|\operatorname{Im}y|}$;
this is exponential type $\le1$.

(4) $\nabla\cdot V(y)=\int i\,\xi\cdot F(\xi)e^{i\xi\cdot y}d\xi$ and
$\xi\cdot F(\xi)=\chi(|\xi|)\,\xi\cdot P_\xi v_0/|\xi|^2=0$ pointwise.

(5) $V(0)=\int F
=v_0\int\frac{\chi(|\xi|)}{|\xi|^2}d\xi
-\int\frac{\chi(|\xi|)\,\xi\,(\xi\cdot v_0)}{|\xi|^4}d\xi$.
Lemma~\ref{app:bandsymcont} with the radial
$g(\xi)=\chi(|\xi|)|\xi|^{-4}$ (integrable against $|\xi|^2$ by (1))
gives $\int\xi_i\xi_jg=\frac{\delta_{ij}}3\int|\xi|^2g
=\frac{\delta_{ij}}3\int\chi(|\xi|)|\xi|^{-2}d\xi$, so the second
integral equals $\frac13v_0\int\chi(|\xi|)|\xi|^{-2}d\xi$. Hence
\[
V(0)=\frac23v_0\int_{\R^3}\frac{\chi(|\xi|)}{|\xi|^2}\,d\xi
=\frac23v_0\cdot4\pi\int_0^\infty\chi(r)\,dr
=\frac{8\pi}{3}\Big(\int_0^\infty\chi\Big)v_0 ,
\]
and $\int_0^\infty\chi\ge\int_0^{1/2}1=\frac12>0$ because $\chi\ge0$ and
$\chi\equiv1$ on $[0,\tfrac12]$.
\end{proof}

\begin{remark}[structural form]\label{app:oseenform}
$F=P_\xi\big(v_0\chi(|\xi|)|\xi|^{-2}\big)$, so
$V=\PP\big(v_0\,g_\chi\big)$ where $g_\chi$ is the inverse transform of
the smoothed Newtonian symbol $\chi(|\cdot|)|\cdot|^{-2}$; solving
$\Delta\phi=g_\chi$ with $\phi$ radial exhibits $V$ in the Oseen normal
form $V(y)=a(r)v_0+b(r)(v_0\cdot\hat y)\hat y$ with
$a=\phi''+\phi'/r$, $b=\phi'/r-\phi''$. This form is not used below; it
explains why the far field of $V$ is exactly the Oseen profile
(Theorem~\ref{app:farfield}).
\end{remark}

We record here the two classical transform facts used later.

\begin{remark}[classical input: Riesz-potential transforms]
\label{app:classical-riesz}
For $0<\alpha<3$, and by analytic continuation for every $\alpha$ at
which the Gamma factors have no pole, the homogeneous function
$|y|^{-\alpha}$ on $\R^3$ is a tempered distribution whose Fourier
transform is the homogeneous distribution
\[
\mathcal F\big[|y|^{-\alpha}\big](\zeta)=c_{3,\alpha}|\zeta|^{\alpha-3},
\qquad
c_{3,\alpha}=\frac{2^{3-\alpha}\pi^{3/2}\,\Gamma\!\big(\frac{3-\alpha}{2}\big)}
{\Gamma\!\big(\frac{\alpha}{2}\big)} .
\]
See, e.g., \cite{SteinWeiss,GelfandShilov,Grafakos} (we give no section
locators, not having verified them). We use
$\Gamma(\tfrac12)=\sqrt\pi$, $\Gamma(-\tfrac12)=-2\sqrt\pi$,
$\Gamma(-\tfrac32)=\tfrac43\sqrt\pi$, which give
\[
c_{3,-1}=-8\pi,\quad c_{3,1}=4\pi,\quad c_{3,2}=2\pi^2,\quad
c_{3,4}=-\pi^2,\quad c_{3,6}=\tfrac{\pi^2}{12},
\]
together with the elementary rules
$\mathcal F[y_ay_bf]=-\partial_a\partial_b\mathcal F[f]$ and
$\mathcal F[y_ay_by_cy_df]=\partial_a\partial_b\partial_c\partial_d\mathcal F[f]$,
valid for tempered $f$.
\end{remark}

\begin{remark}[classical input: products and convolutions of homogeneous
distributions]\label{app:classical-homog}
Let $A,B\in L^1_{\rm loc}(\R^3)$ be homogeneous of degrees $-a,-b$ with
$0<a,b<3$ and $a+b<3$, and smooth away from the origin. Then $AB$ is
locally integrable and homogeneous of degree $-(a+b)$, hence tempered;
$\mathcal FA,\mathcal FB$ are locally integrable, homogeneous of degrees
$a-3,b-3$ and smooth away from the origin; the convolution
$\mathcal FA*\mathcal FB$ converges absolutely at every $\zeta\neq0$
(because $(3-a)+(3-b)>3$) and defines a locally integrable homogeneous
function of degree $3-a-b$; and the exchange formula
\[
\mathcal F[AB]=(2\pi)^{-3}\,\mathcal FA*\mathcal FB
\]
holds as tempered distributions. See, e.g.,
\cite{GelfandShilov,Grafakos}. We use this only in
Theorem~\ref{app:tauformula}, with $a=b=1$.
\end{remark}

\subsection{The far field, and why periodisation is unavailable}
\label{app:farfield-sec}

\begin{theorem}[far-field asymptotics]\label{app:farfield}
Let
\[
W(y):=\frac{v_0}{|y|}+\frac{(v_0\cdot y)\,y}{|y|^3}\qquad(y\neq0).
\]
Then $\widetilde h=\pi^2W$ as tempered distributions, and
\[
V(y)=\pi^2W(y)+\rho(y),\qquad
|\rho(y)|\le C_M(\chi,v_0)|y|^{-M}\quad\text{for }|y|\ge1
\text{ and every }2\le M\le4 .
\]
(Four derivatives are all that Definition~\ref{app:admissible}
supplies, and $M=4$ is all that Remark~\ref{app:noperiod} consumes; if
$\chi\in C^m$ with $m\ge2$ the same proof gives every
$2\le M\le m$.)
\end{theorem}

\begin{proof}
\emph{(a) The identity $\widetilde h=\pi^2W$.} Let
$\mathcal O_{ij}(y)=\frac1{8\pi}\big(\frac{\delta_{ij}}{|y|}
+\frac{y_iy_j}{|y|^3}\big)$ be the Oseen tensor, i.e.\ the fundamental
solution of the steady Stokes system
\cite{LadyzhenskayaBook,Galdi2011}, and
$h_{ij}(\xi)=\frac{\delta_{ij}}{|\xi|^2}-\frac{\xi_i\xi_j}{|\xi|^4}$, so
that $h_i=h_{ij}v_{0j}$. We verify $\mathcal F[\mathcal O_{ij}]=h_{ij}$
from Remark~\ref{app:classical-riesz}. First,
$\mathcal F[|y|^{-1}]=c_{3,1}|\zeta|^{-2}=4\pi|\zeta|^{-2}$, so
$\mathcal F\big[\tfrac{\delta_{ij}}{8\pi|y|}\big]
=\tfrac{\delta_{ij}}{2|\zeta|^{2}}$. Second, differentiating
$\partial_i|y|=y_i/|y|$ once more gives the pointwise identity
\[
\partial_i\partial_j|y|=\frac{\delta_{ij}}{|y|}-\frac{y_iy_j}{|y|^3},
\qquad\text{i.e.}\qquad
\frac{y_iy_j}{|y|^3}=\frac{\delta_{ij}}{|y|}-\partial_i\partial_j|y| .
\]
With $\mathcal F[|y|]=c_{3,-1}|\zeta|^{-4}=-8\pi|\zeta|^{-4}$ we get
$\mathcal F[\partial_i\partial_j|y|]
=(i\zeta_i)(i\zeta_j)\mathcal F[|y|]=8\pi\zeta_i\zeta_j|\zeta|^{-4}$, hence
\[
\mathcal F\Big[\frac{y_iy_j}{|y|^3}\Big]
=\frac{4\pi\delta_{ij}}{|\zeta|^2}-\frac{8\pi\zeta_i\zeta_j}{|\zeta|^4}.
\]
Therefore
$\mathcal F[\mathcal O_{ij}]
=\frac1{8\pi}\big(\frac{4\pi\delta_{ij}}{|\zeta|^2}
+\frac{4\pi\delta_{ij}}{|\zeta|^2}
-\frac{8\pi\zeta_i\zeta_j}{|\zeta|^4}\big)
=\frac{\delta_{ij}}{|\zeta|^2}-\frac{\zeta_i\zeta_j}{|\zeta|^4}=h_{ij}$.
All these are identities between homogeneous tempered distributions;
possible ambiguities are supported at the origin only and are excluded
because both sides are homogeneous of degree $-2>-3$, hence locally
integrable, with no room for a delta term. Fourier inversion
($\mathcal F\widetilde G=(2\pi)^3G$) now gives
$\widetilde h=(2\pi)^3\mathcal O\,v_0=\frac{(2\pi)^3}{8\pi}W=\pi^2W$,
where we used $8\pi\,\mathcal O_{ij}v_{0j}=W_i$.

\emph{(b) The remainder.} Write $G:=(1-\chi(|\cdot|))h$, so that
$F-h=-G$ and $\rho:=-\widetilde G$. Then $G$ vanishes on
$\{|\xi|<\tfrac12\}$ and is $C^4$ on all of $\R^3$; since $h$ is
homogeneous of degree $-2$ and smooth off the origin, and
$1-\chi(|\cdot|)$ is $C^4$ with its first four derivatives bounded, we
get
\[
|\partial^\beta G(\xi)|\le C_\beta(\chi)\|v_0\|\,\langle\xi\rangle^{-2-|\beta|}
\qquad(\xi\in\R^3,\ |\beta|\le4),
\]
which for $2\le|\beta|\le4$ places $\partial^\beta G$ in $L^1(\R^3)$
(the exponent $2+|\beta|>3$). From
the distributional identity
$y^\beta\widetilde G=i^{|\beta|}\,\widetilde{\partial^\beta G}$, the
right-hand side for $2\le|\beta|\le4$ is a bounded continuous function
with supremum at most $\|\partial^\beta G\|_{L^1}$. Choosing
$\beta=(M,0,0),(0,M,0),(0,0,M)$ for a given $2\le M\le4$ and using
$|y|^M\le3^{M/2}\max_i|y_i|^M$ shows that $\widetilde G$ agrees on
$\{|y|\ge1\}$ with a continuous function obeying
$|\widetilde G(y)|\le C_M(\chi,v_0)|y|^{-M}$. Since
$V=\widetilde F=\widetilde h-\widetilde G=\pi^2W+\rho$, the claim
follows.
\end{proof}

Thus the decay of $V$ is \emph{exactly} $|y|^{-1}$, with a
$\chi$-independent leading profile $\pi^2W$ satisfying
\begin{equation}\label{app:Wpos}
v_0\cdot W(y)=\frac{\|v_0\|^2}{|y|}+\frac{(v_0\cdot y)^2}{|y|^3}
\ \ge\ \frac{\|v_0\|^2}{|y|}\ >\ 0 .
\end{equation}

\begin{remark}[the naive periodisation argument is unavailable]
\label{app:noperiod}
It is tempting to reconstruct $u_N$ from $V$ by Poisson summation,
$u_N(x)\stackrel{?}{=}N\sum_{m\in\Z^3}V(N(x+2\pi m))$ with the $m\neq0$
terms uniformly $O(1/N)$. \textbf{The naive absolutely convergent
periodisation argument is unavailable}: the series of translates
diverges, because $V$ decays only like $|y|^{-1}$ and its leading
profile has a sign-definite spherical mean in the $v_0$ direction, so
there is no cancellation to exploit and no absolute convergence to
appeal to. Indeed, by Theorem~\ref{app:farfield} with $M=4$, for
$m\neq0$ and $x\in[-\pi,\pi)^3$ one has $|N(x+2\pi m)|\ge N\pi\ge1$ and
therefore, using \eqref{app:Wpos},
\[
v_0\cdot V\big(N(x+2\pi m)\big)
\ \ge\ \frac{\pi^2\|v_0\|^2}{N|x+2\pi m|}
-\frac{C_4\|v_0\|}{N^4|x+2\pi m|^4}.
\]
The shell $|m|_\infty=n$ contains $O(n^2)$ points with
$|x+2\pi m|\asymp n$, so $\sum_{m\neq0}|x+2\pi m|^{-4}<\infty$ while
$\sum_{m\neq0}|x+2\pi m|^{-1}=\infty$; hence
$\sum_{m\neq0}v_0\cdot V(N(x+2\pi m))=+\infty$ for every $x$ and every
$N$. The terms are moreover eventually positive, so the divergence is
not of a kind that a regular linear summation method would remove: for
$a_m\ge0$ with $\sum a_m=\infty$, monotone convergence gives
$\lim_{\varepsilon\downarrow0}\sum_me^{-\varepsilon|m|}a_m=+\infty$, so
Abel summation in particular does not converge here. We make no
assertion beyond this, and none is needed: the point is only that the
naive absolutely convergent periodisation argument is not available,
because $V$ decays only like $|y|^{-1}$ with a non-oscillating leading
profile whose spherical mean, $\tfrac43v_0/r$, is sign-definite in the
$v_0$ direction. If a periodisation formula
is wanted purely for orientation, the correct \emph{renormalised} one,
\[
u_N(x)=N\Big[V(Nx)+\sum_{m\neq0}\big(V(N(x+2\pi m))-V(2\pi Nm)\big)\Big]
+\text{const},
\]
does converge absolutely, the bracketed differences being $O(|m|^{-2})$
by Theorem~\ref{app:farfield}. \textbf{Nothing in this appendix uses any
periodisation identity}: \S\ref{app:riemann} replaces it by an exact
lattice-Riemann-sum identity, which is elementary, quantitative, and
additionally delivers the derivative statement. This remark, and with it
Theorem~\ref{app:farfield}, is inert for the proof of
Theorem~\ref{capacity}.
\end{remark}

\subsection{The inner rescaling: an exact lattice-Riemann identity with
an explicit rate}
\label{app:riemann}

For $\Phi:\R^3\to\mathbb C^d$ put
\[
\mathcal R_N[\Phi]:=N^{-3}\sum_{k\in\Z^3}\Phi(k/N)
-\int_{\R^3}\Phi(\xi)\,d\xi ,
\]
whenever both terms converge absolutely.

\begin{lemma}[exact inner rescaling]\label{app:innerexact}
For every $y\in\R^3$ and every integer $N\ge2$,
\[
u_N(y/N)=N\big(V(y)+E_N(y)\big),\qquad
\partial_iu_{N,j}(y/N)=N^2\big(\partial_iV_j(y)+E'_{N,ij}(y)\big),
\]
where, with $\Phi_y(\xi):=F(\xi)e^{i\xi\cdot y}$ and
$\Phi'_{y,ij}(\xi):=i\xi_iF_j(\xi)e^{i\xi\cdot y}$ (both set to $0$ at
$\xi=0$),
\[
E_N(y)=\mathcal R_N[\Phi_y],\qquad
E'_{N,ij}(y)=\mathcal R_N[\Phi'_{y,ij}] .
\]
These are identities, not approximations.
\end{lemma}

\begin{proof}
Both series are finite (support $|k|\le N$) and both integrals converge
absolutely by Theorem~\ref{app:Vprops}(1),(3). By \eqref{app:scaling},
\[
u_N(y/N)=\sum_{k\in\Z^3}\hat u_N(k)e^{i(k/N)\cdot y}
=\sum_{k}N^{-2}F(k/N)e^{i(k/N)\cdot y}
=N\cdot N^{-3}\sum_k\Phi_y(k/N),
\]
which equals $N\big(\int\Phi_y+\mathcal R_N[\Phi_y]\big)
=N\big(V(y)+E_N(y)\big)$, using $V(y)=\int\Phi_y$
(Definition~\ref{app:Vdef}) and $\Phi_y(0)=F(0)=0$. For the gradient,
$\partial_iu_{N,j}(x)=\sum_kik_i\hat u_{N,j}(k)e^{ik\cdot x}$; at
$x=y/N$ this is
\[
\sum_kiN^{-2}k_iF_j(k/N)e^{i(k/N)\cdot y}
=N^{-1}\sum_ki(k_i/N)F_j(k/N)e^{i(k/N)\cdot y}
=N^{2}\cdot N^{-3}\sum_k\Phi'_{y,ij}(k/N),
\]
and $\int\Phi'_{y,ij}=\partial_iV_j(y)$ by
Theorem~\ref{app:Vprops}(3).
\end{proof}

Everything now rests on bounding a Riemann-sum defect for an explicit,
compactly supported integrand with a single integrable algebraic
singularity.

\begin{theorem}[$C^1_{\rm loc}$ convergence with explicit rate]
\label{app:rate}
There is an absolute constant $A$ (the proof gives $A=3000$) such that
for every real $R\ge1$ and every integer $N\ge8R$,
\[
\varepsilon_N:=\sup_{|y|\le R}|E_N(y)|
\ \le\ A\,\|v_0\|\,\frac{c_\chi(1+\log_2N)+R}{N},
\qquad
\varepsilon'_N:=\sup_{|y|\le R}|E'_N(y)|
\ \le\ A\,\|v_0\|\,\frac{c_\chi+R}{N} .
\]
In particular $\varepsilon_N+\varepsilon'_N=O(N^{-1}\log N)\to0$ as
$N\to\infty$, for each fixed $R$.
\end{theorem}

\begin{proof}
Partition $\R^3$ into the half-open cubes
$Q_k=\frac kN+[0,\frac1N)^3$, $k\in\Z^3$, of side $1/N$, volume
$N^{-3}$ and diameter $\sqrt3/N$. For any $\Phi$ for which both the sum
and the integral converge absolutely,
\begin{equation}\label{app:cellwise}
\mathcal R_N[\Phi]=\sum_{k\in\Z^3}\int_{Q_k}
\big[\Phi(k/N)-\Phi(\xi)\big]\,d\xi ,
\end{equation}
since $|Q_k|=N^{-3}$. Absolute convergence holds for
$\Phi=\Phi_y$ and $\Phi=\Phi'_{y,ij}$: both are supported in
$\{|\xi|\le1\}$ and dominated there by $\|v_0\||\xi|^{-2}$ respectively
$\|v_0\||\xi|^{-1}$, which are integrable on $\R^3$; and the lattice
sums are finite sums, bounded by
$N^{-3}\sum_{0<|k|\le N}\|v_0\|N^{2}|k|^{-2}
=\|v_0\|N^{-1}S_N\le128\|v_0\|$ and by
$N^{-3}\sum_{0<|k|\le N}\|v_0\|N|k|^{-1}
\le\|v_0\|N^{-2}\cdot N\,S_N\le128\|v_0\|$ respectively, using
$|k|^{-1}\le N|k|^{-2}$ for $|k|\le N$ and Lemma~\ref{app:upper}.

\emph{Pointwise bounds.} On $0<|\xi|\le1$ we have
$|F(\xi)|\le\|v_0\||\xi|^{-2}$, and by the product rule applied to
$F=\chi(|\xi|)|\xi|^{-2}P_\xi v_0$,
\[
|\nabla F(\xi)|\ \le\
\underbrace{c_\chi\|v_0\||\xi|^{-2}}_{\nabla\chi(|\xi|)}
+\underbrace{2\|v_0\||\xi|^{-3}}_{\chi\,\nabla(|\xi|^{-2})}
+\underbrace{4\|v_0\||\xi|^{-3}}_{\chi|\xi|^{-2}\nabla(P_\xi v_0)}
\ \le\ 7c_\chi\|v_0\|\,|\xi|^{-3},
\]
where we used $c_\chi\ge1$, $|\xi|\le1$, $|\nabla(|\xi|^{-2})|
=2|\xi|^{-3}$, and, from
$\partial_a\big(\xi_i\xi_j/|\xi|^2\big)
=(\delta_{ai}\xi_j+\delta_{aj}\xi_i)/|\xi|^2
-2\xi_i\xi_j\xi_a/|\xi|^4$, the honest bound
$|\nabla(P_\xi v_0)|\le4\|v_0\|/|\xi|$. Consequently, for $|y|\le R$,
\begin{equation}\label{app:phibounds}
|\Phi_y|\le\frac{\|v_0\|}{|\xi|^{2}},\quad
|\nabla\Phi_y|\le\frac{7c_\chi\|v_0\|}{|\xi|^{3}}+\frac{R\|v_0\|}{|\xi|^{2}},
\qquad
|\Phi'_y|\le\frac{\|v_0\|}{|\xi|},\quad
|\nabla\Phi'_y|\le\frac{8c_\chi\|v_0\|}{|\xi|^{2}}+\frac{R\|v_0\|}{|\xi|} .
\end{equation}
(For the last one, $\nabla(i\xi_iF_je^{i\xi y})$ has the three terms
$|F|\le\|v_0\||\xi|^{-2}$, $|\xi||\nabla F|\le7c_\chi\|v_0\||\xi|^{-2}$
and $R|\xi||F|\le R\|v_0\||\xi|^{-1}$.)

\emph{The core.} Let
$\mathcal K_0:=\{k:\ Q_k\cap\{|\xi|\le4\sqrt3/N\}\neq\emptyset\}$. Any
such $Q_k$ lies in $\{|\xi|\le5\sqrt3/N\}$ and has $|k|\le5\sqrt3<9$,
i.e.\ $|k|\le8$. Note that $k=0$ is in the core, and that although
$\Phi_y(0)=0$ by convention while $\int_{Q_0}\Phi_y\neq0$, both are
bounded by the two terms below, so no separate treatment is needed. The
core contributes to \eqref{app:cellwise} at most
\[
\int_{|\xi|\le5\sqrt3/N}|\Phi_y|\,d\xi
+N^{-3}\!\!\sum_{0<|k|\le8}\!|\Phi_y(k/N)|
\ \le\ \|v_0\|\Big(\frac{4\pi\cdot5\sqrt3}{N}
+\frac1N\!\!\sum_{0<|k|\le8}\!\frac1{|k|^{2}}\Big)
\ \le\ \frac{1133\,\|v_0\|}{N},
\]
using $\int_{|\xi|\le a}|\xi|^{-2}d\xi=4\pi a$, then
$|\Phi_y(k/N)|\le\|v_0\|N^2|k|^{-2}$, and finally
$\sum_{0<|k|\le8}|k|^{-2}\le S_8\le1024$ by Lemma~\ref{app:upper}
($4\pi\cdot5\sqrt3\le109$, and $109+1024=1133$).

\emph{Dyadic annuli.} Put $J:=\lfloor\log_2\big(N/(4\sqrt3)\big)\rfloor$,
which is $\ge0$ because $N\ge8R\ge8>4\sqrt3$. Let $k\notin\mathcal K_0$
and set $r_k:=\inf_{\xi\in \overline{Q_k}}|\xi|$, so $r_k>4\sqrt3/N>0$.
If $r_k>1$ then $\Phi_y$ vanishes identically on $\overline{Q_k}$ and
that cell contributes $0$; otherwise choose the unique integer $j(k)$
with $2^{-j(k)-1}<r_k\le2^{-j(k)}$. Then $j(k)\ge0$, and
$2^{-j(k)}\ge r_k>4\sqrt3/N$ forces $2^{j(k)}<N/(4\sqrt3)$, i.e.\
$0\le j(k)\le J$; so the annuli
$A_j:=\{2^{-j-1}<|\xi|\le2^{-j}\}$, $0\le j\le J$, together with the
core, do cover $\operatorname{supp}\Phi_y$.

Fix $j$ and consider the cells with $j(k)=j$. They are pairwise
disjoint. Each satisfies $\overline{Q_k}\subset\{|\xi|\ge2^{-j-1}\}$ and
$\overline{Q_k}\subset\{|\xi|\le r_k+\sqrt3/N\le\tfrac54 2^{-j}\}$,
because $\sqrt3/N\le\tfrac14 2^{-j}$ (which follows from
$2^{-j}\ge2^{-J}\ge4\sqrt3/N$). Hence all of them lie in the fixed
annulus $B_j:=\{2^{-j-1}\le|\xi|\le\tfrac54 2^{-j}\}$, whose volume is
$\frac{4\pi}3\big((\tfrac54)^3-(\tfrac12)^3\big)2^{-3j}\le9\cdot2^{-3j}$;
therefore $\sum_{j(k)=j}|Q_k|\le9\cdot2^{-3j}$.

Since $\overline{Q_k}$ is convex, avoids the origin, and $\Phi_y$ is
$C^1$ there, the mean value inequality gives
$\big|\int_{Q_k}[\Phi_y(k/N)-\Phi_y(\xi)]d\xi\big|
\le\frac{\sqrt3}{N}\sup_{\overline{Q_k}}|\nabla\Phi_y|\cdot|Q_k|$. By
\eqref{app:phibounds} and $|\xi|\ge2^{-j-1}$ on $\overline{Q_k}$,
\[
\sup_{\overline{Q_k}}|\nabla\Phi_y|
\le 7c_\chi\|v_0\|2^{3(j+1)}+R\|v_0\|2^{2(j+1)}
=56c_\chi\|v_0\|8^{j}+4R\|v_0\|4^{j}.
\]
Summing over the cells with $j(k)=j$,
\[
\sum_{j(k)=j}\Big|\int_{Q_k}[\Phi_y(k/N)-\Phi_y]\Big|
\le\frac{\sqrt3}{N}\big(56c_\chi\|v_0\|8^{j}+4R\|v_0\|4^{j}\big)
\cdot9\cdot2^{-3j}
\le\frac{\|v_0\|}{N}\big(873\,c_\chi+63\,R\,2^{-j}\big).
\]
Summing over $0\le j\le J\le\log_2N$ and adding the core bound,
\[
|E_N(y)|\le\frac{\|v_0\|}{N}
\Big(1133+873\,c_\chi(1+\log_2N)+126R\Big)
\le\frac{3000\,\|v_0\|}{N}\big(c_\chi(1+\log_2N)+R\big),
\]
using $c_\chi\ge1$ and $1+\log_2N\ge1$. This is the first display.

\emph{The gradient.} Repeat verbatim with $\Phi'_{y,ij}$. The core
contributes at most
\[
\int_{|\xi|\le5\sqrt3/N}\frac{\|v_0\|}{|\xi|}\,d\xi
+N^{-3}\!\!\sum_{0<|k|\le8}\!\frac{\|v_0\|N}{|k|}
=\|v_0\|\Big(\frac{2\pi\cdot75}{N^{2}}
+\frac1{N^{2}}\!\!\sum_{0<|k|\le8}\!\frac1{|k|}\Big)
\le\frac{1816\,\|v_0\|}{N^{2}}\le\frac{1816\,\|v_0\|}{N},
\]
where $\int_{|\xi|\le a}|\xi|^{-1}d\xi=2\pi a^2$ and, by
Lemma~\ref{app:shell},
$\sum_{0<|k|\le8}|k|^{-1}\le\sum_{j=0}^{2}64\cdot8^{j}2^{-j}=1344$
(and $2\pi\cdot75\le472$, $472+1344=1816$). On the annuli,
$\sup_{\overline{Q_k}}|\nabla\Phi'_y|
\le8c_\chi\|v_0\|2^{2(j+1)}+R\|v_0\|2^{j+1}
=32c_\chi\|v_0\|4^{j}+2R\|v_0\|2^{j}$, so the level-$j$ contribution is
at most
\[
\frac{\sqrt3}{N}\big(32c_\chi\|v_0\|4^{j}+2R\|v_0\|2^{j}\big)
\cdot9\cdot2^{-3j}
\le\frac{\|v_0\|}{N}\big(499\,c_\chi2^{-j}+32\,R\,4^{-j}\big),
\]
and now the geometric series converge \emph{without} a logarithm,
summing over all $j\ge0$ to at most
$\frac{\|v_0\|}{N}(998c_\chi+43R)$. Adding the core,
$|E'_N(y)|\le\frac{\|v_0\|}{N}(1816+998c_\chi+43R)
\le\frac{3000\|v_0\|}{N}(c_\chi+R)$. The absence of a logarithm is
structural: the singularity of $\Phi'_y$ is only $|\xi|^{-1}$, the extra
factor $\xi$ taming the origin. The derivative estimate is thus not a
consequence of the $C^0$ one, but it is the easier of the two.
\end{proof}

\subsection{The concentrated test field and the duality lower bound}
\label{app:duality}

\begin{definition}\label{app:testfield}
Let $\Psi\in C_c^\infty(\R^3;\R^3)$ be \emph{real}, with
$\nabla\cdot\Psi=0$, $\operatorname{supp}\Psi\subset\{|y|\le R\}$ for
some $R\ge1$, and $\Psi\not\equiv0$. For $N>R/\pi$ define on $\T^3$
\[
\psi_N(x):=N^{3/2}\sum_{m\in\Z^3}\Psi\big(N(x+2\pi m)\big).
\]
\end{definition}

\begin{lemma}[properties of $\psi_N$]\label{app:psiprops}
Let $N>R/\pi$. Then $\psi_N$ is a well-defined real $C^\infty$ vector
field on $\T^3$; it is divergence-free and has zero mean; on the
fundamental domain $[-\pi,\pi)^3$ exactly the term $m=0$ is nonzero, so
$\operatorname{supp}\psi_N\subset\{|x|\le R/N\}$ there; and
\[
\|\psi_N\|_{L^2(\T^3)}=(2\pi)^{-3/2}\|\Psi\|_{L^2(\R^3)},
\]
which is independent of $N$.
\end{lemma}

\begin{proof}
Local finiteness: $\Psi(N(x+2\pi m))\neq0$ forces
$|x+2\pi m|\le R/N<\pi$; but for $x\in[-\pi,\pi)^3$ and $m\neq0$,
$|x+2\pi m|\ge2\pi|m|_\infty-\pi\ge\pi$, a contradiction. Hence on a
neighbourhood of any point the sum has at most one nonzero term, so
$\psi_N$ is smooth, real, and $2\pi\Z^3$-periodic by construction.
Divergence-free: locally
$\nabla\cdot\psi_N(x)=N^{5/2}(\nabla\cdot\Psi)(N(x+2\pi m))=0$.
Zero mean: with $d\mu=(2\pi)^{-3}dx$ and the substitution $y=Nx$,
\[
\int_{\T^3}\psi_N\,d\mu
=(2\pi)^{-3}N^{3/2}\int_{\R^3}\Psi(Nx)\,dx
=(2\pi)^{-3}N^{-3/2}\int_{\R^3}\Psi ,
\]
and $\int_{\R^3}\Psi_i=\int_{\R^3}\nabla\cdot(y_i\Psi)=0$ since
$\nabla\cdot\Psi=0$ and $y_i\Psi\in C^\infty_c$. (No curl
representation of $\Psi$ is needed: zero mean is automatic for a
compactly supported divergence-free field.) Norm:
\[
\|\psi_N\|_{L^2(\T^3)}^2
=(2\pi)^{-3}\int_{|x|\le R/N}N^{3}|\Psi(Nx)|^2dx
=(2\pi)^{-3}N^{3}N^{-3}\|\Psi\|_{L^2(\R^3)}^2 ,
\]
which is the stated identity.
\end{proof}

\begin{theorem}[duality identity with controlled error]\label{app:pairing}
Let $\Psi$ be as in Definition~\ref{app:testfield} and put
\[
I_\Psi:=\int_{\R^3}\big(V\cdot\nabla V\big)\cdot\Psi\,dy .
\]
Then for every integer $N\ge\max(8R,32)$,
\[
\big\langle\PP(u_N\cdot\nabla u_N),\psi_N\big\rangle_{L^2(\T^3)}
=(2\pi)^{-3}N^{3/2}\big(I_\Psi+\mathcal E_N\big),
\]
where
\[
|\mathcal E_N|\ \le\ \|\Psi\|_{L^1(\R^3)}
\Big[\varepsilon_N\big(6\pi\|v_0\|+\varepsilon'_N\big)
+4\pi\|v_0\|\,\varepsilon'_N\Big]
\ =\ O\Big(\frac{\log N}{N}\Big)
\]
with $\varepsilon_N,\varepsilon'_N$ as in Theorem~\ref{app:rate}.
\end{theorem}

\begin{proof}
By Lemma~\ref{app:psiprops}, $\psi_N$ is zero-mean and divergence-free,
so $\hat\psi_N(0)=0$ and $k\cdot\hat\psi_N(k)=0$ for all $k$, whence
$\PP\psi_N=\psi_N$. Since $\PP$ is an orthogonal projection on
$L^2(\T^3)$, hence self-adjoint,
\[
\big\langle\PP(u_N\cdot\nabla u_N),\psi_N\big\rangle
=\big\langle u_N\cdot\nabla u_N,\PP\psi_N\big\rangle
=\big\langle u_N\cdot\nabla u_N,\psi_N\big\rangle .
\]
Both fields are real. Because $N\ge8R$ and $R\ge1$ we have
$R/N\le\tfrac18<\pi$, so Lemma~\ref{app:psiprops} applies and
$\psi_N(x)=N^{3/2}\Psi(Nx)$ on the fundamental domain, supported in
$\{|x|\le R/N\}$. Using $d\mu=(2\pi)^{-3}dx$ and substituting $x=y/N$
($dx=N^{-3}dy$),
\[
\big\langle u_N\cdot\nabla u_N,\psi_N\big\rangle
=(2\pi)^{-3}N^{3/2}N^{-3}\int_{|y|\le R}
(u_N\cdot\nabla u_N)(y/N)\cdot\Psi(y)\,dy .
\]
By Lemma~\ref{app:innerexact}, for $|y|\le R$,
\[
(u_N\cdot\nabla u_N)_j(y/N)
=u_{N,i}(y/N)\,\partial_iu_{N,j}(y/N)
=N^{3}\big(V_i(y)+E_{N,i}(y)\big)\big(\partial_iV_j(y)+E'_{N,ij}(y)\big),
\]
so the factor $N^{-3}N^{3}$ cancels and
\[
\big\langle u_N\cdot\nabla u_N,\psi_N\big\rangle
=(2\pi)^{-3}N^{3/2}\big(I_\Psi+\mathcal E_N\big),
\qquad
\mathcal E_N=\int_{|y|\le R}
\big[E_{N,i}\partial_iV_j+V_iE'_{N,ij}+E_{N,i}E'_{N,ij}\big]\Psi_j\,dy ,
\]
where we used $\operatorname{supp}\Psi\subset\{|y|\le R\}$ to write the
main term as $I_\Psi$. Bounding each factor by its supremum and using
Theorem~\ref{app:Vprops}(2),(3) --- $\|V\|_\infty\le4\pi\|v_0\|$,
$\|\nabla V\|_\infty\le6\pi\|v_0\|$ --- gives the stated estimate;
$\varepsilon_N,\varepsilon'_N=O(N^{-1}\log N)$ by
Theorem~\ref{app:rate}, $R$ and $\Psi$ being fixed.
\end{proof}

\begin{proposition}[the $N^3$ lower bound, given a good test field]
\label{app:N3}
Suppose some $\Psi$ as in Definition~\ref{app:testfield} has
$I_\Psi\neq0$. Then
\[
\big\|\PP(u_N\cdot\nabla u_N)\big\|_{L^2(\T^3)}^2
\ \ge\ \frac{|I_\Psi+\mathcal E_N|^2}
{(2\pi)^{3}\,\|\Psi\|_{L^2(\R^3)}^{2}}\;N^{3}
\qquad(N\ge\max(8R,32)),
\]
and consequently
$\|\PP(u_N\cdot\nabla u_N)\|_{2}^2\ge c_0N^3$ for all $N\ge N_*$, where
\[
c_0=\frac{|I_\Psi|^2}{2\,(2\pi)^{3}\,\|\Psi\|_{L^2(\R^3)}^{2}},
\qquad
N_*=\min\Big\{N_0\ge\max(8R,32):\ |\mathcal E_N|\le(1-2^{-1/2})|I_\Psi|
\ \ \forall N\ge N_0\Big\},
\]
which is finite because $\mathcal E_N\to0$.
\end{proposition}

\begin{proof}
By Cauchy--Schwarz in $L^2(\T^3)$,
\[
\big\|\PP(u_N\cdot\nabla u_N)\big\|_{2}
\ \ge\ \frac{\big|\langle\PP(u_N\cdot\nabla u_N),\psi_N\rangle\big|}
{\|\psi_N\|_{2}} .
\]
Insert Theorem~\ref{app:pairing} and Lemma~\ref{app:psiprops}: the
right-hand side equals
\[
\frac{(2\pi)^{-3}N^{3/2}|I_\Psi+\mathcal E_N|}
{(2\pi)^{-3/2}\|\Psi\|_{L^2(\R^3)}}
=(2\pi)^{-3/2}N^{3/2}\frac{|I_\Psi+\mathcal E_N|}{\|\Psi\|_{L^2(\R^3)}} .
\]
Squaring gives the display. For $N\ge N_*$ we have
$|I_\Psi+\mathcal E_N|\ge2^{-1/2}|I_\Psi|$, whence
$|I_\Psi+\mathcal E_N|^2\ge\tfrac12|I_\Psi|^2$ and the constant $c_0$ as
stated.
\end{proof}

Everything is now reduced to the single question of whether some
admissible $\Psi$ has $I_\Psi\neq0$.

\subsection{Non-degeneracy of the continuum nonlinearity}
\label{app:nondeg}

\subsubsection*{Reduction to a symbol statement}

\begin{lemma}[symbol of the nonlinearity]\label{app:symbol}
$V\cdot\nabla V=\nabla\cdot(V\otimes V)$ has the representation
\[
(V\cdot\nabla V)(y)=\int_{\R^3}M(\zeta)e^{i\zeta\cdot y}\,d\zeta,
\qquad
M_i(\zeta)=i\,\zeta_j\,T_{ij}(\zeta),\qquad T_{ij}:=F_i*F_j ,
\]
where $T$ is real, symmetric, supported in $\{|\zeta|\le2\}$, and
$T,M\in L^1(\R^3)\cap L^2(\R^3)$. Consequently
$\PP(V\cdot\nabla V)\in L^2(\R^3)$ with symbol $P_\zeta M(\zeta)$, and
\[
\PP(V\cdot\nabla V)\equiv0
\quad\Longleftrightarrow\quad
\zeta\times\big(T(\zeta)\zeta\big)=0\ \text{ for a.e. }\zeta .
\]
\end{lemma}

\begin{proof}
$V=\widetilde F$ with $F\in L^1$, and $\partial_jV_i=\widetilde{i\xi_jF_i}$
with $\xi F\in L^1$ (Theorem~\ref{app:Vprops}(1),(3)). The pointwise
product of inverse transforms of $L^1$ functions is the inverse
transform of the convolution, so
$V_j\partial_jV_i=\widetilde{F_j*(i\,\cdot_j F_i)}$, i.e.\ the symbol of
$V\cdot\nabla V$ is
$i\int F_j(\zeta-\eta)\,\eta_j\,F_i(\eta)\,d\eta$. Since
$(\zeta-\eta)\cdot F(\zeta-\eta)=0$ identically (Definition~\ref{app:Vdef}),
we may replace $\eta_jF_j(\zeta-\eta)$ by $\zeta_jF_j(\zeta-\eta)$
inside the integral, obtaining
$M_i(\zeta)=i\zeta_j(F_i*F_j)(\zeta)=i\zeta_jT_{ij}(\zeta)$.

$T$ is real because $F$ is real, symmetric because convolution is
commutative, and supported in $\{|\zeta|\le2\}$ because
$\operatorname{supp}F\subset\{|\xi|\le1\}$. Integrability: by
Theorem~\ref{app:Vprops}(1), $F\in L^{4/3}$; Young's inequality with
$\frac1r=\frac2{4/3}-1=\frac12$ gives $T=F*F\in L^{2}$, and since $T$
has compact support, $T\in L^1\cap L^2$; the same holds for
$M=i\zeta T$, as $|\zeta|\le2$ on the support. By Plancherel,
$V\cdot\nabla V\in L^2(\R^3)$ and, since $|P_\zeta M|\le|M|$,
$\PP(V\cdot\nabla V)\in L^2(\R^3)$ with symbol $P_\zeta M$. Finally
$\PP(V\cdot\nabla V)=0$ in $L^2$ iff $P_\zeta M(\zeta)=0$ for a.e.
$\zeta$, and for a vector $w$ one has $P_\zeta w=0$ iff $w\parallel\zeta$
iff $\zeta\times w=0$; here $w=i\,T(\zeta)\zeta$.
\end{proof}

\begin{remark}[classical input: density of smooth compactly supported
solenoidal fields]\label{app:classical-density}
The space $C^\infty_{c,\sigma}(\R^3;\R^3)$ of real, smooth, compactly
supported, divergence-free vector fields is dense in
$L^2_\sigma(\R^3;\R^3)$, the closed subspace of real $L^2$ fields that
are distributionally divergence-free; indeed $L^2_\sigma$ is by the
standard definition the $L^2$-closure of $C^\infty_{c,\sigma}$, and on
$\R^3$ the two descriptions agree. See, e.g.,
\cite{Temam,LadyzhenskayaBook,ConstantinFoias}.
\end{remark}

\begin{remark}[classical input: Helmholtz--Leray decomposition on $\R^3$]
\label{app:classical-leray}
For $f\in L^2(\R^3;\R^3)$ one has $f-\PP f=\nabla p$ for some
$p\in L^2_{\rm loc}(\R^3)$ with $\nabla p\in L^2$; equivalently, the
symbol of $f-\PP f$ is $\zeta(\zeta\cdot\hat f)/|\zeta|^2$. Same
references.
\end{remark}

\begin{lemma}[from non-degeneracy to a test field]\label{app:testexists}
If $\PP(V\cdot\nabla V)\not\equiv0$, then there exists a real
$\Psi\in C^\infty_c(\R^3;\R^3)$ with $\nabla\cdot\Psi=0$ and
$I_\Psi=\int_{\R^3}(V\cdot\nabla V)\cdot\Psi\neq0$; after enlarging $R$
if necessary, $\Psi$ satisfies Definition~\ref{app:testfield}.
\end{lemma}

\begin{proof}
By Lemma~\ref{app:symbol}, $\PP(V\cdot\nabla V)$ lies in $L^2(\R^3)$,
is real (its symbol $P_\zeta M$ satisfies the reality condition because
$F$ is real and even), is distributionally divergence-free (its symbol
is orthogonal to $\zeta$), and is nonzero by hypothesis. Hence it is a
nonzero element of the real Hilbert space $L^2_\sigma(\R^3)$. By
Remark~\ref{app:classical-density} there is a real
$\Psi\in C^\infty_{c,\sigma}(\R^3)$ with
$\langle\PP(V\cdot\nabla V),\Psi\rangle_{L^2(\R^3)}\neq0$. By
Remark~\ref{app:classical-leray},
$V\cdot\nabla V-\PP(V\cdot\nabla V)=\nabla p$ with $\nabla p\in L^2$, and
$\int\nabla p\cdot\Psi=-\int p\,\nabla\cdot\Psi=0$ because $\Psi$ is
compactly supported and divergence-free. Therefore
$I_\Psi=\langle\PP(V\cdot\nabla V),\Psi\rangle\neq0$. Choosing
$R\ge\max(1,\sup\{|y|:y\in\operatorname{supp}\Psi\})$ puts $\Psi$ in
Definition~\ref{app:testfield}.
\end{proof}

\subsubsection*{The $\chi$-independent leading singularity}

By rotational equivariance we may and do assume from now on that
$v_0=\|v_0\|e_3$: $F$ and $V$ are linear in $v_0$ and rotate with it,
$T$ and $\tau$ are quadratic, and the conclusion
$\PP(V\cdot\nabla V)\not\equiv0$ is rotation invariant.

\begin{definition}\label{app:tau}
$\tau_{ij}:=h_i*h_j$, with $h(\xi)=P_\xi v_0/|\xi|^2$ the uncut symbol
of Definition~\ref{app:Vdef}.
\end{definition}

\begin{lemma}[$\tau$ is well defined and homogeneous]\label{app:tauwd}
The integral $\tau_{ij}(\zeta)=\int_{\R^3}h_i(\eta)h_j(\zeta-\eta)\,d\eta$
converges absolutely for every $\zeta\neq0$, and $\tau$ is homogeneous of
degree $-1$ and smooth on $\R^3\setminus\{0\}$.
\end{lemma}

\begin{proof}
$|h(\xi)|\le\|v_0\||\xi|^{-2}$. Near $\eta=0$ the integrand is
$O(|\eta|^{-2})$ and near $\eta=\zeta$ it is $O(|\zeta-\eta|^{-2})$, both
locally integrable on $\R^3$; for $|\eta|\ge2|\zeta|$ it is
$O(|\eta|^{-4})$, integrable at infinity. Substituting
$\eta=\lambda\eta'$ and using $h(\lambda\eta)=\lambda^{-2}h(\eta)$ gives
$\tau(\lambda\zeta)=\lambda^{-1}\tau(\zeta)$ for $\lambda>0$. Smoothness
away from $0$ follows from differentiating under the integral sign on
the region where the two singularities are separated.
\end{proof}

\begin{theorem}[closed form of the leading obstruction]
\label{app:tauformula}
With $v_0=\|v_0\|e_3$, for every $\zeta\neq0$,
\[
\zeta\times\big(\tau(\zeta)\zeta\big)
=\frac{3\pi^3}{8}\,\|v_0\|^2\,\frac{\zeta_3}{|\zeta|}\,
\big(\zeta\times e_3\big).
\]
In particular $\zeta\times(\tau(\zeta)\zeta)\neq0$ whenever
$\zeta_3\neq0$ and $\zeta\not\parallel e_3$.
\end{theorem}

\begin{proof}
Both sides are quadratic in $\|v_0\|$, so we normalise $\|v_0\|=1$,
$v_0=e_3$. By Theorem~\ref{app:farfield}, $\widetilde h=\pi^2W$ with
\[
W_i(y)=\frac{\delta_{i3}}{r}+\frac{y_3y_i}{r^3},\qquad r=|y| .
\]
$W_i$ is locally integrable, homogeneous of degree $-1$ and smooth away
from the origin, so Remark~\ref{app:classical-homog} applies with
$a=b=1$ and gives
\[
\tau_{ij}=h_i*h_j=(2\pi)^{-3}\,\mathcal F\big[\widetilde h_i\widetilde h_j\big]
=(2\pi)^{-3}\pi^4\,\mathcal F\big[W_iW_j\big],
\]
where we used $\mathcal F\widetilde h_i=(2\pi)^3h_i$ and the evenness of
$h$. (Both sides are homogeneous of degree $-1$ and locally integrable,
so no distribution supported at the origin can be added.) Expanding,
\[
W_iW_j=\frac{\delta_{i3}\delta_{j3}}{r^{2}}
+\frac{\delta_{i3}y_3y_j+\delta_{j3}y_3y_i}{r^{4}}
+\frac{y_3^{2}y_iy_j}{r^{6}} .
\]
Each bracket is locally integrable and homogeneous of degree $-2$, so by
Remark~\ref{app:classical-riesz} its transform is an honest homogeneous
distribution of degree $-1$, given by
\[
\mathcal F\Big[\frac1{r^{2}}\Big]=\frac{2\pi^{2}}{s},
\qquad
\mathcal F\Big[\frac{y_iy_j}{r^{4}}\Big]
=\pi^{2}\Big(\frac{\delta_{ij}}{s}-\frac{\zeta_i\zeta_j}{s^{3}}\Big),
\qquad
\mathcal F\Big[\frac{y_3^{2}y_iy_j}{r^{6}}\Big]
=\frac{\pi^{2}}{12}\,\partial_3^{2}\partial_i\partial_j\,s^{3},
\]
where $s:=|\zeta|$. (Indeed $\mathcal F[r^{-2}]=c_{3,2}s^{-1}$;
$\mathcal F[y_iy_jr^{-4}]=-\partial_i\partial_j(c_{3,4}s)
=\pi^2\partial_i\partial_js=\pi^2(\delta_{ij}/s-\zeta_i\zeta_j/s^3)$;
$\mathcal F[y_3^2y_iy_jr^{-6}]
=\partial_3^2\partial_i\partial_j(c_{3,6}s^3)$. Trace check on the
middle formula: $\pi^2(3/s-1/s)=2\pi^2/s=\mathcal F[r^{-2}]$, as it must
be.) Direct differentiation of $s^3$ gives
\[
\partial_3^{2}\partial_i\partial_j s^{3}
=3\Big[\frac{\delta_{ij}+2\delta_{i3}\delta_{j3}}{s}
-\frac{\zeta_3^{2}\delta_{ij}
+2\zeta_3(\delta_{i3}\zeta_j+\delta_{j3}\zeta_i)+\zeta_i\zeta_j}{s^{3}}
+\frac{3\zeta_i\zeta_j\zeta_3^{2}}{s^{5}}\Big].
\]
Assembling the three pieces (the middle one contributing
$\pi^2[2\delta_{i3}\delta_{j3}/s-\zeta_3(\delta_{i3}\zeta_j
+\delta_{j3}\zeta_i)/s^3]$) yields
$\mathcal F[W_iW_j]=\pi^{2}G_{ij}$ with
\[
G_{ij}=\frac{\delta_{ij}}{4s}
+\frac{9\,\delta_{i3}\delta_{j3}}{2s}
-\frac{\zeta_3^{2}\delta_{ij}}{4s^{3}}
-\frac{3\zeta_3(\delta_{i3}\zeta_j+\delta_{j3}\zeta_i)}{2s^{3}}
-\frac{\zeta_i\zeta_j}{4s^{3}}
+\frac{3\zeta_i\zeta_j\zeta_3^{2}}{4s^{5}} ,
\]
so that $\tau_{ij}=(2\pi)^{-3}\pi^{4}\cdot\pi^{2}G_{ij}
=\frac{\pi^{3}}{8}G_{ij}$.

Contract with $\zeta_j$, using $\zeta_j\zeta_j=s^2$. The terms
$\frac{\delta_{ij}}{4s}\zeta_j=\frac{\zeta_i}{4s}$ and
$-\frac{\zeta_i\zeta_j}{4s^3}\zeta_j=-\frac{\zeta_i}{4s}$ cancel. The
$\delta_{i3}\zeta_3/s$ coefficients are $\frac92$ from the second term
and $-\frac32$ from the fourth, total $3$. The $\zeta_i\zeta_3^2/s^3$
coefficients are $-\frac14$ (third term), $-\frac32$ (fourth term) and
$+\frac34$ (sixth term), total $-1$. Hence
\[
G(\zeta)\zeta=\frac{3\zeta_3}{s}\,e_3-\frac{\zeta_3^{2}}{s^{3}}\,\zeta ,
\qquad\text{so}\qquad
\zeta\times\big(G(\zeta)\zeta\big)=\frac{3\zeta_3}{s}\,(\zeta\times e_3),
\]
the second term being parallel to $\zeta$. Multiplying by
$\frac{\pi^3}{8}$ and restoring $\|v_0\|^2$ gives the claim. Finally
$\zeta\times e_3=0$ exactly when $\zeta\parallel e_3$.
\end{proof}

\begin{remark}[physical-space cross-check, free of any regularisation]
\label{app:crosscheck}
Theorem~\ref{app:tauformula} passes through the finite-part
continuations of Remark~\ref{app:classical-riesz}. The following
elementary computation confirms it independently, with no
regularisation and no Fourier transform, and is worth recording because
it identifies what the obstruction \emph{is}: the statement that the
Oseen field is not a stationary Euler flow.

For $W(y)=\frac{e_3}{r}+\frac{y_3y}{r^{3}}$ on $\R^3\setminus\{0\}$,
\[
\nabla\cdot W=0,\qquad
\omega:=\nabla\times W=\frac{2\,(e_3\times y)}{r^{3}},\qquad
\nabla\times\big(W\times\omega\big)
=\frac{12\,y_3\,(e_3\times y)}{r^{6}}\ \not\equiv0 .
\]
\emph{Proof.} $\nabla\cdot(e_3/r)=e_3\cdot\nabla r^{-1}=-y_3/r^3$, and
$\partial_i(y_3y_ir^{-3})=y_3r^{-3}+3y_3r^{-3}-3y_3r^{-3}=y_3/r^3$, so
$\nabla\cdot W=0$. Next
$\nabla\times(r^{-1}e_3)=\nabla(r^{-1})\times e_3=-r^{-3}y\times e_3
=r^{-3}(e_3\times y)$, while
$[\nabla\times(y_3y\,r^{-3})]_a
=\epsilon_{abi}\partial_b(y_3y_ir^{-3})
=\epsilon_{abi}(\delta_{b3}y_i+y_3\delta_{bi})r^{-3}
-3\epsilon_{abi}y_3y_iy_br^{-5}
=\epsilon_{a3i}y_ir^{-3}=(e_3\times y)_ar^{-3}$, the last two terms
vanishing by antisymmetry. Adding gives $\omega$. Then, using
$e_3\times(e_3\times y)=y_3e_3-y$ and $y\times(e_3\times y)=r^2e_3-y_3y$,
\[
W\times\omega
=\frac{2(y_3e_3-y)}{r^{4}}+\frac{2y_3e_3}{r^{4}}-\frac{2y_3^{2}y}{r^{6}}
=\frac{4y_3e_3}{r^{4}}-\frac{2y}{r^{4}}-\frac{2y_3^{2}y}{r^{6}} .
\]
Now curl each piece. For any radial $f$, $\nabla\times(f(r)y)
=\nabla f\times y=f'(r)\hat y\times y=0$, so $-2y/r^4$ contributes
nothing. With $g=-2y_3^{2}/r^{6}$ one has
$\nabla g=-4y_3e_3/r^{6}+12y_3^{2}y/r^{8}$, so
$\nabla\times(gy)=\nabla g\times y=-4y_3(e_3\times y)/r^{6}$. With
$f=4y_3/r^{4}$ one has $\nabla f=4e_3/r^{4}-16y_3y/r^{6}$, so
$\nabla\times(fe_3)=\nabla f\times e_3
=-16y_3(y\times e_3)/r^{6}=16y_3(e_3\times y)/r^{6}$. The total is
$12y_3(e_3\times y)/r^{6}$. Since
$W\cdot\nabla W=\nabla\frac{|W|^2}{2}-W\times\omega$, it follows that
$\nabla\times(W\cdot\nabla W)\not\equiv0$, i.e.\ $W\cdot\nabla W$ is not
a gradient, i.e.\ $\PP(W\cdot\nabla W)\neq0$. $\square$

\emph{Consistency with Theorem~\ref{app:tauformula}.} The curl of
$\PP(\widetilde h\cdot\nabla\widetilde h)$ has symbol
$i\zeta\times\big(iP_\zeta\tau\zeta\big)=-\zeta\times(\tau\zeta)$, using
$\zeta\times P_\zeta A=\zeta\times A$. With $\|v_0\|=1$,
Theorem~\ref{app:tauformula} makes this
$-\frac{3\pi^{3}}{8}\zeta_3(\zeta\times e_3)/|\zeta|$. Writing its
$a$-component as $-\frac{3\pi^3}{8}\epsilon_{ab3}\zeta_b\zeta_3/|\zeta|$
and using $\widetilde{|\zeta|^{-1}}=c_{3,1}r^{-2}=4\pi r^{-2}$ together
with $\widetilde{\zeta_b\zeta_3G}=-\partial_b\partial_3\widetilde G$,
\[
\widetilde{\Big[\frac{\zeta_b\zeta_3}{|\zeta|}\Big]}
=-4\pi\,\partial_b\partial_3r^{-2}
=\frac{8\pi\delta_{b3}}{r^{4}}-\frac{32\pi y_3y_b}{r^{6}} ;
\]
since $\epsilon_{ab3}\delta_{b3}=0$, only the second term survives,
leaving $-32\pi y_3(y\times e_3)_a/r^{6}
=+32\pi y_3(e_3\times y)_a/r^{6}$. Multiplying by
$-\frac{3\pi^{3}}{8}$ gives $-12\pi^{4}y_3(e_3\times y)/r^{6}$. On the
other hand $\widetilde h=\pi^2W$ gives
$\widetilde h\cdot\nabla\widetilde h=\pi^{4}W\cdot\nabla W$ and hence
$\nabla\times(\widetilde h\cdot\nabla\widetilde h)
=-\pi^{4}\nabla\times(W\times\omega)
=-12\pi^{4}y_3(e_3\times y)/r^{6}$. The two derivations agree exactly,
constants included.
\end{remark}

\subsubsection*{Transfer to the cutoff family}

\begin{lemma}[the cutoff perturbation is bounded near $\zeta=0$]
\label{app:perturb}
For all $|\zeta|\le\frac18$ and all $i,j$,
\[
\big|T_{ij}(\zeta)-\tau_{ij}(\zeta)\big|\ \le\ C_1:=64\pi\|v_0\|^{2},
\]
uniformly over all admissible $\chi$.
\end{lemma}

\begin{proof}
Since $F=\chi(|\cdot|)h$,
\[
T_{ij}(\zeta)-\tau_{ij}(\zeta)
=\int_{\R^3}h_i(\eta)h_j(\zeta-\eta)
\big[\chi(|\eta|)\chi(|\zeta-\eta|)-1\big]\,d\eta .
\]
The bracket vanishes whenever $|\eta|\le\frac12$ \emph{and}
$|\zeta-\eta|\le\frac12$, because $\chi\equiv1$ on $[0,\frac12]$. For
$|\zeta|\le\frac18$, the condition $|\eta|\le\frac14$ implies
$|\zeta-\eta|\le\frac18+\frac14=\frac38<\frac12$; hence the integrand is
supported in $\{|\eta|\ge\frac14\}$. There
$|\zeta-\eta|\ge|\eta|-\frac18\ge|\eta|-\frac{|\eta|}{2}
=\frac{|\eta|}{2}$. With $|h|\le\|v_0\||\cdot|^{-2}$ and
$|\chi\chi-1|\le1$,
\[
|T_{ij}-\tau_{ij}|
\le\|v_0\|^{2}\int_{|\eta|\ge1/4}\frac{d\eta}{|\eta|^{2}|\zeta-\eta|^{2}}
\le4\|v_0\|^{2}\int_{|\eta|\ge1/4}\frac{d\eta}{|\eta|^{4}}
=4\|v_0\|^{2}\cdot4\pi\int_{1/4}^{\infty}\frac{dr}{r^{2}}
=64\pi\|v_0\|^{2},
\]
which is $C_1$.
\end{proof}

\begin{theorem}[non-degeneracy]\label{app:vnondeg}
For every $v_0\in\R^3\setminus\{0\}$ and every admissible $\chi$,
\[
\PP_{\R^3}\big(V\cdot\nabla V\big)\ \not\equiv\ 0 .
\]
Explicitly, with $v_0=\|v_0\|e_3$, one has
$\zeta\times(T(\zeta)\zeta)\neq0$ --- hence $P_\zeta M(\zeta)\neq0$ ---
at every point of the open set of positive measure
\[
\mathcal U:=\Big\{\zeta\in\R^3:\ 0<|\zeta|<\tfrac{\pi^{2}}{2048},\
|\zeta_3|\ge\tfrac{|\zeta|}{2},\
|\zeta\times e_3|\ge\tfrac{|\zeta|}{2}\Big\}.
\]
\end{theorem}

\begin{proof}
Let $\zeta\in\mathcal U$. By Theorem~\ref{app:tauformula},
\[
\big|\zeta\times(\tau(\zeta)\zeta)\big|
=\frac{3\pi^{3}}{8}\|v_0\|^{2}\,
\frac{|\zeta_3|\,|\zeta\times e_3|}{|\zeta|}
\ \ge\ \frac{3\pi^{3}}{8}\|v_0\|^{2}\cdot\frac{|\zeta|}{4}
=\frac{3\pi^{3}}{32}\|v_0\|^{2}|\zeta| .
\]
Since $\frac{\pi^{2}}{2048}\approx4.8\times10^{-3}<\frac18$,
Lemma~\ref{app:perturb} applies, and the Frobenius norm of
$T(\zeta)-\tau(\zeta)$ is at most $3C_1$ (nine entries each at most
$C_1$). Hence
\[
\big|\zeta\times\big((T(\zeta)-\tau(\zeta))\zeta\big)\big|
\le|\zeta|\cdot3C_1|\zeta|=192\pi\|v_0\|^{2}|\zeta|^{2},
\]
and therefore
\[
\big|\zeta\times(T(\zeta)\zeta)\big|
\ \ge\ \|v_0\|^{2}|\zeta|
\Big(\frac{3\pi^{3}}{32}-192\pi|\zeta|\Big)\ >\ 0
\qquad\text{whenever }
|\zeta|<\frac{3\pi^{3}}{32\cdot192\pi}=\frac{3\pi^{2}}{6144}
=\frac{\pi^{2}}{2048},
\]
which holds on $\mathcal U$. The set $\mathcal U$ is open and nonempty
(it contains, e.g., all small $\zeta$ in the direction
$(1,0,1)/\sqrt2$), hence of positive measure. By
Lemma~\ref{app:symbol}, $P_\zeta M\neq0$ on a set of positive measure,
so $\PP(V\cdot\nabla V)\neq0$ in $L^2(\R^3)$.
\end{proof}

\begin{remark}[uniformity, and where the obstruction lives]
\label{app:uniformity}
The proof of Theorem~\ref{app:vnondeg} is uniform in $\chi$ and in
$v_0$: the factor $\|v_0\|^2$ cancels between the main term and the
perturbation, and the only properties of $\chi$ used are $0\le\chi\le1$
and $\chi\equiv1$ on $[0,\frac12]$ --- no smoothness of $\chi$ enters
\S\ref{app:nondeg} at all. Neither a specific $\chi$ nor any numerically
evaluated convolution integral is required: the obstruction is carried
entirely by the $\chi$-independent leading singularity of $F*F$ at
$\zeta=0$, which is the same for every admissible cutoff. Moreover
Theorem~\ref{app:tauformula} shows the obstruction degenerates exactly
on the two directions $\zeta\parallel v_0$ and $\zeta\perp v_0$, so the
two cone conditions in $\mathcal U$ are essentially optimal.

This uniformity does \textbf{not} propagate to the constants of
Theorem~\ref{app:main}. It is the \emph{qualitative} conclusion
$\PP(V\cdot\nabla V)\not\equiv0$ that holds uniformly in $(\chi,v_0)$,
and that is what licenses the quantifier ``for every admissible $\chi$
and every $v_0\neq0$'' there. The quantitative constants $c_0$ and
$N_*$ are extracted only after a test field $\Psi$ has been chosen for
the particular profile $V=V^{\chi,v_0}$
(Lemma~\ref{app:testexists}), and they depend on $(\chi,v_0)$; no
$(\chi,v_0)$-independent value of either is claimed, and neither is
effective (Remark~\ref{app:noneffective}).
\end{remark}

\subsection{Proof of Theorem~\ref{capacity}}
\label{app:final}

\begin{theorem}[capacity bound; restatement of Theorem~\ref{capacity}]
\label{app:main}
Let $v_0\in\R^3\setminus\{0\}$, let $\chi$ be admissible
(Definition~\ref{app:admissible}), and let $u_N=u_N^{\chi,v_0}$ be the
family \eqref{app:famdef}. Then:
\begin{enumerate}
\item[(1)] $u_N$ is a real, zero-mean, exactly divergence-free
trigonometric polynomial banded at $|k|\le N$, and the exact laws of
Theorem~\ref{app:exactlaws} hold;
\item[(2)] $\dfrac{N}{44544}\le N_0^2(u_N)\le\dfrac{64}{3}N$ for every
$N\ge32$;
\item[(3)] there exist $c_0=c_0(\chi,v_0)>0$ and
$N_*=N_*(\chi,v_0)<\infty$ such that
\[
\big\|\PP(u_N\cdot\nabla u_N)\big\|_{L^2(\T^3)}^{2}\ \ge\ c_0\,N^{3}
\qquad\text{for all }N\ge N_* .
\]
\end{enumerate}
\end{theorem}

\begin{proof}
(1) is Lemma~\ref{app:structure} together with
Theorem~\ref{app:exactlaws}. (2) is Proposition~\ref{app:twosided}.

(3) By Theorem~\ref{app:vnondeg}, $\PP_{\R^3}(V\cdot\nabla V)\not\equiv0$,
where $V$ is the profile of Definition~\ref{app:Vdef} attached to the
given $(\chi,v_0)$. By Lemma~\ref{app:testexists} there is a real
$\Psi\in C^\infty_c(\R^3;\R^3)$, divergence-free, supported in
$\{|y|\le R\}$ for some $R\ge1$, with $I_\Psi\neq0$. Fix such a $\Psi$.
Lemma~\ref{app:innerexact} and Theorem~\ref{app:rate} give the exact
inner rescaling with $\varepsilon_N+\varepsilon'_N=O(N^{-1}\log N)$;
Lemma~\ref{app:psiprops} and Theorem~\ref{app:pairing} convert the
pairing of $\PP(u_N\cdot\nabla u_N)$ against the concentrated field
$\psi_N$ into $(2\pi)^{-3}N^{3/2}(I_\Psi+\mathcal E_N)$ with
$\mathcal E_N\to0$; and Proposition~\ref{app:N3} then yields the stated
bound with
\[
c_0=c_0(\chi,v_0)=\frac{|I_\Psi|^{2}}
{2\,(2\pi)^{3}\,\|\Psi\|_{L^2(\R^3)}^{2}},
\qquad
N_*=N_*(\chi,v_0)
=\min\Big\{N_0\ge\max(8R,32):|\mathcal E_N|\le(1-2^{-1/2})|I_\Psi|
\ \forall N\ge N_0\Big\}.
\]
Both depend on $(\chi,v_0)$ alone, through $V=V^{\chi,v_0}$ and through
the test field $\Psi$ selected for it.
\end{proof}

\begin{remark}[what the constants depend on, and that they are
non-effective]\label{app:noneffective}
The constants $c_0$ and $N_*$ produced above depend on the pair
$(\chi,v_0)$ only, through the profile $V=V^{\chi,v_0}$ and through the
test field $\Psi$ selected in Lemma~\ref{app:testexists}. They are
\textbf{not effective}. The reason is localised precisely at
Lemma~\ref{app:testexists}: $\Psi$ is obtained from the \emph{density}
of $C^\infty_{c,\sigma}(\R^3)$ in $L^2_\sigma(\R^3)$
(Remark~\ref{app:classical-density}), a pure existence argument, so
neither $\Psi$, nor its support radius $R$, nor $I_\Psi$, nor
$\|\Psi\|_{L^2}$, nor $c_0$, nor $N_*$ is exhibited. Every other
constant in this appendix is explicit, namely those of
Lemmas~\ref{app:shell}--\ref{app:lower},
Proposition~\ref{app:twosided}, Theorems~\ref{app:exactlaws},
\ref{app:Vprops}, \ref{app:farfield}, \ref{app:rate},
\ref{app:pairing}, \ref{app:tauformula}, \ref{app:vnondeg},
Lemmas~\ref{app:innerexact}, \ref{app:psiprops}, \ref{app:symbol},
\ref{app:tauwd}, \ref{app:perturb}, and Proposition~\ref{app:N3}
\emph{as a function of} $(\Psi,R,I_\Psi)$.

An effective variant is available in principle --- take $\Psi$ with
$\widehat\Psi$ a smooth bump supported in the explicit set $\mathcal U$
of Theorem~\ref{app:vnondeg}, mollified to compact support in $y$ ---
but it is expensive: non-vanishing is certified only for
$|\zeta|<\pi^2/2048\approx4.8\times10^{-3}$, so such a $\Psi$ has
Fourier support at that scale, hence spatial radius $R\gtrsim10^{3}$,
hence $N_*\ge8R\gtrsim10^{4}$, with a correspondingly small $c_0$. This
is the sole source of non-effectivity in the paper.
\end{remark}

\begin{remark}[the exponent $3$ is sharp, and the sharp cutoff is not
covered]\label{app:sharpness}
The exponent cannot be improved: by Theorem~\ref{main}(d) and
Theorem~\ref{app:exactlaws},
$\|\PP(u_N\cdot\nabla u_N)\|_2^2\le\|u_N\|_\infty^2H_1(u_N)$, and both
$\|u_N\|_\infty$ and $H_1(u_N)$ are $O(N)$ by
Proposition~\ref{app:twosided} and
$\|u_N\|_\infty\le\sum_{k\neq0}|\hat u_N(k)|\le\|v_0\|\Sigma_N^\chi
\le\|v_0\|S_N\le128\|v_0\|N$; so the left-hand side is $O(N^3)$.

Smoothness of $\chi$ is used in exactly one place, namely
Theorem~\ref{app:rate}, through $|\nabla F|\lesssim c_\chi|\xi|^{-3}$.
For the sharply truncated family $\chi=\mathbf 1_{[0,1]}$ this bound
fails, $F$ jumping across $|\xi|=1$; the Riemann-sum defect then carries
a lattice-point (Gauss-circle) term, and $V$ loses the rapidly decaying
tail correction of Theorem~\ref{app:farfield}, acquiring instead an
oscillatory $|y|^{-2}$ tail. Nothing in \S\ref{app:duality} or
\S\ref{app:nondeg} would change, both being insensitive to $\chi$; the
sharp-cutoff case is thus a technical gap in the inner-rescaling step,
not a structural one. It is left open because no statement in this paper
requires it (see Remark~\ref{lstarremark}).
\end{remark}
